% La conscience et l’inconscient — Philosophie 1re Tech — Cours signé OnBuch+
% Programme officiel MINESEC. Compiler avec : tectonic N04-conscience-inconscient.tex
\newcommand{\DOCMATIERE}{Philosophie}
\newcommand{\DOCNIVEAU}{1re Tech}
\newcommand{\DOCMODULE}{Philosophie – classe de Première technique}
\newcommand{\DOCLECON}{N04}
\newcommand{\DOCTITRE}{La conscience et l’inconscient}
% OnBuch+ — préambule « cours signés » ENRICHI (compilé avec tectonic / XeLaTeX)
% Système visuel « Pop » d'OnBuch+ : fond crème (jamais blanc), bordures encre
% épaisses, ombres « dures » décalées, titres Archivo Black en capitales.
% Version MATHÉMATIQUES : graphes (pgfplots/tikz), boîtes pédagogiques
% (définition, propriété, méthode, attention, le savais-tu, exercice résolu,
% exercice, TP, bilan), page de garde et signature OnBuch+ ancrée en bas.
%
% Macros attendues dans main.tex AVANT \input{preamble} :
%   \DOCMATIERE  \DOCNIVEAU  \DOCMODULE  \DOCLECON (numéro)  \DOCTITRE
\documentclass[11pt]{article}

\usepackage{fontspec}
\usepackage[french]{babel}
\usepackage[a4paper,margin=2cm,top=3.4cm,bottom=2.8cm,headheight=1.4cm,footskip=1.3cm]{geometry}
\usepackage{amsmath,amssymb,amsfonts}
\usepackage{mathtools} % \xmapsto, \coloneqq... souvent utilisés par les modèles
\usepackage{cancel} % simplifications barrées (bilans de forces)
% \abs{z} (module, valeur absolue) et \norm{u} (norme), utilisés par les modèles
\providecommand{\abs}{}\let\abs\relax\DeclarePairedDelimiter{\abs}{\lvert}{\rvert}
\providecommand{\norm}{}\let\norm\relax\DeclarePairedDelimiter{\norm}{\lVert}{\rVert}
\usepackage[table]{xcolor} % [table] : fournit aussi \rowcolors (alternance de lignes)
\usepackage{pagecolor}
\usepackage{tikz}
\usetikzlibrary{arrows.meta,positioning,calc,shapes.geometric,shapes.misc,shapes.arrows,decorations.pathmorphing,decorations.pathreplacing,decorations.markings,patterns,shadows,mindmap,trees,fit,backgrounds,matrix,intersections,angles,quotes}
\usepackage{pgfplots}
\usepackage[version=4]{mhchem} % équations nucléaires \ce{^{235}_{92}U}
\usepackage[european,siunitx]{circuitikz} % schémas électriques
\pgfplotsset{compat=1.18}
\usepgfplotslibrary{fillbetween,groupplots} % aires entre courbes (calcul intégral)
\usepackage{enumitem}
\usepackage{fancyhdr}
% [most] charge toutes les bibliothèques tcolorbox (listings, minted, raster,
% documentation...) et gonfle inutilement la mémoire TeX sur un document long
% (~300 boîtes) : on ne charge que ce qui est réellement utilisé.
\usepackage[skins,breakable]{tcolorbox}
\usepackage{graphicx}
\usepackage{colortbl}
\usepackage{booktabs}
\usepackage{tabularx}
\usepackage{multirow}
\usepackage{diagbox} % case d'en-tête diagonale des tableaux à double entrée
% Colonnes extensibles centrée / gauche / droite (C, L, R), souvent utilisées par les modèles
\newcolumntype{C}{>{\centering\arraybackslash}X}
\newcolumntype{L}{>{\raggedright\arraybackslash}X}
\newcolumntype{R}{>{\raggedleft\arraybackslash}X}
\usepackage{array}
\usepackage{multicol}
\usepackage{siunitx}
% parse-numbers=false : accepte aussi \SI{2,5 \times 10^{-3}}{...}
\sisetup{locale=FR,output-decimal-marker={,},group-minimum-digits=4,parse-numbers=false}
\DeclareSIUnit\ohm{\ensuremath{\symup{\Omega}}} % U+2126 absent de la police
\DeclareSIUnit\division{div} % réglages d'oscilloscope (ms/div, V/div)
\DeclareSIUnit\tr{tr} % tours (vitesse de rotation, ex. \SI{1500}{\tr\per\minute})
\DeclareSIUnit\molar{M} % concentration molaire (ex. \SI{3}{\molar}, \SI{1}{\micro\molar})
\DeclareSIUnit\an{an} % année (le modèle confond parfois avec \year, un compteur TeX) — ex. \SI{5}{\kilo\meter\per\an}
\DeclareSIUnit\FCFA{FCFA} % franc CFA (ex. \SI{500}{\FCFA\per\kilo\gram})
\DeclareSIUnit\degreeBrix{\textdegree Brix} % teneur en sucre d'un sirop/jus (réfractométrie)
\DeclareSIUnit\month{mois} % le modèle utilise parfois \month, un compteur TeX, comme unité de durée
\DeclareSIUnit\microliter{\micro\liter} % le modèle écrit parfois \microliter en un seul token
\DeclareSIUnit\million{millions} % dénombrement (ex. \SI{15}{\million\per\mL} spermatozoïdes)
\usepackage{qrcode}
% \degree : commande fréquemment utilisée par les modèles pour un angle ou une
% température en dehors de \SI{}{} (non fournie par les paquets chargés).
\providecommand{\degree}{^{\circ}}
% \qed (fin de démonstration), utilisé par les modèles sans amsthm
\providecommand{\qed}{\hfill$\square$}
% \barre : double barre des valeurs interdites dans un tableau de variations (tkz-tab)
\providecommand{\barre}{\ensuremath{\|}}
% environnement proof (amsthm non chargé) utilisé par les modèles
\ifdefined\proof\else\newenvironment{proof}[1][Démonstration]{\par\noindent\textit{#1.}\ }{\qed\par}\fi
\usepackage{lastpage}
\usepackage{hyperref}
\hypersetup{hidelinks}

% ============================================================
% Palette « Pop » OnBuch+ — mêmes hex que la classe P (plus_ui.dart)
% ============================================================
\definecolor{poporange}{HTML}{F59321}
\definecolor{popdark}{HTML}{E07A0C}
\definecolor{popcream}{HTML}{FFF6E8}
\definecolor{popink}{HTML}{1C1714}
\definecolor{poppurple}{HTML}{7A5AE0}
\definecolor{popgreen}{HTML}{2E9E6A}
\definecolor{poppink}{HTML}{E0457B}
\definecolor{popblue}{HTML}{2F7DE1}
\definecolor{popgold}{HTML}{C98A12}
\definecolor{popmuted}{HTML}{8A7F76}
\definecolor{popline}{HTML}{EADFD2}
\definecolor{popgray}{HTML}{8A7F76}
\colorlet{popgrey}{popgray}
% Alias tolérants (le modèle nomme parfois les couleurs différemment)
\colorlet{popwhite}{white}
\colorlet{popblack}{popink}
\colorlet{popred}{poppink}
\colorlet{popyellow}{popgold}
% Teintes claires (fonds de tableaux/encadrés) — le modèle nomme parfois
% les couleurs avec un suffixe L (ex. popblueL) sans les avoir définies.
\colorlet{poporangeL}{poporange!20}
\colorlet{popdarkL}{popdark!20}
\colorlet{poppurpleL}{poppurple!20}
\colorlet{popgreenL}{popgreen!20}
\colorlet{poppinkL}{poppink!20}
\colorlet{popblueL}{popblue!20}
\colorlet{popgoldL}{popgold!20}
\colorlet{popredL}{popred!20}
\colorlet{popgrayL}{popgray!20}
% Teintes claires pour fonds de tableaux / zones
\colorlet{popblueL}{popblue!10!white}
\colorlet{poporangeL}{poporange!14!white}
\colorlet{popgreenL}{popgreen!12!white}
\colorlet{poppurpleL}{poppurple!12!white}
\colorlet{poppinkL}{poppink!12!white}
\colorlet{popgoldL}{popgold!14!white}

\pagecolor{popcream}

% ============================================================
% Typographie
% ============================================================
\defaultfontfeatures{Ligatures=TeX}
\setmainfont{PlusJakartaSans-Medium.ttf}[Path=../../../fonts/,BoldFont=PlusJakartaSans-Bold.ttf]
% Maths en sans-serif (Fira Math) avec chiffres et lettres droites de Plus
% Jakarta, pour que les formules se fondent dans le texte.
\usepackage[math-style=ISO,bold-style=ISO]{unicode-math}
\setmathfont{FiraMath-Regular.otf}[Path=../../../fonts/]
\setmathfont{PlusJakartaSans-Medium.ttf}[Path=../../../fonts/,range={up/num,up/{Latin,latin},\mathup}]
\usepackage{icomma} % virgule décimale sans espace en mode math
% Caractères mathématiques Unicode absents des polices de texte (Plus Jakarta,
% Archivo Black) : rendus par leur équivalent mathématique, en texte comme en
% formule — sinon ils s'affichent en carré vide (ex. « Arithmétique dans ℤ »).
\usepackage{newunicodechar}
\newunicodechar{ℝ}{\ensuremath{\mathbb{R}}}
\newunicodechar{ℤ}{\ensuremath{\mathbb{Z}}}
\newunicodechar{ℂ}{\ensuremath{\mathbb{C}}}
\newunicodechar{ℕ}{\ensuremath{\mathbb{N}}}
\newunicodechar{ℚ}{\ensuremath{\mathbb{Q}}}
\newunicodechar{𝔻}{\ensuremath{\mathbb{D}}}
\newunicodechar{α}{\ensuremath{\alpha}}
\newunicodechar{β}{\ensuremath{\beta}}
\newunicodechar{γ}{\ensuremath{\gamma}}
\newunicodechar{δ}{\ensuremath{\delta}}
\newunicodechar{ε}{\ensuremath{\varepsilon}}
\newunicodechar{θ}{\ensuremath{\theta}}
\newunicodechar{λ}{\ensuremath{\lambda}}
\newunicodechar{μ}{\ensuremath{\mu}}
\newunicodechar{σ}{\ensuremath{\sigma}}
\newunicodechar{τ}{\ensuremath{\tau}}
\newunicodechar{φ}{\ensuremath{\varphi}}
\newunicodechar{ψ}{\ensuremath{\psi}}
\newunicodechar{ω}{\ensuremath{\omega}}
\newunicodechar{Σ}{\ensuremath{\Sigma}}
% majuscules produites par \MakeUppercase dans les titres (α-aminés -> Α)
\newunicodechar{Α}{\ensuremath{\alpha}}
\newunicodechar{Β}{\ensuremath{\beta}}
\newunicodechar{Γ}{\ensuremath{\Gamma}}
\newunicodechar{Δ}{\ensuremath{\Delta}}
\newunicodechar{Ε}{\ensuremath{\varepsilon}}
\newunicodechar{Θ}{\ensuremath{\Theta}}
\newunicodechar{Λ}{\ensuremath{\Lambda}}
\newunicodechar{Μ}{\ensuremath{\mu}}
\newunicodechar{Τ}{\ensuremath{\tau}}
\newunicodechar{Φ}{\ensuremath{\Phi}}
\newunicodechar{Ψ}{\ensuremath{\Psi}}
\newunicodechar{Ω}{\ensuremath{\Omega}}
\newunicodechar{↦}{\ensuremath{\mapsto}}
\newunicodechar{∈}{\ensuremath{\in}}
\newunicodechar{∉}{\ensuremath{\notin}}
\newunicodechar{∧}{\ensuremath{\wedge}}
\newunicodechar{⇔}{\ensuremath{\Leftrightarrow}}
\newunicodechar{⟺}{\ensuremath{\Longleftrightarrow}}
\newunicodechar{⇒}{\ensuremath{\Rightarrow}}
\newunicodechar{⟹}{\ensuremath{\Longrightarrow}}
\newunicodechar{∘}{\ensuremath{\circ}}
\newunicodechar{∪}{\ensuremath{\cup}}
\newunicodechar{∩}{\ensuremath{\cap}}
\newunicodechar{≡}{\ensuremath{\equiv}}
\newunicodechar{⊂}{\ensuremath{\subset}}
\newunicodechar{⊥}{\ensuremath{\perp}}
\newunicodechar{∃}{\ensuremath{\exists}}
\newunicodechar{∀}{\ensuremath{\forall}}
\newunicodechar{∛}{\ensuremath{\sqrt[3]{\,}}}
\newunicodechar{∜}{\ensuremath{\sqrt[4]{\,}}}
\newunicodechar{⃗}{\ensuremath{{}^{\to}}}
\newunicodechar{ⁿ}{\ifmmode^{n}\else\textsuperscript{\ensuremath{n}}\fi}
\newunicodechar{ˣ}{\ifmmode^{x}\else\textsuperscript{\ensuremath{x}}\fi}
\newunicodechar{ᵇ}{\ifmmode^{b}\else\textsuperscript{\ensuremath{b}}\fi}
\newunicodechar{ʳ}{\ifmmode^{r}\else\textsuperscript{\ensuremath{r}}\fi}
\newunicodechar{ᵘ}{\ifmmode^{u}\else\textsuperscript{\ensuremath{u}}\fi}
\newunicodechar{ʸ}{\ifmmode^{y}\else\textsuperscript{\ensuremath{y}}\fi}
\newunicodechar{ᵃ}{\ifmmode^{a}\else\textsuperscript{\ensuremath{a}}\fi}
\newunicodechar{ᵉ}{\ifmmode^{e}\else\textsuperscript{\ensuremath{e}}\fi}
\newunicodechar{ᵏ}{\ifmmode^{k}\else\textsuperscript{\ensuremath{k}}\fi}
\newunicodechar{ᵖ}{\ifmmode^{p}\else\textsuperscript{\ensuremath{p}}\fi}
\newunicodechar{ᴺ}{\ifmmode^{N}\else\textsuperscript{\ensuremath{N}}\fi}
\newunicodechar{ᶜ}{\ifmmode^{c}\else\textsuperscript{\ensuremath{c}}\fi}
\newunicodechar{ʰ}{\ifmmode^{h}\else\textsuperscript{\ensuremath{h}}\fi}
\newunicodechar{ᵠ}{\ifmmode^{q}\else\textsuperscript{\ensuremath{q}}\fi}
\newunicodechar{ₙ}{\ifmmode_{n}\else\textsubscript{\ensuremath{n}}\fi}
\newunicodechar{ₐ}{\ifmmode_{a}\else\textsubscript{\ensuremath{a}}\fi}
\newunicodechar{ᵢ}{\ifmmode_{i}\else\textsubscript{\ensuremath{i}}\fi}
\newunicodechar{ᵣ}{\ifmmode_{r}\else\textsubscript{\ensuremath{r}}\fi}
\newunicodechar{ₖ}{\ifmmode_{k}\else\textsubscript{\ensuremath{k}}\fi}
\newunicodechar{ᵦ}{\ifmmode_{\beta}\else\textsubscript{\ensuremath{\beta}}\fi}
\newunicodechar{ₓ}{\ifmmode_{x}\else\textsubscript{\ensuremath{x}}\fi}
\newfontfamily\popdisplay{ArchivoBlack-Regular.ttf}[Path=../../../fonts/]
\newfontfamily\poplogo{SpaceGrotesk-Bold.ttf}[Path=../../../fonts/]
\newfontfamily\popsemi{PlusJakartaSans-SemiBold.ttf}[Path=../../../fonts/]
\newfontfamily\popxbold{PlusJakartaSans-ExtraBold.ttf}[Path=../../../fonts/]
\newfontfamily\popxboldls{PlusJakartaSans-ExtraBold.ttf}[Path=../../../fonts/,LetterSpace=3.0]

\setlist[enumerate]{leftmargin=1.7em,itemsep=5pt,topsep=4pt}
\setlist[itemize]{leftmargin=1.7em,itemsep=4pt,topsep=4pt,label={\color{poporange}\textbullet}}
\setlength{\parindent}{0pt}
\setlength{\parskip}{5pt}
\renewcommand{\arraystretch}{1.25}

% pgfplots : style Pop par défaut
\pgfplotsset{
  every axis/.append style={axis line style={popink,line width=1pt},tick style={popink},
    grid style={popline,line width=0.5pt},label style={font=\small},tick label style={font=\footnotesize}},
  popcurve/.style={poporange,line width=1.8pt,no markers,smooth},
}
% Même style disponible dans un \draw TikZ simple (hors pgfplots)
\tikzset{popcurve/.style={poporange,line width=1.8pt}}
\tikzset{popgrid/.style={popline,very thin}}
\colorlet{popcurve}{poporange} % le nom du style est parfois utilisé comme couleur

% ============================================================
% Pill / squiggle
% ============================================================
\newcommand{\poppill}[3]{%
  \tikz[baseline=-0.55ex]\node[draw=popink,line width=1.2pt,rounded corners=2.6pt,%
    fill=#2,inner xsep=6.5pt,inner ysep=3pt]{%
    \popxboldls\fontsize{7}{7}\selectfont\color{#3}\MakeUppercase{#1}};%
}
\newcommand{\popsquiggle}[1]{%
  \tikz[baseline=-0.4ex]{
    \draw[#1,line width=2.6pt,line cap=round]
      plot [smooth] coordinates {(0,0.09) (0.34,0.01) (0.68,0.21) (1.02,0.01) (1.36,0.10)};
  }%
}
% Mot-clé surligné (marqueur orange sous le texte)
\newcommand{\cle}[1]{{\popxbold\color{popdark}#1}}

% ============================================================
% En-tête / pied
% ============================================================
\pagestyle{fancy}
\fancyhf{}
\renewcommand{\headrulewidth}{0pt}
\renewcommand{\footrulewidth}{0pt}
\lhead{\raisebox{-0.15ex}{\poplogo\fontsize{13}{13}\selectfont\color{popink}OnBuch\color{poporange}+}}
\rhead{\poppill{\DOCMATIERE\ \textbullet\ \DOCNIVEAU\ \textbullet\ Leçon \DOCLECON}{white}{popink}}
\lfoot{\popxbold\fontsize{7.6}{7.6}\selectfont\color{popmuted}\MakeUppercase{Cours signé OnBuch+}\ \textbullet\ {\popsemi\DOCTITRE}}
\rfoot{\tikz[baseline=-0.6ex]\node[rounded corners=8.5pt,fill=popink,minimum height=17pt,minimum width=17pt,inner xsep=5pt,inner sep=0]{\popxbold\fontsize{8}{8}\selectfont\color{popcream}\thepage\,/\,\pageref*{LastPage}};}

% ============================================================
% Titres
% ============================================================
\newcommand{\chaptitre}[2]{%
  \begin{tcolorbox}[enhanced,colback=poporange,colframe=popink,boxrule=2.6pt,arc=11pt,
      drop shadow=popink,left=15pt,right=15pt,top=12pt,bottom=11pt,before skip=3pt,after skip=18pt]
    \poppill{#2}{popink}{popcream}\par\vspace{9pt}
    {\raggedright\hyphenpenalty=10000\exhyphenpenalty=10000\popdisplay\fontsize{22}{24}\selectfont\color{popcream}\MakeUppercase{#1}\par}
  \end{tcolorbox}
}
% Section numérotée : pastille orange avec le numéro + titre Archivo
\newcounter{popsec}
\newcommand{\coursec}[1]{%
  \stepcounter{popsec}\setcounter{popsub}{0}%
  \par\vspace{16pt}\noindent
  \begin{minipage}{\linewidth}
  \tikz[baseline=-0.7ex]\node[draw=popink,line width=1.6pt,fill=poporange,rounded corners=4pt,minimum size=22pt,inner sep=0]{\popdisplay\fontsize{12}{12}\selectfont\color{popcream}\thepopsec};%
  \hspace{8pt}{\popdisplay\fontsize{15}{17}\selectfont\color{popink}\MakeUppercase{#1}}\par
  \vspace{4pt}\noindent{\color{poporange}\rule{36pt}{3.6pt}}
  \end{minipage}\par\nopagebreak\vspace{8pt}
}
\newcounter{popsub}
\newcommand{\courssub}[1]{%
  \stepcounter{popsub}%
  \par\vspace{9pt}\noindent{\popxbold\fontsize{11.5}{13}\selectfont\color{popink}{\color{poporange}\thepopsec.\thepopsub}\hspace{6pt}#1}\par\nopagebreak\vspace{4pt}
}

% ============================================================
% Boîtes pédagogiques — même recette : fond blanc, bordure épaisse colorée,
% ombre dure encre, badge plein. Titre optionnel [#1].
% ============================================================
\tcbset{popbase/.style={breakable,enhanced,colback=white,boxrule=2.3pt,arc=9pt,
  shadow={3pt}{-3pt}{0pt}{popink},left=14pt,right=14pt,top=11pt,bottom=11pt,before skip=11pt,after skip=15pt,
  fontupper=\popsemi\fontsize{10.6}{14.5}\selectfont\color{popink},
  fontlower=\popsemi\fontsize{10.6}{14.5}\selectfont\color{popink}}}
% Coche dessinée (le glyphe ✓ n'existe pas dans les polices du document)
\newcommand{\popcheck}[1]{\tikz[baseline=-0.5ex]\draw[#1,line width=1.6pt,line cap=round,line join=round](-0.13,0.0)--(-0.04,-0.09)--(0.13,0.1);}
\providecommand{\coche}{\popcheck{popgreen}}
\providecommand{\resultat}[1]{\par\smallskip\noindent\fcolorbox{popgreen}{popgreenL}{\parbox{\dimexpr\linewidth-2\fboxsep-2\fboxrule\relax}{\textbf{Résultat.} #1}}\par\smallskip}
\newcommand{\popboxhead}[3]{\poppill{#1}{#2}{white}\ifx\relax#3\relax\else\hspace{6pt}{\popxbold\fontsize{10.6}{12}\selectfont\color{popink}#3}\fi\par\vspace{7pt}}

\newtcolorbox{aretenir}[1][]{popbase,colframe=popgreen,before upper={\popboxhead{À retenir}{popgreen}{#1}}}
\newtcolorbox{exemplebox}[1][]{popbase,colframe=poporange,before upper={\popboxhead{Exemple}{poporange}{#1}}}
\newtcolorbox{definition}[1][]{popbase,colframe=popblue,before upper={\popboxhead{Définition}{popblue}{#1}}}
\newtcolorbox{propriete}[1][]{popbase,colframe=poppurple,before upper={\popboxhead{Propriété}{poppurple}{#1}}}
\newtcolorbox{methode}[1][]{popbase,colframe=popgold,before upper={\popboxhead{Méthode}{popgold}{#1}}}
\newtcolorbox{attention}[1][]{popbase,colframe=poppink,before upper={\popboxhead{Attention}{poppink}{#1}}}
\newtcolorbox{experience}[1][]{popbase,colframe=popblue,colback=popblueL,before upper={\popboxhead{Expérience}{popblue}{#1}}}
% « Le savais-tu ? » : carte violette pleine (contexte, culture, Cameroun)
\newtcolorbox{savaistu}[1][]{popbase,colback=poppurple,colframe=popink,
  fontupper=\popsemi\fontsize{10.4}{14.2}\selectfont\color{white},
  before upper={\poppill{Le savais-tu\,?}{white}{poppurple}\ifx\relax#1\relax\else\hspace{6pt}{\popxbold\color{white}#1}\fi\par\vspace{7pt}}}
% Exercice résolu : énoncé en haut, \tcblower puis la solution sur fond crème
\newcounter{popexo}\newcounter{popexr}
\newtcolorbox{exoresolu}[1][]{popbase,colframe=poporange,colbacklower=popcream,
  segmentation style={popink,line width=1.2pt,dashed},
  before upper={\stepcounter{popexr}\popboxhead{Exercice résolu \thepopexr}{poporange}{#1}},
  before lower={\poppill{Solution}{popink}{popcream}\par\vspace{6pt}}}
% Exercice (non résolu en place) — la correction est dans la partie Corrigés
\newtcolorbox{exercice}[1][]{popbase,colframe=popink,
  before upper={\stepcounter{popexo}\popboxhead{Exercice \thepopexo}{popink}{#1}}}
% Corrigé
\newtcolorbox{corrige}[1][]{popbase,colframe=popgreen,colback=popgreenL,
  before upper={\popboxhead{Corrigé}{popgreen}{#1}}}
% Objectifs / prérequis / bilan
\newtcolorbox{objectifs}{popbase,colframe=popink,colback=poporangeL,
  before upper={\popboxhead{Objectifs de la leçon}{popink}{}}}
\newtcolorbox{prerequis}{popbase,colframe=popink,colback=popblueL,
  before upper={\popboxhead{Prérequis}{popblue}{}}}
\newtcolorbox{bilan}{popbase,colframe=popink,colback=white,boxrule=2.8pt,
  before upper={\popboxhead{Fiche bilan}{poporange}{L'essentiel à connaître pour l'examen}}}
% Figure encadrée avec légende
\newtcolorbox{popfigure}[1][]{enhanced,breakable=false,colback=white,colframe=popline,boxrule=1.4pt,arc=8pt,
  left=10pt,right=10pt,top=10pt,bottom=8pt,before skip=10pt,after skip=12pt,halign=center,
  lower separated=false,halign lower=center,
  fontlower=\popsemi\fontsize{9}{11.5}\selectfont\color{popmuted},
  #1}
\newcommand{\legende}[1]{\par\vspace{6pt}{\popsemi\fontsize{9}{11.5}\selectfont\color{popmuted}#1}}

% ============================================================
% Signature OnBuch+
% ============================================================
\newcommand{\poplogoinline}[1]{{\poplogo\fontsize{#1}{#1}\selectfont\color{popink}OnBuch\color{poporange}+}}
\newcommand{\signatureonbuch}{%
  \begin{tcolorbox}[enhanced,colback=popink,colframe=popink,boxrule=2.2pt,arc=10pt,
      drop shadow=poporange,left=14pt,right=14pt,top=10pt,bottom=10pt,before skip=0pt,after skip=0pt]
    \tikz[baseline=-0.7ex]\node[circle,fill=popgreen,draw=popcream,line width=1.3pt,minimum size=22pt,inner sep=0]{\popcheck{white}};%
    \hspace{9pt}\begin{minipage}[c]{0.62\linewidth}
      {\popxbold\fontsize{12}{13}\selectfont\color{popcream}Signé\ \textbullet\ {\poplogo OnBuch\color{poporange}+}}\par\vspace{2pt}
      {\raggedright\popsemi\fontsize{8}{10.5}\selectfont\color{popline}\MakeUppercase{Équipe pédagogique OnBuch+}\\\MakeUppercase{Conforme au programme officiel MINESEC}\par}
    \end{minipage}\hfill
    {\poplogo\fontsize{22}{22}\selectfont\color{popcream}OnBuch\color{poporange}+}
  \end{tcolorbox}%
}

% ============================================================
% Page de garde
% #1 titre · #2 matière · #3 niveau · #4 module · #5 n° leçon · #6 durée ·
% #7 description / objectifs (texte ou itemize)
% ============================================================
\newcommand{\pagegarde}[7]{%
  \thispagestyle{empty}
  \newgeometry{a4paper,margin=1.7cm,top=1.6cm,bottom=1.4cm}%
  \begin{tikzpicture}[remember picture,overlay]
    \fill[poporange!18] ([xshift=-3.2cm,yshift=-3.6cm]current page.north east) circle (4.6cm);
    \fill[poppurple!14] ([xshift=2.4cm,yshift=7.6cm]current page.south west) circle (3.8cm);
  \end{tikzpicture}%
  {\poplogo\fontsize{30}{30}\selectfont\color{popink}OnBuch\color{poporange}+}\hfill
  \raisebox{0.6ex}{\poppill{Cours signé \textbullet\ Premium}{popink}{popcream}}\par
  \vspace{1pt}\popsquiggle{poppurple}\par
  \vspace{8pt}
  \begin{tcolorbox}[enhanced,colback=poporange,colframe=popink,boxrule=2.8pt,arc=13pt,
      drop shadow=popink,left=18pt,right=18pt,top=12pt,bottom=13pt,after skip=10pt]
    \poppill{#2 \textbullet\ #3}{popink}{popcream}\hspace{5pt}\poppill{#4}{white}{popink}\par\vspace{8pt}
    {\popdisplay\fontsize{14}{14}\selectfont\color{popink}LEÇON #5}\par\vspace{4pt}
    {\raggedright\hyphenpenalty=10000\exhyphenpenalty=10000\popdisplay\fontsize{25}{27}\selectfont\color{popcream}\MakeUppercase{#1}\par}
    \vspace{8pt}
    {\popxbold\fontsize{9.5}{9.5}\selectfont\color{popink}\MakeUppercase{Durée conseillée : #6}}
  \end{tcolorbox}
  \begin{tcolorbox}[enhanced,colback=white,colframe=popink,boxrule=2.2pt,arc=10pt,
      drop shadow=popink,left=16pt,right=16pt,top=10pt,bottom=10pt,after skip=8pt]
    \poppill{Dans cette leçon}{poppurple}{white}\par\vspace{5pt}
    \popsemi\fontsize{9.3}{12.6}\selectfont\color{popink}#7
  \end{tcolorbox}
  \noindent\begin{minipage}[t]{0.487\linewidth}
    \begin{tcolorbox}[enhanced,colback=popgreen,colframe=popink,boxrule=2.2pt,arc=10pt,
        drop shadow=popink,left=11pt,right=11pt,top=10pt,bottom=10pt,height=3.1cm,valign=center]
      \tcbox[colback=white,colframe=popink,boxrule=1.3pt,arc=5pt,boxsep=0pt,nobeforeafter,
        left=3pt,right=3pt,top=3pt,bottom=3pt,valign=center]{\qrcode[height=1.9cm]{https://play.google.com/store/apps/details?id=com.luvvix.onbuch&hl=fr}}%
      \hspace{8pt}\begin{minipage}[c]{0.5\linewidth}
        {\popdisplay\fontsize{11}{12.5}\selectfont\color{white}\MakeUppercase{Télécharge OnBuch}}\par\vspace{4pt}
        {\raggedright\popsemi\fontsize{8.4}{11}\selectfont\color{white}Gratuit \textbullet\ Google Play\\Tuteur IA, annales, concours, résultats\par}
      \end{minipage}
    \end{tcolorbox}
  \end{minipage}\hfill
  \begin{minipage}[t]{0.487\linewidth}
    \begin{tcolorbox}[enhanced,colback=poppurple,colframe=popink,boxrule=2.2pt,arc=10pt,
        drop shadow=popink,left=11pt,right=11pt,top=10pt,bottom=10pt,height=3.1cm,valign=center]
      \tikz[baseline=-0.7ex]\node[circle,fill=white,draw=popink,line width=1.3pt,minimum size=1.6cm,inner sep=0]{\popdisplay\fontsize{24}{24}\selectfont\color{poppurple}+};%
      \hspace{8pt}\begin{minipage}[c]{0.6\linewidth}
        {\popdisplay\fontsize{11}{12.5}\selectfont\color{white}\MakeUppercase{Passe à OnBuch+}}\par\vspace{4pt}
        {\raggedright\popsemi\fontsize{8.4}{11}\selectfont\color{white}Tous les cours signés, Tuteur IA illimité, fiches PDF — onglet OnBuch+ de l'app\par}
      \end{minipage}
    \end{tcolorbox}
  \end{minipage}
  \vspace{8pt}
  \popcontenu
  \vfill
  \signatureonbuch
  \restoregeometry
  \newpage
}

% Rangée « Ce cours contient » (page de garde)
\newcommand{\poptile}[3]{% #1 couleur #2 gros texte #3 légende
  \begin{tcolorbox}[enhanced,colback=#1,colframe=popink,boxrule=2pt,arc=9pt,shadow={2.5pt}{-2.5pt}{0pt}{popink},
      left=6pt,right=6pt,top=9pt,bottom=9pt,halign=center,valign=center,height=1.75cm,width=\linewidth,nobeforeafter]
    {\popdisplay\fontsize{12.5}{13}\selectfont\color{white}\MakeUppercase{#2}}\par\vspace{3pt}
    {\centering\popsemi\fontsize{7.8}{9.5}\selectfont\color{white}#3\par}
  \end{tcolorbox}}
\newcommand{\popcontenu}{%
  \par\noindent{\popxboldls\fontsize{7.5}{7.5}\selectfont\color{popmuted}CE COURS SIGNÉ CONTIENT}\par\vspace{7pt}
  \noindent\begin{minipage}[t]{0.235\linewidth}\poptile{popblue}{Cours}{complet, illustré, pas à pas}\end{minipage}\hfill
  \begin{minipage}[t]{0.235\linewidth}\poptile{poporange}{Méthodes}{et exercices résolus}\end{minipage}\hfill
  \begin{minipage}[t]{0.235\linewidth}\poptile{poppink}{Exercices}{type examen + corrigés}\end{minipage}\hfill
  \begin{minipage}[t]{0.235\linewidth}\poptile{popgreen}{Fiche bilan}{l'essentiel à retenir}\end{minipage}\par}

% Sommaire
\newtcolorbox{sommaire}{enhanced,colback=white,colframe=popink,boxrule=2.2pt,arc=10pt,
  drop shadow=popink,left=18pt,right=18pt,top=13pt,bottom=13pt,before skip=6pt,after skip=20pt,
  before upper={\poppill{Sommaire}{poppurple}{white}\par\vspace{9pt}\popsemi\fontsize{10.6}{14.5}\selectfont\color{popink}}}

% Signature de fin de cours — poussée tout en bas de la dernière page
\newcommand{\findecours}{%
  \par\vfill\nopagebreak
  \begin{center}\popsquiggle{poporange}\end{center}\vspace{4pt}
  \signatureonbuch
}
\AtBeginDocument{\def\llbracket{\mathopen{[\mkern-3.5mu[}}\def\rrbracket{\mathclose{]\mkern-3.5mu]}}} % intervalles d'entiers (glyphe ⟦ absent de la police)
% Glyphes absents de Fira Math : redéfinis à partir de glyphes disponibles
\AtBeginDocument{%
  \def\setminus{\mathbin{\backslash}}\let\smallsetminus\setminus
  \def\vdots{\mathord{\vbox{\baselineskip3.2pt\lineskiplimit0pt\kern4pt\hbox{.}\hbox{.}\hbox{.}}}}%
  \def\ddots{\mathinner{\mkern1mu\raise6.5pt\hbox{.}\mkern2mu\raise3.6pt\hbox{.}\mkern2mu\raise0.7pt\hbox{.}\mkern1mu}}%
  \def\triangle{\mathord{\scalebox{0.9}{$\symup{\Delta}$}}}%
  \def\nabla{\mathord{\raisebox{\depth}{\scalebox{1}[-1]{$\symup{\Delta}$}}}}%
  \def\checkmark{\popcheck{popdark}}%
  \def\star{\mathbin{\ast}}\def\bigstar{\mathord{\ast}}%
  \def\diamond{\mathbin{\scalebox{0.6}{$\Diamond$}}}%
  \def\bigcup{\mathop{\mathchoice{\raisebox{-0.3em}{\scalebox{1.7}{$\cup$}}}{\scalebox{1.25}{$\cup$}}{\cup}{\cup}}\displaylimits}%
  \def\bigcap{\mathop{\mathchoice{\raisebox{-0.3em}{\scalebox{1.7}{$\cap$}}}{\scalebox{1.25}{$\cap$}}{\cap}{\cap}}\displaylimits}%
  \def\lesseqqgtr{\mathrel{\lessgtr}}\def\bigoplus{\mathop{\oplus}\displaylimits}%
}
\providecommand{\ssi}{si et seulement si} % abréviation fréquente
\AtBeginDocument{\providecommand{\wideparen}{\overparen}} % arc de cercle

\begin{document}

\pagegarde{La conscience et l’inconscient}{Philosophie}{1re Tech}{Philosophie – classe de Première technique}{N04}{2 séances (fiche de progression harmonisée)}{Cette leçon interroge la nature de la conscience humaine, ses types, sa valeur et ses limites, puis présente la découverte de l'inconscient et la théorie psychanalytique de Freud, objet de vifs débats. Elle conduit l'élève à se demander si la conscience suffit à rendre compte de la vie psychique humaine.
\vspace{4pt}
\begin{multicols}{2}\small
\begin{itemize}
\item À la fin de cette leçon, je sais définir la conscience et distinguer ses différents types (conscience spontanée, réfléchie, morale)
\item Identifier les caractéristiques de la conscience (intentionnalité, unité, temporalité, spontanéité)
\item Évaluer la valeur de la conscience comme connaissance de soi et comme fondement de la liberté et de la responsabilité
\item Reconnaître les limites de la conscience (erreurs des sens, illusions, déterminismes sociaux et biologiques)
\item Expliquer comment l'hypothèse de l'inconscient a été historiquement découverte (hypnose, hystérie, Freud)
\item Distinguer les types d'inconscient (inconscient psychologique freudien, inconscient collectif de Jung)
\end{itemize}\end{multicols}}

\begin{sommaire}
\begin{multicols}{2}
\begin{enumerate}
\item La conscience : notion et types
\item Caractéristiques de la conscience
\item Valeur de la conscience
\item Limites de la conscience
\item La découverte de l'inconscient et ses types
\item La psychanalyse freudienne et sa critique
\item Synthèse : conscience et inconscient, complémentaires ou adversaires ?
\item Activité d'intégration
\item Méthodes et exercices résolus
\item Exercices
\item Corrigés des exercices
\item Fiche bilan
\end{enumerate}
\end{multicols}
\end{sommaire}

\begin{objectifs}
\begin{itemize}
\item À la fin de cette leçon, je sais définir la conscience et distinguer ses différents types (conscience spontanée, réfléchie, morale)
\item Identifier les caractéristiques de la conscience (intentionnalité, unité, temporalité, spontanéité)
\item Évaluer la valeur de la conscience comme connaissance de soi et comme fondement de la liberté et de la responsabilité
\item Reconnaître les limites de la conscience (erreurs des sens, illusions, déterminismes sociaux et biologiques)
\item Expliquer comment l'hypothèse de l'inconscient a été historiquement découverte (hypnose, hystérie, Freud)
\item Distinguer les types d'inconscient (inconscient psychologique freudien, inconscient collectif de Jung)
\item Décrire la théorie freudienne de la psychanalyse (Ça, Moi, Surmoi ; refoulement ; interprétation des rêves et des actes manqués)
\item Discuter les critiques adressées à la psychanalyse (Sartre, Popper, sciences cognitives)
\item Appliquer ces notions à des situations concrètes de la vie courante et rédiger une argumentation justifiée
\end{itemize}
\end{objectifs}

\begin{prerequis}
\begin{itemize}
\item Notion de sujet et d'identité personnelle (classe de Seconde)
\item Distinction entre vérité et erreur, connaissance et illusion
\item Notions de liberté et de responsabilité morale
\item Vocabulaire de base : psyché, comportement, désir, pulsion
\end{itemize}
\end{prerequis}

\newpage

\coursec{La conscience : notion et types}

\courssub{Situation d'accroche et problématique}

À Yaoundé, dans un lycée technique de la capitale, un élève de Première jure à son professeur qu'il n'a \og pas fait exprès \fg{} d'oublier ses devoirs trois fois de suite. Sa mère, convoquée au collège, affirme qu'il ment : \og Tu savais très bien que tu devais rendre ce travail ! \fg{}. Son meilleur ami, lui, pense qu'il dit vrai : \og Il a vraiment oublié, il était distrait ces jours-ci. \fg{} Mais le plus troublant, c'est que l'élève lui-même hésite : était-ce un simple oubli, un accident de la mémoire, ou bien un refus caché de travailler, une paresse qu'il n'ose pas s'avouer ? Sigmund Freud, le fondateur de la psychanalyse, dirait que cet oubli n'est pas un hasard : il révèle un \cle{désir inconscient}, celui d'échapper à une tâche désagréable. Qui a raison ? L'élève, sa mère, son ami, ou Freud ?

Cette petite scène de la vie scolaire camerounaise pose en réalité l'un des plus grands problèmes de la philosophie : l'homme est-il \cle{maître de ses actes} grâce à sa conscience, ou est-il \cle{déterminé} par des forces psychiques qui lui échappent ? Pour répondre à cette question, il faut d'abord savoir ce qu'est la conscience, combien de types de conscience existent, et ce que chacun d'eux nous permet de faire. C'est l'objet de cette première section.

\textbf{Problématique :} qu'est-ce que la conscience ? Se réduit-elle à une simple connaissance de soi, ou comporte-t-elle plusieurs formes distinctes ?

\courssub{Étymologie et sens du mot « conscience »}

Commençons par le mot lui-même, car en philosophie, l'étymologie éclaire souvent le sens profond d'une notion. Le mot \emph{conscience} vient du latin \emph{conscientia}, formé à partir du verbe \emph{cum-scire}, qui signifie littéralement \og \cle{connaître avec} \fg{}. Avoir conscience, c'est donc d'abord \emph{connaître avec soi-même} : être présent à ce que l'on fait, à ce que l'on pense, à ce que l'on ressent, comme si l'on était à la fois l'acteur et le témoin de sa propre vie.

Cette étymologie révèle un \cle{double sens} fondamental du mot, que la langue française a conservé :

\begin{enumerate}
\item Le sens \cle{psychologique} (ou cognitif) : la conscience comme \emph{connaissance}. Dire \og j'ai conscience du danger \fg{}, c'est dire que je le connais, que je m'en rends compte. C'est la conscience comme lumière intérieure qui éclaire mes états et le monde.
\item Le sens \cle{moral} : la conscience comme \emph{témoignage intérieur} sur le bien et le mal. Dire \og j'ai mauvaise conscience \fg{} ou \og ma conscience me reproche cela \fg{}, c'est évoquer un juge intérieur qui évalue mes actes. En latin, \emph{conscientia} désignait déjà ce témoignage que l'on se rend à soi-même de ses propres actions, bonnes ou mauvaises.
\end{enumerate}

\begin{attention}[Ne pas confondre les deux sens du mot « conscience »]
Erreur fréquente à l'examen : confondre la \emph{conscience} au sens psychologique (la connaissance que le sujet a de lui-même) et la \emph{conscience morale} (le sens du devoir, le jugement du bien et du mal). Les deux sens coexistent en français dans un seul mot, ce qui n'est pas le cas dans toutes les langues (l'allemand distingue \emph{Bewusstsein} et \emph{Gewissen}). Dans une dissertation, précisez toujours de quel sens vous parlez : \og la conscience comme connaissance de soi \fg{} ou \og la conscience comme sens moral \fg{}. Une copie qui mélange les deux sans les distinguer perd en rigueur.
\end{attention}

\courssub{Définition générale de la conscience}

Partons d'une intuition simple. Fermez les yeux un instant : vous entendez des bruits, vous sentez votre corps sur la chaise, vous vous souvenez peut-être de ce que vous avez mangé ce matin. Or, vous ne faites pas que percevoir ces choses : vous \emph{savez} que vous les percevez. Vous êtes présent à vous-même. Cette présence à soi, c'est la conscience. Un animal perçoit aussi le monde, mais l'homme, lui, sait qu'il perçoit ; il peut se dire \og c'est moi qui vois, c'est moi qui souffre, c'est moi qui décide \fg{}.

\begin{definition}[La conscience]
La \cle{conscience} est la \emph{connaissance immédiate} que le sujet a de lui-même, de ses états (pensées, émotions, sensations), de ses actes et du monde extérieur. Elle est à la fois \emph{connaissance} (je sais ce que je fais et ce que je vis) et \emph{témoignage intérieur} (je me rends compte que c'est bien moi qui agis et qui pense).
\end{definition}

Trois éléments de cette définition méritent d'être soulignés :

\begin{itemize}
\item La conscience est une connaissance \emph{immédiate} : elle ne passe pas par un raisonnement. Je n'ai pas besoin de démontrer que j'ai mal à la tête : je le sais directement, par le seul fait de le vivre.
\item Elle porte sur \emph{trois objets} : moi-même (mes états intérieurs), mes actes (ce que je fais) et le monde extérieur (ce qui m'entoure).
\item Elle suppose un \emph{sujet} : il n'y a de conscience que pour un \og je \fg{}, un être capable de se rapporter à lui-même.
\end{itemize}

\courssub{Les trois types de conscience}

La conscience n'est pas un bloc uniforme. Les philosophes distinguent classiquement \cle{trois types de conscience}, que l'on peut imaginer comme trois cercles concentriques, du plus large (le plus fondamental) au plus restreint (le plus élaboré).

\textbf{1. La conscience spontanée (ou immédiate).} C'est le niveau le plus simple : la \emph{présence simple à soi et aux choses}, sans réflexion. Quand je marche dans la rue de Douala en évitant les flaques, quand je réponds \og présent ! \fg{} à l'appel, quand je goûte un plat d'eru bien assaisonné, je suis conscient de ce que je fais et de ce que je sens, mais je n'y pense pas spécialement. Ma conscience accompagne mes actes comme une lumière qui suit le voyageur sans qu'il la regarde. C'est la conscience de l'action vécue, avant toute analyse.

\textbf{2. La conscience réfléchie.} C'est le \emph{retour du sujet sur lui-même} : au lieu de simplement vivre mes pensées, je les examine. Je me demande : \og pourquoi ai-je dit cela ? Ai-je bien fait ? Que dois-je penser de cette idée ? \fg{}. La conscience réfléchie est comme un miroir que l'esprit se tend à lui-même. C'est elle qui fait de l'homme un être capable de \emph{douter}, de \emph{s'interroger} et de se connaître. Elle est au cœur de la philosophie de René Descartes (1596-1650).

\begin{exemplebox}[Le « Cogito » de Descartes : fondement de la conscience réfléchie]
Dans le \emph{Discours de la méthode} (1637), Descartes décide de douter de tout : les sens peuvent nous tromper, les rêves imitent la réalité, les raisonnements peuvent être faux. Mais en doutant de tout, il découvre une vérité indubitable : \og \cle{je pense, donc je suis} \fg{} (\emph{Cogito, ergo sum}). Même si je me trompe sur tout, je ne peux pas me tromper sur le fait que je pense, car pour me tromper, il faut d'abord que je pense ! Le Cogito est l'acte par lequel la conscience se saisit elle-même : c'est le modèle même de la conscience réfléchie, ce retour du sujet sur sa propre pensée qui fonde toute certitude.
\end{exemplebox}

\textbf{3. La conscience morale.} C'est la \emph{capacité de juger le bien et le mal de ses propres actions}. Elle ne se contente pas de connaître : elle \emph{évalue}. On l'appelle parfois le \og \cle{tribunal intérieur} \fg{} : après une faute, elle nous juge et nous condamne (remords) ; avant l'action, elle nous avertit ; après une bonne action, elle nous acquitte et nous apaise.

\begin{definition}[La conscience morale]
La \cle{conscience morale} est la faculté par laquelle le sujet juge ses propres actions (passées, présentes ou à venir) selon les normes du bien et du mal, éprouvant l'approbation (bonne conscience) ou la réprobation (mauvaise conscience, remords) envers lui-même.
\end{definition}

\begin{popfigure}
\begin{center}
\begin{tikzpicture}[scale=1.1, every node/.style={font=\small}]
% Cercles concentriques (cercles pleins avec opacité pour voir les couches)
\fill[popgoldL, opacity=0.6] (0,0) circle (3.0);
\fill[poporangeL, opacity=0.7] (0,0) circle (2.0);
\fill[popblueL, opacity=0.8] (0,0) circle (1.0);

% Contours
\draw[popdark, thick] (0,0) circle (3.0);
\draw[popdark, thick] (0,0) circle (2.0);
\draw[popdark, thick] (0,0) circle (1.0);

% Étiquettes placées dans les anneaux correspondants
\node[align=center, text=popdark, font=\bfseries] at (0, -2.3) {Conscience\\spontanée};
\node[align=center, text=popdark, font=\bfseries] at (0, -0.5) {Conscience\\réfléchie};
\node[align=center, text=popdark, font=\bfseries] at (0, 0) {Conscience\\morale};

% Flèche explicative
\draw[->, thick, popdark] (3.5, -2.5) -- (3.0, -1.5);
\node[font=\footnotesize, text=popmuted, align=left, anchor=west] at (3.7, -2.5) {Du plus fondamental\\au plus élaboré};
\end{tikzpicture}
\end{center}
\legende{Les trois types de conscience comme cercles concentriques : la conscience spontanée (présence à soi, cercle large) est le fondement ; la conscience réfléchie (examen de soi, cercle moyen) l'inclut et la dépasse ; la conscience morale (jugement de valeur, cercle central) les couronne en les orientant vers le bien.}
\end{popfigure}

\begin{center}
\begin{tabularx}{\linewidth}{|l|X|X|X|}
\hline
\rowcolor{popblueL} \textbf{Type} & \textbf{Définition} & \textbf{Exemple de la vie courante (Cameroun)} & \textbf{Philosophe associé} \\
\hline
Conscience spontanée & Présence immédiate à soi et aux choses, sans réflexion & Je conduis ma moto à Bafoussam en évitant les nids-de-poule sans y penser & Bergson (la conscience comme durée vécue) \\
\hline
Conscience réfléchie & Retour du sujet sur lui-même, examen de ses pensées et de ses actes & Le soir, je repasse ma journée et je me demande pourquoi j'ai répondu mal à mon père & Descartes (\og je pense, donc je suis \fg{}) \\
\hline
Conscience morale & Jugement du bien et du mal de ses propres actions & Après avoir menti à un ami, j'éprouve du remords et je décide de lui dire la vérité & Kant (le devoir), Rousseau (la conscience, voix de l'âme) \\
\hline
\end{tabularx}
\end{center}

\courssub{Conscience psychologique et conscience morale : une distinction essentielle}

Il faut maintenant distinguer clairement deux grandes familles. La \cle{conscience psychologique} regroupe la conscience spontanée et la conscience réfléchie : elle concerne la \emph{connaissance} (que se passe-t-il en moi et autour de moi ?). La \cle{conscience morale}, elle, concerne la \emph{valeur} (ce que je fais est-il bien ou mal ?). On peut avoir une conscience psychologique très vive et une conscience morale endormie : le cambrioleur habile sait parfaitement ce qu'il fait (conscience psychologique éveillée), mais il ne juge pas son acte comme mauvais (conscience morale absente ou dévoyée). Inversement, une personne simple, peu portée à l'analyse d'elle-même, peut avoir un sens moral très délicat.

\begin{exoresolu}[Appliquer les types de conscience à une situation concrète]
\textbf{Énoncé.} Un candidat au Probatoire technique, à Yaoundé, hésite à recopier la réponse de son voisin pendant l'épreuve. D'un côté, il a peur d'échouer ; de l'autre, une voix intérieure lui dit que tricher est mal. Identifiez, dans cette situation, les trois types de conscience et montrez le conflit qui s'y joue.
\tcblower
\textbf{Solution détaillée.}
\begin{enumerate}
\item \emph{Conscience spontanée :} le candidat est présent à la situation — il sent sa peur de l'échec, il voit la copie du voisin, il perçoit la présence du surveillant. C'est la conscience immédiate de ce qu'il vit.
\item \emph{Conscience réfléchie :} il se retourne sur lui-même et examine la situation : \og si je triche et que je suis pris, je suis exclu ; si je ne triche pas, je risque l'échec. \fg{} Il pèse, calcule, délibère : c'est le retour du sujet sur ses pensées et ses actes possibles.
\item \emph{Conscience morale :} au-delà du calcul des risques, une voix intérieure juge : \og tricher est mal, c'est un manque d'honnêteté envers moi-même et envers les autres candidats. \fg{} C'est le tribunal intérieur qui évalue l'acte selon le bien et le mal.
\end{enumerate}
\emph{Conclusion :} le conflit vécu par le candidat oppose la conscience spontanée (la peur immédiate d'échouer, le désir de réussir à tout prix) à la conscience morale (le devoir d'honnêteté). La conscience réfléchie est le théâtre de ce débat intérieur. S'il choisit l'honnêteté, il aura la \emph{bonne conscience} ; s'il triche, il risque le \emph{remords}, signe que la conscience morale continue de juger après l'acte.
\end{exoresolu}

\begin{savaistu}[Conscience et responsabilité en droit camerounais]
En droit pénal camerounais, comme dans la plupart des systèmes juridiques, on ne peut condamner que quelqu'un qui avait conscience de ce qu'il faisait : c'est l'exigence du \cle{discernement}. Ainsi, le Code pénal prévoit que le fou, qui a agi sans conscience de ses actes, n'est pas pénalement responsable : il sera soigné, non puni. De même, la responsabilité pénale des mineurs est aménagée selon leur capacité de discernement. Le droit rejoint ici la philosophie : punir suppose que le coupable pouvait se représenter son acte et en juger la valeur. Sans conscience, pas de culpabilité.
\end{savaistu}

\courssub{Bilan de la section}

Nous avons vu que le mot \emph{conscience}, issu du latin \emph{cum-scire} (\og connaître avec \fg{}), désigne la connaissance immédiate que le sujet a de lui-même, de ses états, de ses actes et du monde. Cette conscience se déploie en trois types : la conscience \emph{spontanée} (présence simple à soi), la conscience \emph{réfléchie} (retour sur soi, dont le Cogito de Descartes est le modèle) et la conscience \emph{morale} (tribunal intérieur jugeant le bien et le mal). Mais cette conscience est-elle toute-puissante ? Suffit-elle à expliquer nos actes ? L'élève de notre situation d'accroche, qui a \og oublié \fg{} trois fois ses devoirs, était-il vraiment maître de son oubli ? C'est à ces questions — la valeur, puis les limites de la conscience — que répondront les sections suivantes.

\coursec{Caractéristiques de la conscience}

\courssub{Transition : de la définition aux structures vécues}

Dans la section précédente, nous avons distingué la conscience morale (jugement de valeur), la conscience psychologique (connaissance immédiate) et la conscience réflexive (retour sur soi). Nous nous attachons maintenant à décrire les \textbf{structures invariantes} qui font de la conscience ce qu'elle est, quel que soit son type. Ces caractéristiques ne sont pas des propriétés accidentelles : elles constituent l'essence même de la vie consciente. Comme l'écrit Husserl, « toute conscience est conscience de quelque chose » — c'est par là que nous commençons.

\courssub{L'intentionnalité : la conscience comme visée}

\begin{definition}[Intentionnalité (Husserl)]
L'\cle{intentionnalité} est la propriété fondamentale de la conscience d'être toujours \emph{conscience de quelque chose}. Elle désigne la structure par laquelle tout acte psychique (percevoir, juger, désirer, se souvenir, imaginer) vise un objet, qu'il existe réellement ou non. L'objet visé est l'\cle{objet intentionnel}, distinct de l'objet réel.
\end{definition}

\textbf{Intuition première.}\par Quand je vois l'arbre devant moi, ma vue n'est pas un état clos sur lui-même : elle \emph{file} vers l'arbre. Quand je me souviens de mon enfance, mon souvenir \emph{vise} un passé qui n'est plus présent. Quand je crains l'échec, ma peur \emph{porte} sur un avenir possible. Dans tous les cas, la conscience sort d'elle-même : elle est \emph{extase} (ek-stasis, « sortie hors de soi »).

\textbf{Analyse rigoureuse.}\par Contre le psychologisme qui réduit la conscience à un simple ensemble d'états mentaux juxtaposés, Husserl montre que l'intentionnalité n'est pas une relation entre deux choses réelles (un sujet et un objet), mais la \emph{structure même} de l'acte conscient. L'objet intentionnel est \emph{constitutif} de l'acte : il n'y a pas de perception sans perçu, pas de jugement sans jugé. La \cle{réduction phénoménologique} (epoché) consiste à mettre entre parenthèses l'existence de l'objet pour décrire purement la manière dont il est donné à la conscience.

\begin{exemplebox}[L'intentionnalité au volant]
Le chauffeur de taxi à Douala conduit « machinalement » tout en pensant à son repas du soir. Sa conscience de la route est une \emph{conscience marginale} (arrière-plan) : elle vise la trajectoire, les feux, les piétons, mais sans thématisation explicite. Sa conscience du repas est une \emph{conscience thématique} (premier plan) : elle vise un objet absent (le futur plat). Les deux sont intentionnelles : l'une vise un présent perçu, l'autre un avenir imaginé. Si un enfant traverse soudain, la conscience de la route devient thématique instantanément — la visée intentionnelle se recentre.
\end{exemplebox}

\begin{attention}[Piège : intentionnalité $\neq$ intention volontaire]
En langage courant, « intention » signifie « projet délibéré ». En philosophie, l'intentionnalité concerne \emph{tout} acte conscient, y compris les passions, les perceptions involontaires, les rêves. La peur intentionnelle n'est pas une peur « voulue ». Ne confondez pas \cle{intentionnalité} (structure de la conscience) et \cle{volonté} (faculté pratique).
\end{attention}

\courssub{L'unité et l'identité : le « moi » comme pôle unificateur}

\begin{propriete}[Unité et flux : la tension fondamentale]
La conscience présente une double structure paradoxale :
\begin{enumerate}
\item \textbf{Unité synthétique :} à tout instant, mes perceptions, pensées, affects, souvenirs forment \emph{une seule} conscience, rattachée à \emph{un seul} « je » (l'\cle{ego} pôle d'identité chez Husserl, le \emph{moi} empirique en psychologie). Cette unité permet l'identité personnelle dans le temps.
\item \textbf{Flux continu (Bergson, James) :} la conscience n'est pas une substance figée, mais un \emph{devenir} perpétuel. « La conscience est un flux » (William James) : les états se succèdent, se pénètrent, sans coupure nette.
\end{enumerate}
La tension entre \emph{unité} (identité) et \emph{flux} (différence) est le cœur de la subjectivité.
\end{propriete}

\textbf{Expérience vécue.}\par Fermez les yeux. Vous entendez le bruit de la pluie, vous sentez le contact du siège, vous pensez à ce cours, une émotion passe. Tout cela est \emph{votre} conscience maintenant — unité. Mais dans une seconde, la pensée aura changé, le bruit aura varié — flux. Pourtant, c'est le \emph{même} « je » qui les traverse.

\begin{exoresolu}[Analyse de l'unité dans la distraction]
\textbf{Énoncé :} Le chauffeur de taxi conduit machinalement et pense à son repas. Y a-t-il une ou deux consciences ?
\tcblower
\textbf{Solution :} Il n'y a qu'\emph{une} conscience (unité du moi), mais avec \emph{deux champs thématiques} distincts : un champ focal (le repas) et un champ marginal (la conduite). L'unité se manifeste par la capacité de bascule instantanée : si un danger survient, le champ marginal devient focal sans rupture du « je ». C'est la structure \emph{figure/fond} (forme/fond) appliquée à la conscience.
\end{exoresolu}

\courssub{La temporalité : la conscience comme synthèse du temps}

La conscience ne se trouve pas \emph{dans} le temps comme un objet dans une boîte ; elle \emph{constitue} le temps vécu. Husserl analyse la \cle{conscience interne du temps} comme une triple synthèse :

\begin{popfigure}
\begin{tikzpicture}[>=Stealth, font=\small]
% Axe temporal
\draw[popdark, thick, ->] (-0.5,0) -- (13,0) node[right] {Flux de la conscience};
% Zones
\fill[popblue!15] (-0.3,-1.2) rectangle (4.5,1.2);
\fill[popgreen!15] (4.5,-1.2) rectangle (8.5,1.2);
\fill[poporange!15] (8.5,-1.2) rectangle (12.8,1.2);
% Étiquettes zones
\node[popblue, align=center] at (2.1,1.5) {\textbf{Rétention}\\(Passé immédiat\\encore retenu)};
\node[popgreen, align=center] at (6.5,1.5) {\textbf{Impression originaire}\\(Perception présente,\\Maintenant vif)};
\node[poporange, align=center] at (10.65,1.5) {\textbf{Protention}\\(Attente, Projet,\\Avenir anticipé)};
% Flèches internes
\draw[popblue, thick, <->] (1.5,-0.8) -- (3.5,-0.8) node[midway, below] {Retenir};
\draw[popgreen, thick, <->] (5.5,-0.8) -- (7.5,-0.8) node[midway, below] {Vivre};
\draw[poporange, thick, <->] (9.5,-0.8) -- (11.5,-0.8) node[midway, below] {Anticiper};
% Lien unification
\draw[poppurple, very thick, decorate, decoration={brace, amplitude=5pt, mirror}] (-0.3,-1.4) -- (12.8,-1.4) node[midway, below=8pt, poppurple] {\textbf{Unité de la conscience (Ego) synthétisant le temps}};
% Point "Maintenant"
\draw[popgreen, dashed, thick] (6.5,-1.2) -- (6.5,1.2);
\node[popgreen] at (6.5, -1.7) {Point d'origine (maintenant)};
\end{tikzpicture}
\legende{Structure tripartite de la conscience temporelle selon Husserl : rétention (passé retenu), impression originaire (présent vif), protention (avenir anticipé). L'ego unifie ce flux.}
\end{popfigure}

\begin{itemize}
\item \textbf{Rétention} : non pas un rappel de mémoire (représentation d'un passé lointain), mais le \emph{maintien} de l'impression qui vient de s'éteindre (la note de musique qui « dure » encore dans l'écoute).
\item \textbf{Impression originaire} : le « maintenant » vivant, point d'ancrage, jamais saisi comme objet (on ne peut pas « avoir » le présent, on ne peut que le \emph{vivre}).
\item \textbf{Protention} : l'anticipation vide mais déterminée (j'attends la note suivante, je tends mon action vers le but).
\end{itemize}

\begin{exemplebox}[Temporalité du chauffeur]
Le chauffeur \emph{retient} la trajectoire qu'il vient de suivre (rétention corporelle), \emph{perçoit} le feu qui passe au vert (impression originaire), \emph{anticipe} le freinage au carrefour suivant (protention). Son projet de repas (futur personnel) s'entrelace avec la protention de conduite (futur immédiat). La conscience tisse ainsi plusieurs temporalités.
\end{exemplebox}

\courssub{La spontanéité et la réflexivité : la conscience se prend pour objet}

\begin{definition}[Réflexivité]
La \cle{réflexivité} est la capacité de la conscience à se retourner sur elle-même pour se prendre comme objet : « Je sais que je sais », « Je pense que je pense ». Elle transforme la conscience \emph{vécue} (pré-réflexive) en conscience \emph{thématisée} (réfléchie).
\end{definition}

\textbf{Spontanéité.}\par La conscience n'est pas seulement réceptive (passivité de la perception) ; elle est \emph{spontanée} : elle initie des actes (juger, décider, imaginer, douter). Chez Sartre, la conscience est \emph{néantisation} : elle introduit le néant dans l'être par la question, le doute, le projet.

\textbf{Niveaux de réflexivité.}\par
\begin{enumerate}
\item \textbf{Conscience pré-réflexive (ou non-positionnelle) :} conscience \emph{de} l'objet sans conscience thématique \emph{de soi} (le chauffeur conduit « sans y penser »).
\item \textbf{Conscience réflexive positionnelle :} je me dis « je conduis » — le « je » devient objet.
\item \textbf{Conscience réflexive ultérieure :} je me dis « je me dis que je conduis » — régression à l'infini possible mais non réalisée.
\end{enumerate}

\begin{attention}[Piège : réflexivité $\neq$ introspection psychologique]
L'introspection est une \emph{méthode} (regarder à l'intérieur de soi). La réflexivité est une \emph{structure} : la conscience est \emph{par essence} conscience d'elle-même (Sartre : « La conscience est conscience d'être conscience »). On ne « fait » pas de la réflexivité ; on \emph{est} réflexif.
\end{attention}

\courssub{Le caractère privé : l'accès de première personne}

\begin{propriete}[Privilège épistémique de la première personne]
Seul le sujet a un accès \emph{direct}, \emph{immédiat}, \emph{incorrigible} (on ne peut pas se tromper sur le fait que l'on a mal, même si on se trompe sur la cause) à sa propre conscience. Autrui n'accède à ma conscience que par \emph{inférence} (comportement, langage, empathie) — accès médiat, incertain, public.
\end{propriete}

\textbf{Conséquence : le problème des autres esprits.}\par Si la conscience est essentiellement privée, comment savoir que les autres sont conscients ? L'argument par analogie (je suis conscient et j'ai un corps ; autrui a un corps semblable, donc il est conscient) est faible, car il repose sur un seul cas : moi. Pour la phénoménologie (Husserl, Lévinas), l'altérité est donnée \emph{originairement} dans l'expérience de la rencontre (le visage d'autrui), non déduite par raisonnement.

\begin{exemplebox}[Privacité dans le taxi]
Le chauffeur \emph{éprouve} sa fatigue, son envie de manger, son agacement contre le trafic. Le passager ne voit que : mains sur le volant, regard sur la route, éventuellement un soupir. Le passager \emph{infère} : « Il est fatigué ». Le chauffeur \emph{sait} : « Je suis fatigué ». L'écart est infranchissable par l'observation externe.
\end{exemplebox}

\courssub{La conscience comme présence et comme éveil : les degrés de la conscience}

La conscience n'est pas un tout-ou-rien. Elle admet des \cle{degrés} (niveaux d'intensité, de clarté, d'étendue). Cette gradation rejoint la physiologie (vigilance) mais la dépasse par la description de l'expérience vécue.

\begin{popfigure}
\begin{tikzpicture}
\begin{axis}[
    ybar,
    bar width=18pt,
    width=\linewidth,
    height=8cm,
    ymin=0, ymax=100,
    ylabel={Niveau de clarté de la conscience (échelle qualitative)},
    symbolic x coords={Coma, Sommeil profond, Rêve, Veille distraite, Attention concentrée},
    xtick=data,
    x tick label style={rotate=30, anchor=east, font=\small},
    ymajorgrids=true,
    grid style={popline},
    nodes near coords,
    nodes near coords align={vertical},
    every node near coord/.append style={font=\small, popdark},
    title style={font=\normalsize\bfseries, popdark},
    title={Degrés de la conscience : de l'absence à la plénitude},
    bar shift=0pt,
]
\addplot[popblue!70, fill=popblue!50] coordinates {
    (Coma, 5)
    (Sommeil profond, 15)
    (Rêve, 45)
    (Veille distraite, 65)
    (Attention concentrée, 95)
};
\end{axis}
\end{tikzpicture}
\legende{Échelle qualitative (et non mesurée) des degrés de conscience. Coma : conscience quasi nulle. Sommeil profond : conscience éteinte, sans souvenir au réveil. Rêve : conscience présente mais déconnectée du réel. Veille distraite (conduite « automatique ») : conscience marginale large, focale étroite. Attention concentrée : conscience focale intense.}
\end{popfigure}

\begin{savaistu}[Un tiers de vie en conscience altérée]
Un adulte dort en moyenne \num{8} heures sur \num{24} : il passe \textbf{un tiers de sa vie} en sommeil (conscience altérée, discontinue). Sur \num{80} ans, cela représente environ \num{26,7} années de « non-conscience » ordinaire. La conscience \emph{pleine} (veille attentive) n'occupe qu'une fraction du temps d'éveil. \textbf{La conscience n'est pas un attribut permanent du sujet, mais un processus intermittent.}
\end{savaistu}

\begin{attention}[Conscience (philosophie) $\neq$ Vigilance (médecine)]
La \emph{vigilance} est un état physiologique (éveil cortical, mesurable par électroencéphalogramme). La \emph{conscience} au sens philosophique est l'\emph{expérience vécue} (intentionnalité, réflexivité). Un patient en état végétatif peut présenter des cycles veille/sommeil (vigilance) sans conscience au sens plein (pas de contenu intentionnel accessible). Inversement, le rêveur a une conscience vive (le rêve) sans vigilance éveillée. Ne confondez pas le \emph{substrat physiologique} et le \emph{phénomène vécu}.
\end{attention}

\courssub{Synthèse des caractéristiques : tableau récapitulatif}

\begin{center}
\begin{tabularx}{\linewidth}{|l|X|X|}\hline
\rowcolor{popblueL}
\textbf{Caractéristique} & \textbf{Énoncé essentiel} & \textbf{Auteur / Référence clé} \\ \hline
Intentionnalité & Toute conscience vise un objet (réel ou irréel) & Husserl \\ \hline
Unité / Flux & Un seul « je » pour un flux continu d'états & James, Bergson, Husserl \\ \hline
Temporalité & Synthèse rétention -- impression -- protention & Husserl \\ \hline
Spontanéité / Réflexivité & La conscience s'initie et se prend pour objet & Sartre \\ \hline
Privacité & Accès direct, incorrigible, de première personne & Descartes \\ \hline
Degrés (Éveil) & Continuum du coma à l'attention maximale & Bergson \\ \hline
\end{tabularx}
\end{center}

\begin{methode}[Analyser une caractéristique de la conscience]
Pour traiter une caractéristique dans une dissertation ou un commentaire :
\begin{enumerate}
\item \textbf{Identifiez l'expérience vécue} qui la manifeste (ex. : « Je me souviens » $\rightarrow$ temporalité / rétention).
\item \textbf{Nommez la structure} avec le terme technique exact (intentionnalité, protention, réflexivité positionnelle...).
\item \textbf{Citez l'auteur de référence} (Husserl pour l'intentionnalité et la temporalité, Sartre pour la réflexivité, Bergson pour le flux et les degrés).
\item \textbf{Problématisez} : cette caractéristique est-elle une certitude (Descartes) ou une construction (Husserl) ? Est-elle absolue (privacité) ou relative (degrés) ?
\item \textbf{Illustrez par un exemple concret} (le chauffeur, le musicien, l'élève en examen) pour ancrer l'abstraction.
\end{enumerate}
\end{methode}

\begin{exoresolu}[Application de la méthode : « La conscience est-elle une ? »]
\textbf{Énoncé :} À partir de l'expérience du chauffeur de taxi, discutez l'unité de la conscience.
\tcblower
\textbf{Solution :}
\begin{enumerate}
\item \textbf{Expérience :} conduite automatique + pensée sur le repas. Deux contenus, une seule personne.
\item \textbf{Structure :} \cle{champ de conscience} avec \emph{focal} (repas) et \emph{marginal} (route). Unité structurelle (un ego), pluralité des actes intentionnels.
\item \textbf{Auteurs :} \emph{James} (flux unifié), \emph{Husserl} (ego pôle d'identité), \emph{Janet} (désagrégation possible : névrose, hypnose).
\item \textbf{Problématisation :} l'unité est-elle \emph{donnée} (substance cartésienne) ou \emph{construite} (synthèse husserlienne) ? La distraction montre une \emph{unité fonctionnelle} (basculement possible) mais une \emph{disjonction thématique}. L'unité n'est pas l'uniformité.
\item \textbf{Conclusion :} la conscience est \emph{unité synthétique d'un flux hétérogène} — identité dans la différence.
\end{enumerate}
\end{exoresolu}

\begin{savaistu}[Le « chauffeur automatique » et les neurosciences]
Ce que le philosophe appelle « conscience marginale » ou « habitude corporelle » (Merleau-Ponty : \emph{corps propre}, \emph{habitude motrice}), les neurosciences le modélisent par l'\emph{automatisme procédural} (circuits de l'habitude, cervelet). Le chauffeur ne « pense pas » à changer de vitesse : son corps \emph{sait}. La philosophie décrit l'expérience vécue (intentionnalité marginale) ; les neurosciences expliquent le mécanisme. Les deux niveaux sont complémentaires, non réductibles l'un à l'autre.
\end{savaistu}

\courssub{Transition vers la section suivante : valeur de la conscience}

Nous avons décrit \emph{comment} la conscience est structurée (intentionnelle, unitaire, temporelle, réflexive, privée, graduée). Nous devons maintenant nous demander : \emph{que vaut} cette conscience ? Est-elle source de vérité (critère cartésien) ou d'illusion ? Est-elle le fondement de la liberté et de la responsabilité morale ? C'est l'objet de la section suivante : \textbf{La valeur de la conscience}.

\coursec{Valeur de la conscience}

Après avoir analysé les \emph{caractéristiques} de la conscience — son unité, sa transparence, son intentionalité et sa réflexivité — nous devons maintenant nous demander : \emph{à quoi sert la conscience ?} Quelle est sa \textbf{valeur} ? Non pas sa valeur marchande, bien sûr, mais sa valeur \emph{ontologique} (ce qu'elle nous fait être), \emph{épistémologique} (ce qu'elle nous fait savoir) et \emph{axiologique} (ce qu'elle nous fait valoir). C'est parce que la conscience possède ces caractéristiques qu'elle peut jouer le rôle de fondement : c'est \emph{en étant} une lumière tournée vers soi qu'elle permet la connaissance de soi ; c'est \emph{en étant} intentionnelle et délibérative qu'elle rend possible la liberté.

\courssub{La conscience comme connaissance de soi : l'héritage socratique}

\emph{« Connais-toi toi-même »} (\emph{Gnôthi seauton}). Cette injonction, gravée au fronton du temple d'Apollon à Delphes, n'est pas une simple maxime morale : elle pose la conscience comme la \textbf{condition de toute sagesse}. Pour Socrate, l'ignorance de soi est la racine de tous les maux ; l'homme qui ne se connaît pas prend ses désirs pour des besoins, ses opinions pour des vérités, ses préjugés pour des jugements. La conscience, comme \emph{retour réflexif sur soi}, est l'unique voie pour discerner le vrai du faux, le bien du mal, l'essentiel de l'accessoire.

\begin{savaistu}[L'oracle de Delphes]
La devise « Connais-toi toi-même » (\emph{Gnôthi seauton}) était inscrite au fronton du temple d'Apollon à Delphes, en Grèce antique. Elle résumait la sagesse delphique : l'homme ne peut accéder à la vérité sur le monde qu'en commençant par se connaître lui-même. Socrate en fit le principe de sa méthode : le philosophe n'est pas celui qui sait, mais celui qui \emph{cherche} à se connaître, et qui aide les autres à accoucher de leur propre vérité (maïeutique).
\end{savaistu}

Mais cette connaissance de soi n'est pas un miroir passif. Elle est \emph{activité} : examen de conscience, dialogue intérieur, mise à l'épreuve de ses motifs. C'est ce que Socrate appelle le « soin de l'âme » (\emph{epiméleia tês psychês}). Sans conscience, l'homme est un étranger pour lui-même ; avec elle, il devient \emph{sujet} de sa vie.

\courssub{Descartes : la conscience comme certitude première et fondement de toute vérité (le Cogito)}

Au XVII\textsuperscript{e} siècle, René Descartes radicalise cette intuition. Dans ses \emph{Méditations métaphysiques}, il met en doute tout : le monde extérieur, son corps, ses sens, même les vérités mathématiques (hypothèse du \emph{malin génie}). Mais une chose résiste au doute hyperbolique : \emph{je doute, donc je pense, donc je suis}. Cette certitude — le \textbf{Cogito} — n'est pas une déduction logique (un syllogisme), mais une \emph{intuition immédiate} que la conscience saisit en elle-même.

\begin{definition}[Le Cogito cartésien]
« Je pense, donc je suis » (\emph{Cogito, ergo sum}) : vérité indubitable que la conscience saisit en elle-même, sans médiation d'aucun raisonnement. Pour Descartes, le \emph{cogito} est le \textbf{premier principe} de la philosophie : il est certain parce qu'il est \emph{clair et distinct}, perçu par la conscience dans son acte même de penser. Tout le reste (Dieu, le monde, les sciences) ne sera reconstruit qu'à partir de ce fondement.
\end{definition}

La valeur épistémologique de la conscience est ici suprême : elle est \emph{l'archimède} sur lequel tout l'édifice du savoir peut s'appuyer. Mais attention : ce « je » n'est pas une substance mystérieuse ; c'est l'acte même de conscience (\emph{res cogitans}). La conscience n'est pas \emph{un} objet parmi d'autres : elle est la \emph{condition de possibilité} de tout objet.

\begin{attention}[Piège fréquent : confondre Cogito et syllogisme]
Beaucoup d'élèves écrivent : « Je pense, donc je suis : c'est un syllogisme (tout ce qui pense existe ; je pense ; donc je suis) ». \textbf{Faux.} Descartes précise expressément (\emph{Réponses aux secondes objections}) que le Cogito n'est \emph{pas} un raisonnement déductif, mais une \emph{intuition simple et immédiate} : « Quand quelqu'un dit "je pense, donc je suis", il ne déduit pas l'existence de la pensée par un syllogisme, mais il la reconnaît comme une chose connue par elle-même. » Ne confondez pas \emph{intuition intellectuelle} et \emph{inférence logique}.
\end{attention}

\courssub{La conscience comme condition de la liberté : délibération et choix}

Si la conscience nous donne la vérité sur nous-mêmes, elle nous donne aussi le \emph{pouvoir} sur nous-mêmes. \textbf{Être conscient, c'est pouvoir choisir.} Un automate, une plante, un animal guidé par le seul instinct ne \emph{délibèrent} pas : ils subissent. L'homme conscient, lui, \emph{représente} des fins possibles, les \emph{pèse}, les \emph{compare} et \emph{décide}. Cette capacité de \emph{délibération} — mettre à distance ses désirs immédiats pour les évaluer à la lumière de raisons — est la marque de la \textbf{liberté}.

\begin{propriete}[Conscience et liberté]
Sans conscience, pas de liberté. La liberté suppose la capacité de se représenter des alternatives, de les évaluer et de s'engager dans l'une d'elles par un acte de volonté. L'inconscient (pulsions, réflexes, conditionnements) détermine ; la conscience \emph{libère} en instaurant un espace de \emph{délibération} entre le stimulus et la réponse.
\end{propriete}

Cet espace est ce que les philosophes appellent la \emph{puissance de se déterminer soi-même} (Spinoza, Kant). La conscience ne \emph{garantit} pas que nous choisirons bien (on peut user de sa liberté pour le mal), mais elle \emph{rend possible} le choix. C'est pourquoi l'aliénation (addiction, manipulation, aliénation sociale) se définit souvent comme une \emph{perte de conscience} : on agit sans se savoir agir, on devient spectateur de sa propre vie.

\begin{exemplebox}[La délibération morale : un exemple camerounais]
Un jeune ingénieur à Douala reçoit une proposition : falsifier un rapport d'impact environnemental pour un projet minier, en échange d'une forte somme. Son \emph{premier mouvement} (inconscient, pulsionnel) est l'appât du gain. Mais sa \emph{conscience} intervient : il se représente les conséquences (pollution du fleuve Wouri, maladies pour les riverains, violation de la loi), il entend la voix de sa formation d'ingénieur, il se souvient du serment professionnel. Il \emph{délibère} : le gain immédiat vs la responsabilité à long terme. C'est \emph{parce qu'il est conscient} qu'il peut dire \emph{non}. Un robot programmé pour « maximiser le profit » n'aurait pas cette liberté.
\end{exemplebox}

\courssub{La conscience comme fondement de la responsabilité morale et juridique}

La liberté n'est pas une fin en soi : elle engage. \textbf{Qui est libre est responsable.} La responsabilité (\emph{respondere} : répondre de ses actes) suppose que l'agent \emph{savait} ce qu'il faisait, \emph{voulait} le faire, et \emph{pouvait} faire autrement. Or, cette triple condition — connaissance, volonté, pouvoir — est exactement ce que la conscience assure.

\begin{propriete}[Conscience et responsabilité]
Sans conscience, pas de responsabilité — on ne peut être tenu pour coupable d'un acte accompli sans discernement. Le droit comme la morale distinguent l'\emph{acte humain} (conscient, libre, imputable) de l'\emph{acte de l'homme} (réflexe, contrainte, folie, sommeil). Seuls les actes conscients engagent la responsabilité.
\end{propriete}

C'est ce que le droit camerounais, comme la plupart des systèmes juridiques, consacre. Le Code pénal (Loi n°2016/007 du 12 juillet 2016) prévoit l'\emph{irresponsabilité pénale} pour trouble mental abolissant le discernement (art. 73) et l'\emph{atténuation} pour trouble partiel. Le juge examine : l'accusé \emph{avait-il conscience} de la portée de son acte ? \emph{Pouvait-il} se déterminer autrement ?

\begin{exemplebox}[Le tribunal et la valeur juridique de la conscience]
Un homme tue son rival dans une dispute. Deux scénarios :
\begin{enumerate}
\item \textbf{Crime prémédité :} Il a acheté une arme, guetté la victime, tiré de sang-froid. La conscience est \emph{pleine} : il a représenté l'acte, l'a voulu, l'a planifié. \emph{Responsabilité maximale} (assassinat, peine capitale possible).
\item \textbf{Accident / Légitime défense / Folie :} Il a poussé l'autre qui est tombé mortellement (accident) ; ou il a riposté à une agression mortelle (légitime défense) ; ou il était en crise psychotique aigüe, croyant tuer un démon (folie). La conscience est \emph{absente, altérée ou contrainte}. \emph{Responsabilité nulle ou atténuée}.
\end{enumerate}
Le tribunal \emph{applique la valeur juridique de la conscience} : il juge non pas l'acte matériel (un mort), mais l'\emph{état de conscience} de l'auteur.
\end{exemplebox}

\courssub{La conscience comme source de la dignité humaine : l'homme comme fin (Kant)}

Pour Emmanuel Kant, la conscience morale (la \emph{raison pratique}) est ce qui élève l'homme au-dessus de la nature. Dans les \emph{Fondements de la métaphysique des mœurs}, il écrit : \emph{« L'homme, et généralement tout être rationnel, existe comme fin en soi, non pas seulement comme moyen. »} Pourquoi ? Parce que la conscience permet à l'homme de \emph{se représenter des fins}, de se \emph{donner sa propre loi} (autonomie), de ne pas être \emph{seulement} poussé par des causes naturelles (désirs, instincts, conditionnements) mais d'agir \emph{par devoir}, c'est-à-dire par respect pour la loi morale qu'il se prescrit à lui-même.

\begin{aretenir}[Dignité et conscience chez Kant]
\begin{itemize}
\item \textbf{Prix vs Dignité :} Ce qui a un \emph{prix} peut être remplacé (marchandises, outils). Ce qui a une \emph{dignité} n'a pas d'équivalent, il est \emph{fin en soi}.
\item \textbf{Source de la dignité :} La \emph{conscience morale} (raison pratique) qui permet l'\emph{autonomie} : l'homme se donne la loi morale à lui-même, au lieu de la subir de l'extérieur (hétéronomie).
\item \textbf{Conséquence :} Traiter un homme \emph{uniquement comme moyen} (esclavage, exploitation, manipulation, torture) nie sa conscience et viole sa dignité.
\end{itemize}
\end{aretenir}

Cette thèse fonde les \textbf{Droits de l'Homme} modernes : la \emph{Déclaration universelle de 1948} (art. 1) proclame que « tous les êtres humains naissent libres et égaux en dignité et en droits. Ils sont doués de raison et de conscience ». La conscience n'est pas un luxe : elle est le \emph{titre de noblesse} de l'humanité.

\courssub{Application : au Cameroun, l'objection de conscience et la liberté de conscience garanties par la Constitution}

Le Cameroun, dans sa \textbf{Constitution du 18 janvier 1996} (Préambule), inscrit explicitement la valeur juridique et politique de la conscience :

\begin{quote}
\emph{« La liberté de conscience, de culte, d'opinion et d'expression est garantie à tout citoyen, sous réserve du respect des libertés d'autrui et de l'ordre public. »} (Préambule, al. 3)

\emph{« Nul ne peut être contraint d'accomplir un acte contraire à sa conscience. »} (Préambule, al. 12 — droit à l'objection de conscience)
\end{quote}

Ces dispositions ont une portée concrète majeure :
\begin{itemize}
\item \textbf{Liberté de conscience :} droit de croire ou de ne pas croire, de changer de religion, de manifester ses convictions (port du voile, jeûne, prières, symboles religieux) sans discrimination.
\item \textbf{Objection de conscience :} un médecin peut refuser de pratiquer un avortement (sauf danger vital pour la mère), un militaire peut refuser un ordre manifestement illégal (crime de guerre), un fonctionnaire peut refuser de signer un acte contraire à ses convictions profondes (corruption, fraude électorale).
\item \textbf{Laïcité de l'État :} l'État ne favorise aucune religion, garantit l'égalité devant la loi indépendamment des convictions.
\end{itemize}

\begin{exoresolu}[Cas pratique : l'objection de conscience d'un infirmier]
\textbf{Énoncé :} Paul, infirmier dans un hôpital public à Yaoundé, refuse de participer à une intervention d'avortement thérapeutique (autorisé par la loi pour sauver la vie de la mère), invoquant sa foi catholique et l'article 12 du Préambule de la Constitution. Le directeur de l'hôpital menace de le sanctionner pour « insubordination » et « refus de service ». Paul a-t-il raison sur le plan juridique ?
\tcblower
\textbf{Solution détaillée :}
\begin{enumerate}
\item \textbf{Identification du droit applicable :} Préambule de la Constitution de 1996, al. 12 : « Nul ne peut être contraint d'accomplir un acte contraire à sa conscience. » Loi n°90/053 du 19 décembre 1990 relative à la liberté d'association (et jurisprudence du Conseil constitutionnel).
\item \textbf{Analyse :} L'objection de conscience est un \emph{droit fondamental} attaché à la dignité de la personne. Elle s'applique aux actes qui heurtent gravement les convictions intimes (vie, mort, intégrité corporelle). L'avortement, même thérapeutique, touche à la vie : l'objection est recevable.
\item \textbf{Limitations :} Le droit à l'objection ne doit pas mettre en danger la vie d'autrui (ici, la mère). Mais l'hôpital a l'obligation d'\emph{organiser le remplacement} (un autre soignant non objecteur) pour garantir le droit à la santé de la patiente. La sanction pour insubordination serait \emph{inconstitutionnelle} si l'objection est sincère et si le service est assuré par un collègue.
\item \textbf{Conclusion :} Paul a raison. Le directeur doit respecter son objection et organiser la continuité des soins. Toute sanction serait annulée par le juge administratif.
\end{enumerate}
\end{exoresolu}

\courssub{La conscience comme lumière intérieure guidant l'action : la tradition morale africaine}

La philosophie occidentale n'a pas le monopole de la valeur de la conscience. Dans les \textbf{traditions morales africaines} (Bamiléké, Bassa, Béti, Douala, Peul, etc.), la conscience est souvent désignée par le terme de \emph{cœur} (\emph{tí} en bamiléké, \emph{ngón} en bassa, \emph{zúm} en béti). Ce « cœur » n'est pas l'organe pompe, mais le \emph{siège de la pensée, de la volonté, du discernement moral}.

\begin{exemplebox}[Le « cœur » juge dans la tradition bamiléké]
Chez les Bamiléké, on dit : \emph{« Tí wú fó »} — « Le cœur ne ment pas. » Avant un acte important (mariage, alliance, guerre, jugement), les anciens demandent : \emph{« Ton cœur est-il pur ? »} (\emph{Tí wú fó ?}). Celui qui agit avec un « cœur trouble » (\emph{tí wú kù}) — c'est-à-dire avec arrière-pensée, hypocrisie, intention cachée — perd sa \emph{dignité} (\emph{ńkù}) et son \emph{autorité morale}. Le « cœur » est cette \emph{conscience intérieure} qui juge l'acte \emph{avant} qu'il ne soit posé, et qui continue de juger \emph{après} (remords ou paix intérieure). C'est exactement la fonction de la \emph{conscience morale} décrite par Kant, mais ancrée dans une anthropologie communautaire : le « cœur » relie l'individu à ses ancêtres, à sa lignée, à la communauté.
\end{exemplebox}

Cette convergence est frappante : \emph{partout où l'homme pense sa condition}, il découvre que la conscience est ce qui le rend \emph{humain} — capable de vérité, de liberté, de responsabilité, de dignité. Qu'elle s'appelle \emph{cogito}, \emph{raison pratique}, \emph{discernement} ou \emph{cœur}, elle est la \emph{lumière intérieure} sans laquelle l'homme n'est qu'une « chose » parmi les choses.

\begin{methode}[Dans une dissertation, comment traiter la « valeur » d'une notion ?]
La « valeur » d'une notion (conscience, liberté, travail, technique, etc.) ne se raconte pas : elle se \emph{démontre} en montrant \textbf{ce qu'elle rend possible}. Structure type :
\begin{enumerate}
\item \textbf{Définir} la notion (ici : conscience = connaissance réflexive).
\item \textbf{Identifier les puissances} qu'elle actualise : \emph{épistémique} (connaissance de soi, certitude), \emph{éthique} (liberté, responsabilité, dignité), \emph{juridique/politique} (droits, objection de conscience), \emph{anthropologique} (humanité, « cœur »).
\item \textbf{Illustrer} chaque puissance par un auteur (Socrate, Descartes, Kant), un texte (Constitution), une situation concrète (tribunal, hôpital, vie courante).
\item \textbf{Nuancer :} La conscience \emph{garantit-elle} le bien ? Non (on peut user de sa liberté pour le mal). Est-elle \emph{infaillible} ? Non (erreur, mauvaise foi, inconscient). La valeur n'est pas une \emph{vertu automatique}, mais une \emph{capacité} qu'il faut \emph{cultiver} (éducation, examen de conscience).
\item \textbf{Conclure :} La conscience est la \emph{condition de possibilité} de l'humanité de l'homme.
\end{enumerate}
\end{methode}

\begin{popfigure}
\begin{tikzpicture}[
  node distance=2.5cm and 3cm,
  center node/.style={circle, draw=popblue, thick, fill=popblue!10, minimum width=3.2cm, align=center, font=\bfseries},
  ray node/.style={rectangle, rounded corners, draw=popdark, thick, fill=popcream, minimum width=3.5cm, minimum height=1.6cm, align=center, font=\small},
  arrow/.style={-Stealth, thick, popblue}
]
\node[center node, align=center] (C){La\\ Conscience};
\node[ray node, align=center] (K1) [above left=of C]{Connaissance\\ de soi\\(Socrate, Descartes)};
\node[ray node, align=center] (K2) [above right=of C]{Liberté\\ \& Délibération\\(Choix, autonomie)};
\node[ray node, align=center] (K3) [below right=of C]{Responsabilité\\ morale \& juridique\\(Imputabilité, droit)};
\node[ray node, align=center] (K4) [below left=of C]{Dignité\\ humaine\\(Fin en soi, Kant)};

\draw[arrow] (C) -- (K1);
\draw[arrow] (C) -- (K2);
\draw[arrow] (C) -- (K3);
\draw[arrow] (C) -- (K4);

% Petit décor : ancrage constitutionnel et tradition
\node[draw=popgreen, dashed, thick, fill=popgreen!5, rounded corners, fit=(K1)(K2)(K3)(K4), inner sep=0.8cm, label={[font=\tiny\bfseries, popgreen]above:Valeurs rendues possibles par la conscience}] {};
\end{tikzpicture}
\legende{Schéma : la conscience comme source centrale de quatre valeurs fondamentales — connaissance de soi, liberté, responsabilité, dignité. Chaque flèche indique une « puissance » que la conscience actualise.}
\end{popfigure}

\courssub{Synthèse intermédiaire : la conscience, fondement de l'humain}

Avant d'aborder l'inconscient (section 5), retenons l'essentiel : la conscience n'est pas un « luxe » psychologique ni un épiphénomène cérébral. Elle est la \textbf{condition de possibilité} de tout ce qui fait de l'homme un \emph{sujet} :
\begin{itemize}
\item \textbf{Sujet de savoir :} par le \emph{Cogito} (Descartes) et l'examen de soi (Socrate).
\item \textbf{Sujet de droit :} par la responsabilité imputable et l'objection de conscience (Constitution camerounaise).
\item \textbf{Sujet moral :} par la liberté de délibérer et la dignité de fin en soi (Kant, tradition du « cœur »).
\end{itemize}

Mais cette lumière a des \emph{limites} (section 4) : elle ne voit pas tout, elle peut se tromper, elle ignore ce qui la fonde. C'est là qu'intervient la découverte de l'\textbf{inconscient} — non pas comme ennemi, mais comme \emph{partie cachée} de nous-mêmes que la conscience doit apprendre à connaître pour être pleinement lucide et libre.

\coursec{Limites de la conscience}

Après avoir examiné la \emph{valeur} de la conscience — sa capacité à nous donner prise sur le réel, à fonder la responsabilité et à structurer l'unité du moi —, nous devons maintenant faire droit à ses \emph{failles}. La conscience n'est pas un œil parfait qui tout voit et tout sait ; elle est une faculté finie, située, vulnérable. Reconnaître ses limites n'est pas la nier, c'est la comprendre pour ce qu'elle est : une lumière partielle, non une transparence totale.

\courssub{Les erreurs des sens et les illusions perceptives : la conscience peut se tromper sur le monde}

La première limite de la conscience est \emph{perceptive}. Nous croyons voir le monde « tel qu'il est » ; en réalité, notre conscience construit une représentation à partir de données sensorielles fragmentaires, souvent ambiguës, que le cerveau interprète selon des inférences inconscientes. Lorsque les conditions de l'observation dévient de la norme, l'interprétation déraille : c'est l'illusion.

\begin{experience}[L'illusion de Müller-Lyer : quand la conscience « voit » faux]
\textbf{Matériel} : figure de deux segments de même longueur, l'un terminé par des pointes rentrantes ($>$--$<$), l'autre par des pointes sortantes ($<$--$>$).
\textbf{Protocole} : observer la figure ; mesurer les segments à la règle ; constater l'écart entre le jugement visuel immédiat et la mesure objective.
\textbf{Observation} : le segment aux pointes sortantes \emph{paraît} plus long, bien qu'il soit rigoureusement égal à l'autre.
\textbf{Interprétation} : la conscience perceptive applique des heuristiques de profondeur (coins de pièce, coins de bâtiment) à une figure plane. Elle « corrige » la taille rétinienne en fonction d'une profondeur \emph{supposée}. L'erreur n'est pas dans l'œil, mais dans l'interprétation \emph{inconsciente} qui précède la conscience visuelle.
\end{experience}

\begin{popfigure}
\begin{tikzpicture}[scale=1.2]
% Segment supérieur : pointes rentrantes (paraît plus court)
\draw[thick, popblue] (0,0) -- (5,0);
\draw[thick, popblue] (0,0) -- (-0.4,0.3) (0,0) -- (-0.4,-0.3);
\draw[thick, popblue] (5,0) -- (5.4,0.3) (5,0) -- (5.4,-0.3);
% Flèches de mesure
\draw[<->, popdark, thick] (0,-0.7) -- (5,-0.7);
\node[popdark] at (2.5,-1.0) {Longueur réelle : L};
% Segment inférieur : pointes sortantes (paraît plus long)
\draw[thick, poporange] (0,-2) -- (5,-2);
\draw[thick, poporange] (0,-2) -- (0.4,-1.7) (0,-2) -- (0.4,-2.3);
\draw[thick, poporange] (5,-2) -- (4.6,-1.7) (5,-2) -- (4.6,-2.3);
\draw[<->, popdark, thick] (0,-2.7) -- (5,-2.7);
\node[popdark] at (2.5,-3.0) {Longueur réelle : L (identique)};
% Légendes
\node[align=center, popblue] at (2.5,0.7) {Pointes rentrantes $\rightarrow$ \emph{paraît plus court}};
\node[align=center, poporange] at (2.5,-1.3) {Pointes sortantes $\rightarrow$ \emph{paraît plus long}};
\end{tikzpicture}
\legende{Illusion de Müller-Lyer : deux segments égaux perçus comme inégaux. La conscience perceptive est trompée par un traitement inconscient des indices de profondeur.}
\end{popfigure}

\begin{aretenir}[La perception n'est pas une photographie]
La conscience perceptive est une \emph{construction} (inférence inconsciente, hypothèses bayésiennes). Elle est adaptée à la survie (taille constante, forme constante), non à la vérité métrique brute. \textbf{La conscience du monde est médiatisée par des traitements pré-conscients qui peuvent l'induire en erreur.}
\end{aretenir}

\courssub{L'illusion de la transparence à soi : on croit se connaître alors qu'on s'ignore}

Au-delà du monde extérieur, la conscience se trompe sur \emph{elle-même}. L'intuition cartésienne (« je pense, donc je suis ») suggère une transparence immédiate : je sais ce que je pense, je sais ce que je veux. L'histoire de la philosophie (et la psychologie moderne) montre le contraire.

\begin{savaistu}[Nietzsche et la conscience comme surface]
Dans \emph{Le Gai Savoir} (\S 354) et \emph{Par-delà bien et mal}, \textbf{Friedrich Nietzsche} écrit : \emph{« La conscience est la dernière et la plus tardive évolution de l'organique [...] Elle ne concerne que les rapports de l'homme à l'homme [...] La conscience n'est au fond qu'un filet de communication entre les hommes. »} Il compare la conscience à la \emph{surface d'une mer} dont les profondeurs (instincts, pulsions, physiologie, histoire évolutive) nous échappent. Nous ne sommes pas les « capitaines » de notre âme, mais les observateurs tardifs de processus qui nous précèdent et nous dépassent.
\end{savaistu}

\begin{definition}[Illusion de la transparence à soi]
C'est la croyance naïve selon laquelle l'esprit aurait un accès direct, complet et infaillible à ses propres états (motifs, désirs, raisons d'agir). Les sciences cognitives et la psychanalyse montrent que nos \emph{vraies} motivations sont souvent opaques ; nous les \emph{reconstruisons} a posteriori sous forme de récits cohérents (confabulation, rationalisation). \cle{transparence à soi}
\end{definition}

\begin{exemplebox}[L'expérience du « choix libre » (Libet, Soon et al.)]
Des expériences en neuroimagerie (IRMf) montrent que l'activité cérébrale préfigurant une décision simple (lever le bras gauche ou droit) est détectable \emph{jusqu'à 7 à 10 secondes} avant que le sujet n'en ait la conscience subjective. La conscience du « je décide » survient \emph{après} que le cerveau a déjà enclenché le processus. La décision consciente serait un \emph{récit rétrospectif} justifiant un processus déjà lancé.
\end{exemplebox}

\courssub{Les déterminismes biologiques : hérédité, fatigue, maladie}

La conscience repose sur un substrat biologique (cerveau, système nerveux, endocrine). Ce substrat impose ses lois : la conscience ne flotte pas au-dessus de la matière, elle \emph{en} émerge et en subit les contraintes.

\begin{itemize}
    \item \textbf{Hérédité / tempérament} : prédispositions anxieuses, impulsivité, seuil de réactivité émotionnelle — la conscience « reçoit » un caractère qu'elle n'a pas choisi.
    \item \textbf{Fatigue, sommeil, rythmes circadiens} : la vigilance, la clarté du jugement, la maîtrise de soi fluctuent selon l'heure, le manque de sommeil, l'accumulation d'adénosine.
    \item \textbf{Pathologies neurologiques / psychiatriques} : lésion du lobe frontal (perte de l'initiative, désinhibition), dépression (ralentissement psychomoteur, filtre négatif sur la réalité), manie (euphorie, grandiosité), tumeurs (changements de personnalité). Phineas Gage (1848) : un barreau de fer traverse son lobe frontal ; il survit, mais « Gage n'était plus Gage » — sa conscience morale et sociale est transformée par une lésion physique.
\end{itemize}

\begin{attention}[Réductionnisme vs. prise en compte du substrat]
\emph{Erreur fréquente} : conclure que « tout est biologique, la conscience n'est qu'un épiphénomène ». \textbf{Nuance} : le biologique est une \emph{condition nécessaire} (sans cerveau, pas de conscience), non une \emph{explication suffisante} du sens des contenus conscients. La philosophie de l'esprit distingue \emph{corrélation} et \emph{identité}. Au Probatoire / Baccalauréat technique, on attend une analyse dialectique : la conscience est \emph{incarnée}, donc limitée par son corps, mais elle n'est pas \emph{réductible} à son corps.
\end{attention}

\courssub{Les déterminismes sociaux : l'éducation, la coutume, l'opinion}

Si le corps nous limite par le dedans, la société nous limite par le dehors — qui devient dedans par \emph{intériorisation}. \textbf{Émile Durkheim} (sociologie) et \textbf{Karl Marx} (idéologie) montrent que nos jugements « personnels » sont souvent des échos sociaux.

\begin{itemize}
    \item \textbf{Socialisation primaire (famille, école)} : acquisition de normes, de valeurs, de catégories de perception (espace, temps, causalité, sacré/profane) sans examen critique. \emph{Ex.} : la notion de « sale » ou « pur » varie radicalement selon les cultures ; la conscience morale l'applique comme une évidence.
    \item \textbf{Opinion publique, conformisme (Asch, Milgram)} : la conscience cède sous la pression du groupe. Dans l'expérience d'Asch, 75\% des sujets donnent au moins une fois une réponse fausse (évidente) pour s'aligner sur la majorité.
    \item \textbf{Idéologie (Marx)} : les idées dominantes sont les idées de la classe dominante. La conscience individuelle se prend pour l'auteur de pensées qui sont des \emph{représentations déformées} des rapports sociaux réels (ex. : « le pauvre est paresseux » masque l'exploitation structurelle).
    \item \textbf{Habitus (Bourdieu)} : systèmes de dispositions durables, inculqués par la classe sociale, qui génèrent des pratiques et des perceptions « naturelles » (goûts, posture, rapport au langage) sans passage par la conscience délibérante. \cle{habitus}
\end{itemize}

\begin{savaistu}[Conscience et déterminismes sociaux au Cameroun]
Au Cameroun, l'\emph{habitus} de classe et les représentations ethno-régionales structurent souvent la conscience sans qu'on s'en aperçoive. Exemple : un élève brillant, fils de paysan dans le Nord, peut s'auto-exclure des filières d'excellence (médecine, grandes écoles) non par manque de capacités, mais parce qu'il a intériorisé l'idée que « ces places sont pour les enfants de notables ou de fonctionnaires urbains ». Inversement, le conformisme linguistique (français/anglais vs langues nationales) ou le poids de la « grande famille » sur les choix d'orientation illustrent comment la conscience individuelle est habitée par des normes collectives. \textbf{La lucidité philosophique commence par repérer ces déterminismes pour s'en affranchir partiellement.}
\end{savaistu}

\begin{exoresolu}[Analyser une situation de déterminisme social]
\textbf{Énoncé} : « Paul, fils d'ouvrier, refuse d'intégrer une classe préparatoire prestigieuse où il a été admis, disant : "Ce n'est pas pour les gens comme moi." Identifiez le déterminisme social à l'œuvre. »
\textbf{Solution} :
\begin{enumerate}
    \item \textbf{Fait observable} : auto-exclusion malgré la réussite objective (admission).
    \item \textbf{Mécanisme} : \emph{habitus} de classe (Bourdieu) / \emph{intériorisation} de la division sociale (Durkheim). Paul a incorporé la croyance que « les grandes écoles sont pour les autres ».
    \item \textbf{Rapport à la conscience} : Paul \emph{croit} choisir librement (authenticité, fidélité à soi), mais sa conscience est le théâtre d'une contrainte sociale qu'il ne perçoit pas comme telle. C'est une \emph{liberté illusoire} (Spinoza : « les hommes se croient libres parce qu'ils sont conscients de leurs actions et ignorants des causes qui les déterminent »).
\end{enumerate}
\end{exoresolu}

\courssub{Le mensonge à soi-même et la mauvaise foi (Sartre) : la conscience peut se tromper elle-même}

C'est la limite la plus troublante : la conscience n'est pas seulement \emph{passive} (subissant illusions, déterminismes), elle peut être \emph{active} dans sa propre cécité. \textbf{Jean-Paul Sartre} (\emph{L'Être et le Néant}, 1943) analyse la \textbf{mauvaise foi} comme une structure ontologique de la conscience humaine.

\begin{definition}[Mauvaise foi (Sartre)]
Attitude par laquelle la conscience \emph{se ment à elle-même} pour fuir l'angoisse de sa liberté radicale. Elle consiste à se faire \emph{être-en-soi} (chose déterminée, nature fixe) alors qu'on est \emph{être-pour-soi} (néant, liberté, projet). La mauvaise foi n'est pas un mensonge ordinaire (le menteur sait la vérité et la cache à autrui) : le menteur à soi-même \emph{sait et ne sait pas} simultanément ; il est à la fois le trompeur et le trompé. \cle{mauvaise foi}
\end{definition}

\begin{exemplebox}[Le garçon de café (Sartre)]
Son mouvement est \emph{trop} précis, \emph{trop} empressé ; il \emph{joue} à être un garçon de café comme on joue un rôle au théâtre. Il fuit sa liberté de \emph{ne pas être} ce garçon de café (de pouvoir démissionner, de rêver d'autre chose) en se figeant dans une \emph{essence} (« je \emph{suis} un garçon de café »). Sa conscience se trompe elle-même en prenant un \emph{fait} (son emploi actuel) pour une \emph{nature} (son être nécessaire).
\end{exemplebox}

\begin{exemplebox}[La femme au rendez-vous (Sartre)]
Elle laisse sa main dans celle de l'homme « sans s'en apercevoir », ni consentir ni refuser. Elle reporte sa conscience sur la conversation intellectuelle pour \emph{ne pas décider} du sens charnel du geste. Elle se ment : elle \emph{sait} (son corps sait) mais elle \emph{refuse de savoir} « consciemment ». C'est la structure même de la mauvaise foi : \emph{disjonction} entre le savoir corporel/pré-réfléchi et le savoir thétique/réfléchi.
\end{exemplebox}

\begin{attention}[Mauvaise foi $\neq$ simple erreur ou ignorance]
\emph{Erreur fréquente} : confondre mauvaise foi et simple méconnaissance. La mauvaise foi \emph{exige} une lucidité minimale (on ne peut se mentir que si l'on sait, confusément, la vérité qu'on fuit). C'est une \emph{complicité} de la conscience avec elle-même. Au Probatoire / Baccalauréat technique, citez Sartre et distinguez : \emph{ignorance} (je ne sais pas) / \emph{mensonge à autrui} (je sais, je cache) / \emph{mauvaise foi} (je sais et je ne sais pas, je me cache à moi-même).
\end{attention}

\courssub{Les actes inexpliqués : oublis, lapsus, actes manqués, rêves}

La vie quotidienne fourmille d'événements que la conscience \emph{ne commande pas} et \emph{ne comprend pas} sur le moment. Ils sont les « symptômes » cliniques d'une autre instance psychique.

\begin{itemize}
    \item \textbf{Oubli de noms, de rendez-vous} : pas simple défaillance mnésique ; souvent sélectif (nom d'une personne qu'on évite, rendez-vous angoissant).
    \item \textbf{Lapsus linguae (parole)} : dire « beau-père » au lieu de « patron » ; « je suis \emph{malheureux} » au lieu de « \emph{heureux} ». Le mot « vrai » surgit à la place du mot « convenu ».
    \item \textbf{Actes manqués (Freud)} : perdre un objet qu'on devait remettre, faire une faute de frappe révélatrice, rater un examen qu'on voulait réussir. L'acte « rate » son but conscient pour accomplir un but \emph{inconscient}.
    \item \textbf{Rêves} : « La voie royale vers l'inconscient » (Freud). La conscience endormie produit des scénarios bizarres, symboliques, chargés d'affects, qui échappent au contrôle volontaire et à la logique diurne.
\end{itemize}

\begin{exemplebox}[Le lapsus révélateur]
Un employé, convoqué par son directeur autoritaire, dit en entrant : \emph{« Bonjour, beau-père »} au lieu de \emph{« Bonjour, patron »}.
\textbf{Analyse freudienne} : le mot « beau-père » condense \emph{autorité}, \emph{figure paternelle}, \emph{hostilité refoulée} (le beau-père stéréotypé est un rival, un intrus). La conscience voulait dire « patron » (politesse sociale) ; l'inconscient a fait « glisser » le signifiant vers sa vérité affective. La conscience est \emph{dépossédée} de sa parole. \cle{lapsus} \cle{refoulement}
\end{exemplebox}

\begin{popfigure}
\begin{tikzpicture}[scale=0.9]
% Iceberg
\draw[thick, popblue!80!popdark, fill=popblue!10] 
    (0,0) ellipse ({4} and {1.5}) % waterline
    .. controls (-4,0) and (-5,-1) .. (-5,-4)
    .. controls (-5,-4) and (5,-4) .. (5,-4)
    .. controls (5,-4) and (4,-1) .. (4,0);
% Water line
\draw[popblue, thick] (-4.5,0) -- (4.5,0);
% Conscious part (above water)
\node[align=center, popdark, font=\bfseries] at (0,0.5) {CONSCIENCE\\ (Partie émergée)};
\node[align=center, popmuted, font=\small] at (0,-0.3) {Pensées claires, volontés,\\ perceptions, moi conscient};
% Unconscious part (below water)
\node[align=center, poppurple, font=\bfseries] at (0,-2) {INCONSCIENT\\ (Partie immergée)};
\node[align=center, popmuted, font=\small] at (0,-2.8) {Pulsions, refoulements,\\ souvenirs traumatiques,\\ désirs inavoués,\\ automatismes};
% Arrows
\draw[->, poporange, thick] (0,-1.2) -- (0,-1.7) node[midway, right, poporange] {Refoulement};
\draw[<-, popgreen, thick] (0,-3.2) -- (0,-3.7) node[midway, right, popgreen, align=center]{Retour du refoulé\\(rêves, lapsus, symptômes)};
% Boat (ego)
\draw[popdark, fill=popcream] (-0.3,0.2) -- (0.3,0.2) -- (0.2,0.5) -- (-0.2,0.5) -- cycle;
\node[popdark, font=\tiny] at (0,0.35) {Moi};
\end{tikzpicture}
\legende{Schéma de l'iceberg (métaphore freudienne/topique) : la conscience n'est que la partie émergée d'un appareil psychique dont l'essentiel (inconscient) échappe au regard direct.}
\end{popfigure}

\courssub{Synthèse : ces limites ouvrent la voie à l'hypothèse de l'inconscient}

\begin{aretenir}[Bilan des limites de la conscience]
\begin{enumerate}
    \item \textbf{Perceptive} : illusions (Müller-Lyer) $\rightarrow$ la conscience construit, elle ne miroite pas.
    \item \textbf{Introspective} : illusion de transparence (Nietzsche) $\rightarrow$ le moi est opaque à lui-même.
    \item \textbf{Biologique} : fatigue, lésions, tempérament $\rightarrow$ la conscience est \emph{incarnée}, dépendante.
    \item \textbf{Sociale} : habitus, idéologie, conformisme $\rightarrow$ la conscience est \emph{intériorisation} de l'extérieur.
    \item \textbf{Éthique/Existentielle} : mauvaise foi (Sartre) $\rightarrow$ la conscience \emph{se} trompe pour fuir sa liberté.
    \item \textbf{Clinique/Quotidienne} : lapsus, actes manqués, rêves $\rightarrow$ la conscience ne \emph{maîtrise} pas la vie psychique.
\end{enumerate}
\end{aretenir}

\begin{methode}[Argumenter les limites d'une notion (Probatoire / Baccalauréat technique Philo)]
\begin{enumerate}
    \item \textbf{Ne pas confondre « limite » et « néant »}. Dire que la conscience a des limites, ce n'est pas dire qu'elle est inutile ou illusoire. C'est dire qu'elle est \emph{finie}, \emph{située}, \emph{partielle}.
    \item \textbf{Apporter des faits observables}, pas des opinions. \emph{Ex.} : « L'illusion de Müller-Lyer est reproductible en laboratoire » ; « L'expérience d'Asch montre un taux de conformité de 75\% » ; « Le lapsus "beau-père/patron" est un fait de langage documenté ».
    \item \textbf{Citer les auteurs pertinents} comme \emph{témoins} ou \emph{théoriciens} de ces faits : Nietzsche (transparence), Marx/Durkheim/Bourdieu (social), Freud (lapsus/rêves), Sartre (mauvaise foi), Spinoza (ignorance des causes).
    \item \textbf{Articuler les niveaux} : biologique $\neq$ social $\neq$ existentiel. Ne les réduisez pas les uns aux autres.
    \item \textbf{Conclure sur le sens philosophique} : ces limites ne détruisent pas la conscience ; elles la \emph{relativisent} et appellent une \emph{vigilance} (éthique, épistémologique, psychanalytique). Elles justifient l'hypothèse d'un \emph{inconscient} comme instance nécessaire pour rendre compte de ce qui, en nous, pense et agit sans passer par la conscience.
\end{enumerate}
\end{methode}

\begin{exercice}[Entraînement dissertation]
\textbf{Sujet} : « La conscience est-elle maîtresse en sa propre maison ? » (d'après Freud).
\textbf{Consignes} :
\begin{enumerate}
    \item Analysez les termes : « maîtresse » (pouvoir, contrôle, transparence), « en sa propre maison » (intériorité, vie psychique).
    \item Faites un plan dialectique : Thèse (oui, conscience = unité, responsabilité, lucidité) / Antithèse (non, limites ci-dessus : inconscient, déterminismes, mauvaise foi) / Synthèse (la conscience n'est pas \emph{tout-puissante} mais \emph{responsable} : elle doit \emph{conquérir} sa maîtrise par le travail sur soi, l'analyse, la vigilance critique).
    \item Rédigez l'introduction (accroche, définition, thèse, plan) et le premier paragraphe de la thèse.
\end{enumerate}
\end{exercice}

\begin{corrige}[Exercice n°1 - Éléments de correction]
\textbf{Introduction} : Freud, \emph{Introduction à la psychanalyse} (1917) : « Le moi n'est pas maître en sa propre maison. » Accroche : la métaphore de la maison suggère l'intimité, la propriété, le contrôle. Problème : la conscience a-t-elle un pouvoir souverain sur la vie psychique, ou bien d'autres instances (inconscient, corps, société) y font la loi ? Thèse : la conscience n'est pas maîtresse \emph{par nature} (donnée immédiate), mais elle peut le devenir \emph{par conquête} (travail de vérité).
\textbf{Premier paragraphe (Thèse)} : La conscience \emph{est} le lieu de la maîtrise \emph{possible}. 1) Unité du moi (Descartes : \emph{cogito} fondateur). 2) Responsabilité morale/juridique : la loi suppose un sujet conscient de ses actes (libre arbitre). 3) Projet sartrien : la conscience est liberté ; elle \emph{peut} s'assumer comme telle (authenticité). Mais cette maîtrise est une \emph{tâche}, non un don.
\end{corrige}

\medskip
\noindent\textbf{Transition vers la section 5 :} Si la conscience n'est pas maîtresse en sa maison, \emph{qui} y habite ? Les limites que nous venons de cartographier (illusions, déterminismes, refoulements, lapsus) ne sont pas de simples « trous » dans la conscience : elles sont les \emph{signes}, les \emph{effets de surface} d'une instance profonde, structurée, qui a sa logique propre : l'\textbf{inconscient}. C'est à sa découverte, à ses types et à la théorie freudienne que nous allons maintenant nous attacher.

\coursec{La découverte de l'inconscient et ses types}

\courssub{Transition : des limites de la conscience à la nécessité de l'inconscient}

Nous avons vu que la conscience, malgré sa valeur de lumière et de maîtrise, possède des \cle{limites structurelles} : elle est étroite, intermittente (sommeil, distraction), et surtout impuissante face à certains de nos propres actes (lapsus, oublis, symptômes névrotiques, rêves). Si je ne suis pas toujours maître en ma propre maison, comme le disait déjà Leibniz, c'est qu'il existe une autre scène psychique, hors de mon regard intérieur. C'est la découverte de cette \emph{autre scène} — l'inconscient — qui va révolutionner la compréhension de l'homme à la fin du XIXe siècle. Passons maintenant de la \emph{phénoménologie} de la conscience à la \emph{métapsychologie} de l'inconscient.

\textbf{Contexte historique : l'hystérie, l'hypnose et la naissance d'une science nouvelle}\par

\paragraph{La leçon de la Salpêtrière (1885).} Tout commence non dans un cabinet de philosophe, mais à l'hôpital de la \textbf{Salpêtrière} à Paris, sous la direction du neurologue \textbf{Jean-Martin Charcot}. Charcot étudie l'\cle{hystérie} : des paralysies, des cécités, des contractures sans aucune lésion organique du système nerveux. Il découvre que sous \cle{hypnose}, ces symptômes disparaissent ou peuvent être suggérés. La suggestion hypnotique prouve que l'esprit peut commander au corps sans passer par la conscience volontaire.

\begin{experience}[L'expérience fondatrice : l'hystérie et la suggestion]
\textbf{Observation (Charcot, 1885).} Une patiente présente une paralysie du bras gauche. Aucun dommage neurologique n'est trouvé. Sous hypnose, le médecin suggère : « Votre bras est guéri. » La patiente bouge le bras. Au réveil, elle ignore comment elle a fait, mais la paralysie a disparu.
\textbf{Interprétation.} Le symptôme n'est pas une lésion matérielle, mais une \emph{conversion} d'un conflit psychique en langage corporel. La guérison par suggestion prouve l'existence de processus psychiques \emph{efficaces} mais \emph{inconscients}.
\textbf{Conséquence pour Freud.} Le jeune \textbf{Sigmund Freud} (1856–1939), interne à la Salpêtrière en 1885, comprend que l'hypnose n'est qu'un outil provisoire. Il faut une méthode pour accéder à l'inconscient \emph{sans} hypnose : ce sera la \cle{méthode cathartique} (avec Breuer) puis la \cle{psychanalyse} (association libre, interprétation des rêves).
\end{experience}

\begin{popfigure}
\begin{tikzpicture}[scale=0.9, every node/.style={font=\small}]
% Frise chronologique
\draw[popblue, thick] (0,0) -- (16,0);
% Dates et événements
\foreach \x/\date/\event/\detail in {
    1/1885/Charcot \& l'hypnose/Freud interne à la Salpêtrière,
    4/1895/Études sur l'hystérie/Breuer \& Freud : méthode cathartique,
    7/1900/L'Interprétation des rêves/Fondation de la psychanalyse,
    10/1915/Métapsychologie/Topique : Inconscient - Préconscient - Conscient,
    13/1923/Le Moi et le Ça/Deuxième topique : Ça - Moi - Surmoi} {
    \draw[popblue, thick] (\x,-0.15) -- (\x,0.15);
    \node[above, align=center, popdark] at (\x,0.5) {\textbf{\date}\\ \event};
    \node[above, align=center, popmuted] at (\x,1.5) {\detail};
}
% Flèches d'évolution
\draw[->, poporange, thick] (1.5,-0.5) -- (3.5,-0.5) node[midway, below] {Rejet de l'hypnose};
\draw[->, poporange, thick] (4.5,-0.5) -- (6.5,-0.5) node[midway, below] {Association libre};
\draw[->, poporange, thick] (7.5,-0.5) -- (9.5,-0.5) node[midway, below] {Théorie du rêve};
\draw[->, poporange, thick] (10.5,-0.5) -- (12.5,-0.5) node[midway, below] {Structuration théorique};
\end{tikzpicture}
\legende{Frise chronologique : de la neurologie charcotienne à la métapsychologie freudienne (1885–1923).}
\end{popfigure}

\textbf{Les preuves de l'inconscient selon Freud : la « voie royale »}\par

Freud ne postule pas l'inconscient par spéculation métaphysique ; il l'infère comme une \textbf{hypothèse explicative nécessaire} pour rendre compte de phénomènes cliniques observables. Il en dresse l'inventaire dans \emph{La Psychopathologie de la vie quotidienne} (1901) et \emph{L'Interprétation des rêves} (1900).

\begin{definition}[Preuves de l'inconscient (Freud)]
Ce sont des formations de l'inconscient (\emph{formations de compromis}) où un contenu refoulé (désir, souvenir, pulsion) parvient à s'exprimer de manière déguisée, contournant la censure du conscient. Les quatre preuves majeures sont :
\begin{enumerate}
    \item \textbf{Les rêves} : « voie royale vers l'inconscient ». Le rêve accomplit un désir refoulé sous forme hallucinatoire déguisée (travail du rêve : condensation, déplacement, figuration).
    \item \textbf{Les lapsus (lapsus linguae, calami, lectionis)} : erreurs de parole, d'écriture, de lecture qui trahissent une intention cachée contraire à l'intention consciente.
    \item \textbf{Les actes manqués} : oublis de noms, promesses non tenues, objets perdus, méprises d'identité — actes en apparence accidentels mais motivés par un mobile inconscient.
    \item \textbf{Les symptômes névrotiques} : phobies, obsessions, compulsions, conversions hystériques. Ils sont des \emph{formations de compromis} entre le désir inconscient et la défense du Moi.
\end{enumerate}
\end{definition}

\begin{exemplebox}[Le lapsus du « père qui gronde »]
\textbf{Situation.} Un élève, convoqué par le proviseur pour une bêtise mineure, entre dans le bureau et dit : « Bonjour Monsieur, je viens pour que vous me \textbf{grondiez} » (au lieu de « receviez » ou « entendiez »).
\textbf{Analyse freudienne.} Ce n'est pas une simple fatigue. Le lapsus révèle un \textbf{désir inconscient} : l'élève \emph{veut} être grondé, peut-être pour expier une faute plus grave non découverte, ou par besoin de reconnaissance négative (mieux vaut être grondé qu'ignoré). Le mot « grondiez » a glissé à la place du mot socialement attendu : c'est l'inconscient qui parle.
\end{exemplebox}

\begin{exoresolu}[Analyse d'un acte manqué : l'oubli motivé]
\textbf{Énoncé.} Un étudiant en Première Technique prépare soigneusement son exposé sur « La conscience ». Le jour J, il arrive en classe, ouvre son sac… et réalise qu'il a oublié sa clé USB \emph{uniquement} pour cet exposé (il a bien ses autres cours). Il n'a pas dormi la veille, angoissé. Expliquez cet oubli à la lumière de la psychanalyse.
\tcblower
\textbf{Solution.}
\begin{enumerate}
    \item \textbf{Constat :} L'oubli est sélectif (seul l'exposé phobogène est concerné) et coïncide avec un état d'angoisse anticipatoire.
    \item \textbf{Hypothèse freudienne :} C'est un \textbf{oubli motivé} (\emph{Vergessenwollen}). L'inconscient refuse l'épreuve (peur de l'échec, du jugement, castration symbolique). L'acte manqué (oublier le support) réalise le désir inconscient : \emph{ne pas faire l'exposé}.
    \item \textbf{Mécanisme :} Le désir inconscient (éviter l'angoisse) l'emporte sur l'intention consciente (réussir l'exposé). L'oubli n'est pas une défaillance de la mémoire, mais une \emph{réussite} de la défense.
    \item \textbf{Conclusion :} L'inconscient n'est pas une absence de conscience, mais une \emph{autre conscience} avec ses propres finalités, ses lois (plaisir, refoulement) et son efficacité réelle.
\end{enumerate}
\end{exoresolu}

\textbf{Structure de l'appareil psychique : la première topique (1900–1920)}\par

Avant 1923, Freud organise la vie psychique en trois \textbf{systèmes} (topique) distingués par leur rapport à la conscience et au refoulement. C'est le modèle de l'\textbf{iceberg}.

\begin{popfigure}
\begin{tikzpicture}[scale=1.1, every node/.style={font=\small}]
% Eau
\fill[popblueL!30] (-3,-4) rectangle (3,0.5);
\draw[popblue, thick] (-3,0.5) -- (3,0.5) node[midway, above, sloped, popblue] {Ligne de flottaison = Seuil de conscience};
% Iceberg émergé : Conscient
\draw[popgreen, thick, fill=popgreen!20] (-1.5,0.5) -- (-1.5,2.5) -- (1.5,2.5) -- (1.5,0.5) -- cycle;
\node[popgreen, align=center] at (0,1.5) {\textbf{CONSCIENT}\\(Cs)\\Perceptions, pensées actuelles\\« Je sais que je sais »};
% Iceberg juste sous l'eau : Préconscient
\draw[poporange, thick, fill=poporange!20] (-2,0.5) -- (-2,-1) -- (2,-1) -- (2,0.5) -- cycle;
\node[poporange, align=center] at (0,-0.25) {\textbf{PRÉCONSCIENT}\\(Pcs)\\Souvenirs, savoirs accessibles\\« Je ne pense pas à mon nom,\\mais je peux le retrouver »};
% Iceberg profond : Inconscient
\draw[poppurple, thick, fill=poppurple!20] (-2.5,-1) -- (-2.5,-3.5) -- (2.5,-3.5) -- (2.5,-1) -- cycle;
\node[poppurple, align=center] at (0,-2.25) {\textbf{INCONSCIENT}\\(Ics)\\Pulsions refoulées, désirs interdits,\\traces mnésiques traumatiques\\« Je ne sais pas que je sais,\\et je ne \emph{veux} pas le savoir »};
% Flèches
\draw[<->, popline, thick] (2.2,1.5) -- (2.2,0.6) node[midway, right] {Attention};
\draw[<->, popline, thick] (2.7,-0.25) -- (2.7,-1.1) node[midway, right, align=center]{Effort de mémoire\\/ Association libre};
\draw[->, popred, thick, dashed] (0,-1) -- (0,-1.5) node[midway, right, popred] {REFOULEMENT};
\draw[->, popred, thick, dashed] (0,-2.25) .. controls (1.5,-2) and (1.5,1) .. (1,1.5) node[right, popred, align=center]{Retour du refoulé\\(symptômes, rêves, lapsus)};
\end{tikzpicture}
\legende{Schéma de la première topique freudienne (modèle de l'iceberg). L'inconscient est la partie immergée, la plus vaste et la plus déterminante.}
\end{popfigure}

\begin{definition}[Le Préconscient (Pcs)]
Système intermédiaire. Contenus psychiques \emph{momentanément absents} de la conscience (souvenirs, connaissances, automatismes) mais \textbf{accessibles} par un simple effort d'attention ou d'association (ex. : votre numéro de téléphone, la capitale du Cameroun). Il n'y a pas de refoulement, seulement un retrait de l'attention. C'est le \emph{réservoir} de la conscience.
\end{definition}

\begin{definition}[L'Inconscient (Ics) — Sens freudien strict]
Ensemble des contenus psychiques (\textbf{représentations-inconscientes} chargées d'affect / \textbf{pulsions}) qui ont été \textbf{refoulés} (\emph{Verdrängung}) hors de la conscience parce qu'ils sont incompatibles avec les exigences du Moi (morale, réalité, image de soi). L'inconscient \textbf{ignore le temps, la contradiction, la négation} : il ne connaît que le processus primaire (plaisir immédiat, condensation, déplacement). Il est \emph{dynamique} : il pousse sans cesse pour revenir (retour du refoulé).
\end{definition}

\begin{attention}[Piège majeur : l'inconscient n'est PAS une « seconde conscience »]
\begin{itemize}
    \item \textbf{Erreur fréquente :} « L'inconscient, c'est comme un petit homme caché dans ma tête qui tire les ficelles. »
    \item \textbf{Correction :} L'inconscient freudien n'a \textbf{pas de sujet} (pas de « je » inconscient). C'est un \emph{système de processus} (pulsions, mécanismes) sans conscience de soi. Il ne « pense » pas, il \emph{pousse}. Il obéit au \textbf{principe de plaisir} (décharge immédiate), non au principe de réalité.
    \item \textbf{Conséquence :} On ne « prend conscience » de l'inconscient que par ses \emph{effets} (rêves, symptômes, lapsus), jamais par introspection directe. D'où la nécessité de la \textbf{psychanalyse} (tiers, association libre, transfert).
\end{itemize}
\end{attention}

\textbf{Les types d'inconscient : Freud et Jung}\par

Le programme exige de distinguer l'inconscient \textbf{psychologique (individuel)} de l'inconscient \textbf{collectif}. Cette distinction est cruciale pour l'épreuve du Probatoire.

\begin{propriete}[L'inconscient psychologique (freudien) — Individuel et historique]
\begin{itemize}
    \item \textbf{Origine :} Histoire \emph{personnelle} du sujet (enfance, traumatismes, conflits œdipiens, refoulements spécifiques).
    \item \textbf{Contenu :} \textbf{Complexes} (noyaux affectifs refoulés), désirs incestueux, pulsions sexuelles et agressives (libido, Thanatos), souvenirs traumatiques.
    \item \textbf{Structure :} Organisé en \emph{complexes} autonomes (ex. : complexe d'Œdipe, complexe de castration, complexe d'infériorité).
    \item \textbf{Accessibilité :} Par la cure psychanalytique individuelle (association libre, analyse du transfert).
    \item \textbf{Exemple camerounais :} Un jeune ingénieur de Douala refuse toute promotion par peur inconsciente de surpasser son père (rivalité œdipienne non résolue). Son symptôme (auto-sabotage, maladies psychosomatiques avant les entretiens) ne se comprend qu'à la lumière de \emph{son} histoire familiale singulière.
\end{itemize}
\end{propriete}

\begin{propriete}[L'inconscient collectif (jungien) — Universel et héréditaire]
\textbf{Carl Gustav Jung} (1875–1961), disciple puis dissident de Freud, élargit la notion.
\begin{itemize}
    \item \textbf{Origine :} \emph{Phylogénétique} (héritage de l'espèce humaine), non ontogénétique (histoire personnelle). « L'inconscient collectif est la partie du psychisme qui n'est pas acquise individuellement, mais innée, héréditaire. »
    \item \textbf{Contenu :} Les \textbf{archétypes} (\emph{Ur-bilder} : images primordiales) — formes vides, matrices universelles qui structurent l'imaginaire humain : l'Ombre, l'Anima/Animus, le Sage, la Grande Mère, le Héros, le Soi, le Trickster.
    \item \textbf{Manifestation :} Mythes, contes, religions, rêves « grands » (non personnels), symboles culturels récurrents, comportements collectifs (mouvements de foule, renaissances mystiques).
    \item \textbf{Accessibilité :} Par l'analyse des symboles culturels, l'amplification mythologique, l'étude des rêves « typiques ».
    \item \textbf{Exemple camerounais :} Le motif de l'\textbf{Ancêtre protecteur} ou du \textbf{Génie de l'eau} (Mami Wata, Jengu) dans les cultures Sawa, Bassa, Bamiléké. Ce n'est pas un souvenir personnel refoulé, mais une \emph{structure symbolique universelle} (archétype de la Grande Mère / de l'Anima) qui organise le rapport au sacré, à la fécondité, au danger. Tout Camerounais, même non initié, reconnaît la puissance de ce symbole.
\end{itemize}
\end{propriete}

\begin{center}
\begin{tabularx}{\linewidth}{|l|X|X|}
\hline
\rowcolor{popblueL}
\textbf{Critère} & \textbf{Inconscient personnel (Freud)} & \textbf{Inconscient collectif (Jung)} \\
\hline
\textbf{Origine} & Biographique (ontogenèse) & Phylogénétique (espèce) \\
\hline
\textbf{Contenu} & Refoulés personnels, complexes, pulsions & Archétypes (formes universelles) \\
\hline
\textbf{Structure} & Dynamique (conflit désir/défense) & Structurale (matrices symboliques) \\
\hline
\textbf{Temporalité} & Ignore le temps personnel & Intemporel, éternel \\
\hline
\textbf{Méthode d'accès} & Psychanalyse (transfert, association libre) & Herméneutique symbolique, amplification \\
\hline
\textbf{Exemple} & Phobie des ascenseurs (trauma enfermement) & Rêve du Dragon / du Grand Initié (mythe du Héros) \\
\hline
\end{tabularx}
\end{center}

\textbf{L'inconscient comme hypothèse explicative : statut épistémologique}\par

\textbf{Problème.}\par L'inconscient, par définition, \emph{échappe à la conscience}. Comment peut-on le connaître scientifiquement ? N'est-ce pas une \emph{pétition de principe} (expliquer l'inconnu par l'inconnaissable) ?

\begin{methode}[Valider l'hypothèse de l'inconscient : démarche d'inférence à la meilleure explication]
\begin{enumerate}
    \item \textbf{Observation de faits résistants :} Symptômes sans cause organique, lapsus systématiques, rêves répétitifs, comportements irrationnels répétés. La conscience \emph{ne s'explique pas elle-même}.
    \item \textbf{Hypothèse théorique :} Postuler une instance psychique \emph{autre} (l'inconscient) dotée de lois propres (refoulement, processus primaire, symbolisation).
    \item \textbf{Pouvoir prédictif et thérapeutique :} L'hypothèse permet de \emph{prédire} la structure des rêves (travail du rêve), de \emph{retrouver} l'origine oubliée d'un symptôme, et surtout de le \emph{guérir} par la levée du refoulement (catharsis, insight).
    \item \textbf{Critère de scientificité (Popper / épistémologie) :} L'inconscient n'est pas directement observable (comme un électron), mais ses \emph{effets} le sont. C'est une \textbf{entité théorique} validée par sa fécondité clinique et explicative. \emph{Preuve par l'efficacité thérapeutique.}
\end{enumerate}
\end{methode}

\begin{aretenir}[Statut de l'inconscient en philosophie et en science]
\begin{itemize}
    \item L'inconscient n'est pas une \emph{métaphore} : c'est un \textbf{concept opératoire} qui rend compte de la \emph{causalité psychique} là où la conscience faillit.
    \item Il \textbf{ne se prouve pas par l'introspection} (puisqu'il est invisible pour le sujet), mais par l'\textbf{extrospection clinique} (analyse des productions : rêves, paroles, actes, symptômes).
    \item \textbf{Critique possible (Bac) :} L'inconscient freudien est-il \emph{falsifiable} ? (Popper : non, car il explique tout et son contraire). Réponse : La psychanalyse n'est pas une science de la nature (nomothétique) mais une \emph{herméneutique} (science de l'interprétation) validée par la \emph{transformation subjective} du patient.
\end{itemize}
\end{aretenir}

\begin{savaistu}[Freud et le Cameroun : une réception critique]
\begin{itemize}
    \item \textbf{1900 :} Parution de \emph{Die Traumdeutung} (L'Interprétation des rêves). Tiré à 600 exemplaires ; 8 ans pour écouler la première édition. Aujourd'hui : texte fondateur de la modernité psychique.
    \item \textbf{Années 1950–60 :} La psychanalyse arrive au Cameroun via les médecins coloniaux (hôpital de la Maternité d'Ebolowa, service de neuropsychiatrie de l'Hôpital Central de Yaoundé). Les premiers psychiatres camerounais formés en France (Dr. M'Bida, Dr. Owona) l'utilisent pour comprendre les « syndromes culture-bound » (ex. : \emph{bouffée délirante aiguë}, possession).
    \item \textbf{Aujourd'hui :} La psychanalyse freudienne et lacanienne irrigue la formation des psychologues cliniciens (ESSEC Yaoundé, Université de Yaoundé I). Elle est débattue : certains y voient un outil d'émancipation subjective ; d'autres (anthropologues, sociologues) critiquent son \emph{ethnocentrisme} (l'Œdipe comme universel ?). Le débat est vif : l'inconscient est-il \emph{universel} (Freud/Jung) ou \emph{culturellement construit} ?
\end{itemize}
\end{savaistu}

\textbf{Synthèse intermédiaire : conscience et inconscient, adversaires ou partenaires ?}\par

Nous avons découvert l'inconscient comme \textbf{réponse aux limites de la conscience}. Mais la relation ne s'arrête pas à l'opposition.
\begin{itemize}
    \item \textbf{Adversaires ?} Oui, au sens du \textbf{refoulement} : le Moi conscient se construit \emph{contre} l'inconscient (censure, défense). La névrose est le conflit.
    \item \textbf{Partenaires ?} Oui, au sens de la \textbf{sublimation} et de la \emph{psychanalyse} : l'inconscient fournit l'énergie (libido), la créativité, la profondeur ; la conscience apporte la forme, la réalité, la parole. « \emph{Wo Es war, soll Ich werden} » — « Là où c'était le Ça, le Moi doit advenir » (Freud, 1933).
\end{itemize}
La section suivante (6/7) approfondira la \textbf{théorie de la psychanalyse freudienne} (deuxième topique : Ça - Moi - Surmoi ; pulsions ; mécanismes de défense) et ses \textbf{critiques} (philosophiques, scientifiques, féministes, anthropologiques), pour préparer l'exercice de dissertation finale.

\coursec{La psychanalyse freudienne et sa critique}

Dans la section précédente, nous avons vu comment la découverte de l'inconscient — à travers l'hypnose, l'hystérie et les travaux de Charcot — a conduit à distinguer plusieurs types d'inconscient. Mais c'est un homme qui a donné à cette notion sa forme la plus célèbre et la plus discutée : Sigmund \cle{Freud} (1856-1939), médecin viennois, fondateur de la \cle{psychanalyse}. Nous allons maintenant présenter sa théorie de l'appareil psychique, sa méthode thérapeutique, puis les critiques majeures dont elle a fait l'objet, notamment celles de Sartre et de Popper. La question qui guidera toute cette section est la suivante : la psychanalyse est-elle une science de l'esprit humain, ou seulement une manière de l'interpréter ?

\courssub{La première topique : conscient, préconscient, inconscient}

Commençons par une intuition simple. En ce moment, vous êtes conscient de lire ce cours. Vous n'êtes pas en train de penser à votre adresse, mais si on vous la demande, elle vous revient immédiatement : elle était disponible, « en réserve ». En revanche, certains souvenirs ou désirs semblent totalement inaccessibles, comme enfouis, et pourtant ils agissent sur nous. Freud traduit cette intuition en un premier modèle de l'esprit, appelé \cle{première topique} (du grec \emph{topos}, le lieu) : il distingue trois « lieux » psychiques.

\begin{definition}[La première topique]
Selon Freud, la vie psychique se répartit en trois systèmes :
\begin{itemize}
\item le \cle{conscient} : ce dont le sujet a actuellement connaissance (perceptions, pensées présentes) ;
\item le \cle{préconscient} : ce qui n'est pas actuellement pensé mais peut devenir conscient sans difficulté (souvenirs, savoirs disponibles) ;
\item l'\cle{inconscient} : ensemble des représentations et des désirs \emph{refoulés}, maintenus hors de la conscience parce qu'ils sont inacceptables pour le sujet, mais qui continuent d'agir sur sa conduite.
\end{itemize}
\end{definition}

Le point essentiel est que l'inconscient freudien n'est pas un simple « non-conscient » passif : il est \emph{dynamique}. Les contenus refoulés cherchent sans cesse à revenir à la conscience, et ils y parviennent sous des formes déguisées : rêves, lapsus, actes manqués, symptômes. L'inconscient se manifeste donc indirectement, et c'est précisément ce qui permet de l'étudier. Freud résume cette découverte par une formule restée célèbre : le Moi « n'est pas maître dans sa propre maison ».

\courssub{La seconde topique : le Ça, le Moi, le Surmoi}

Vers 1920, Freud perfectionne son modèle et propose une \cle{seconde topique}, qui décrit l'esprit non plus comme trois lieux, mais comme trois \emph{instances} en conflit permanent. Ce second modèle ne remplace pas le premier : il le précise, en montrant quelles forces s'affrontent dans la vie psychique.

\begin{definition}[Les trois instances de la seconde topique]
\begin{itemize}
\item Le \cle{Ça} : réservoir des \emph{pulsions} (désirs, besoins, agressivité). Il obéit au \emph{principe de plaisir} : il veut la satisfaction immédiate, sans tenir compte de la réalité ni de la morale. Il est entièrement inconscient.
\item Le \cle{Moi} : instance de médiation. Il obéit au \emph{principe de réalité} : il cherche à satisfaire les désirs du Ça, mais en tenant compte du monde extérieur, des dangers et des contraintes sociales.
\item Le \cle{Surmoi} : instance psychique issue de l'\emph{intériorisation des interdits parentaux et sociaux}. Il fonctionne comme un juge intérieur : il approuve ou condamne, et il est la source du sentiment de \emph{culpabilité}.
\end{itemize}
\end{definition}

Prenons une image célèbre : Freud compare le Moi à un cavalier (le Ça étant le cheval). Le cheval fournit l'énergie, mais c'est le cavalier qui doit diriger — et il arrive que le cheval l'emporte. Le Moi est donc pris entre deux exigences opposées : les pulsions du Ça (« je veux, tout de suite ») et les interdits du Surmoi (« tu ne dois pas »). Un élève tenté de tricher à un examen illustre ce conflit : le Ça désire la réussite facile, le Surmoi condamne la tricherie, et le Moi doit négocier une issue réaliste — par exemple, travailler davantage.

\begin{popfigure}
\begin{center}
\begin{tikzpicture}[font=\small]
% Zones
\fill[poporangeL] (0,0) circle (2.6);
\fill[popblueL,opacity=0.85] (-1.5,1.1) circle (1.5);
\fill[poppurpleL,opacity=0.85] (1.5,1.1) circle (1.5);
\fill[popgreenL,opacity=0.9] (0,-0.6) circle (1.3);
% Labels
\node[popdark,font=\bfseries] at (-1.5,1.1) {Le Ça};
\node[popdark] at (-1.5,0.55) {\footnotesize pulsions};
\node[popdark] at (-1.5,0.2) {\footnotesize principe de plaisir};
\node[popdark,font=\bfseries] at (1.5,1.1) {Le Surmoi};
\node[popdark] at (1.5,0.55) {\footnotesize interdits};
\node[popdark] at (1.5,0.2) {\footnotesize culpabilité};
\node[popdark,font=\bfseries] at (0,-0.6) {Le Moi};
\node[popdark] at (0,-1.15) {\footnotesize principe de réalité};
% Flèches de conflit et de médiation
\draw[-{Stealth},very thick,popink] (-0.9,0.2) -- (-0.35,-0.15);
\draw[-{Stealth},very thick,popink] (0.9,0.2) -- (0.35,-0.15);
\draw[{Stealth}-{Stealth},very thick,popgold] (-0.4,1.6) -- (0.4,1.6);
\node[popgold,font=\footnotesize] at (0,1.95) {conflit};
\node[popink,font=\footnotesize] at (-1.35,-0.55) {désirs};
\node[popink,font=\footnotesize] at (1.35,-0.55) {interdits};
\end{tikzpicture}
\end{center}
\legende{La seconde topique freudienne : le Moi (vert) doit arbitrer entre les pulsions du Ça (bleu) et les interdits du Surmoi (violet), eux-mêmes en conflit direct (flèche dorée). Schéma simplifié.}
\end{popfigure}

\courssub{Le refoulement et le complexe d'Œdipe}

Comment des contenus psychiques deviennent-ils inconscients ? Par un mécanisme central de la théorie freudienne : le \cle{refoulement}.

\begin{definition}[Le refoulement]
Le \cle{refoulement} est le mécanisme de défense par lequel le Moi maintient dans l'inconscient les représentations et les désirs incompatibles avec la morale, avec l'image que le sujet a de lui-même ou avec les interdits sociaux. Le refoulé ne disparaît pas : il continue d'agir et cherche à revenir sous des formes déguisées.
\end{definition}

Le refoulement trouve son origine, selon Freud, dans les conflits de la petite enfance, et notamment dans le \cle{complexe d'Œdipe}. Freud emprunte ce nom au mythe grec d'Œdipe, qui tue son père et épouse sa mère sans le savoir. De manière sobre : entre trois et cinq ans environ, l'enfant éprouverait un attachement privilégié pour le parent de sexe opposé et un sentiment de rivalité envers le parent de même sexe. Ce conflit, impossible à assumer, est refoulé ; sa résolution (l'identification au parent de même sexe) structurerait durablement la personnalité et contribuerait à la formation du Surmoi. Qu'on adhère ou non à cette thèse, elle illustre l'idée fondamentale de Freud : \emph{l'enfant est le père de l'homme} — la vie psychique adulte porte la marque des conflits infantiles refoulés.

\courssub{La méthode psychanalytique : comment accéder à l'inconscient ?}

Si l'inconscient est par définition inaccessible à la conscience directe, comment le connaître ? Freud élabore une méthode originale, fondée sur l'idée que l'inconscient se trahit dans les productions apparemment insignifiantes de la vie psychique.

\begin{methode}[Les voies d'accès à l'inconscient selon Freud]
\begin{enumerate}
\item La \cle{libre association} : le patient, allongé, dit tout ce qui lui passe par l'esprit, sans trier ni censurer. Les enchaînements apparemment absurdes révèlent des liens inconscients.
\item L'\cle{interprétation des rêves} : Freud l'appelle la « voie royale vers l'inconscient ». Pendant le sommeil, la censure du Surmoi se relâche, et les désirs refoulés s'expriment sous forme déguisée.
\item L'analyse des \cle{actes manqués} : lapsus, oublis, erreurs. Ils ne sont pas de simples accidents : ils expriment un désir ou une hostilité inconscients.
\item Le \cle{transfert} : au cours de la cure, le patient reporte sur l'analyste des sentiments (amour, hostilité) issus de ses relations infantiles ; l'analyse de ce transfert permet de remonter aux conflits refoulés.
\end{enumerate}
\end{methode}

\begin{propriete}[Le rêve comme accomplissement déguisé d'un désir]
Pour Freud, le rêve est « l'accomplissement déguisé d'un désir refoulé ». Le contenu manifeste (ce qu'on se rappelle du rêve) est la traduction déformée d'un contenu latent (le désir inconscient), déformation imposée par la censure morale. (Propriété admise, non démontrée : elle relève d'une interprétation, non d'une preuve expérimentale.)
\end{propriete}

\begin{exemplebox}[L'acte manqué : l'oubli du cahier de mathématiques]
Un élève de Première technique « oublie » systématiquement son cahier de mathématiques les jours où il a cours dans cette matière, alors qu'il n'oublie jamais ses autres affaires. Pour Freud, cet oubli n'est pas un hasard : c'est un \emph{acte manqué} qui révèle une hostilité inconsciente envers la matière (peur de l'échec, conflit avec un professeur, sentiment d'incompétence refoulé). L'oubli « arrange » le sujet : il lui évite la situation redoutée, tout en lui épargnant d'avouer cette aversion.
\end{exemplebox}

\begin{exoresolu}[Analyser une situation avec les concepts freudiens]
Énoncé. Une jeune femme se prépare à un entretien d'embauche important à Douala. Le matin venu, elle se trompe de bus, oublie son dossier à la maison et arrive en retard. Elle affirme : « Je tenais vraiment à ce poste, c'est juste de la malchance. » Analysez cette situation à l'aide de deux concepts freudiens.
\tcblower
Solution détaillée.
Premier concept : l'\emph{acte manqué}. Selon Freud, les erreurs et oublis ne sont jamais de purs accidents : ils réalisent un désir inconscient. Ici, l'accumulation des « malchances » (mauvais bus, dossier oublié) peut s'interpréter comme l'expression d'une peur inconsciente de l'entretien — peur d'échouer, ou même peur de réussir et de devoir changer de vie.
Second concept : le \emph{refoulement}. Cette peur est inacceptable pour le sujet (« je tiens vraiment à ce poste ») : elle est donc maintenue hors de la conscience. Le Moi croit sincèrement vouloir réussir, tandis que le désir refoulé d'éviter l'épreuve agit dans l'ombre et sabote la conduite.
Conclusion : la psychanalyse montre que le sujet n'est pas « maître en sa propre maison » ; sa conduite peut obéir à des motifs qu'il ignore.
\end{exoresolu}

\courssub{Les critiques : Sartre et Popper}

La psychanalyse a profondément marqué la culture du \textsc{xx}\textsuperscript{e} siècle, mais elle a aussi suscité des objections philosophiques et scientifiques majeures. Nous retiendrons les deux plus célèbres, car elles attaquent la théorie sur deux fronts différents : la liberté (Sartre) et la scientificité (Popper).

\textbf{La critique de Sartre : l'inconscient détruit la liberté.} Jean-Paul \cle{Sartre} (1905-1980) reproche à Freud de diviser le sujet et de faire de l'homme une marionnette de forces obscures. Si mes actes sont déterminés par des désirs inconscients, alors je ne suis plus libre ni responsable : l'inconscient devient une excuse commode (« ce n'est pas moi, c'est mon inconscient »). Or, pour Sartre, l'homme est \emph{condamné à être libre} : il est entièrement responsable de ce qu'il fait de sa vie. Pour expliquer le mensonge à soi-même (par exemple la jeune femme qui se dit « malchanceuse »), Sartre n'a pas besoin de l'inconscient : il lui suffit de la notion de \cle{mauvaise foi}, attitude par laquelle la conscience se ment à elle-même en se faisant passer pour une chose déterminée. Le problème, souligne Sartre, est logique : qui, dans la théorie freudienne, opère la censure ? Il faudrait une instance qui \emph{connaît} le désir pour le refouler — on retombe donc sur une forme de conscience.

\textbf{La critique de Popper : une théorie irréfutable.} Karl \cle{Popper} (1902-1994), philosophe des sciences, formule une objection différente. Pour lui, une théorie scientifique doit être \emph{réfutable} (falsifiable) : elle doit pouvoir être mise en échec par l'expérience. Or la psychanalyse, remarque Popper, semble pouvoir expliquer \emph{tout} et son contraire. Si le patient accepte l'interprétation, cela confirme la théorie ; s'il la refuse, on dira qu'il « résiste », ce qui confirme encore la théorie. Une théorie compatible avec tous les faits possibles n'en explique réellement aucun : la psychanalyse ne serait donc pas une \emph{science} au sens strict, mais un système d'interprétation invérifiable.

\begin{center}
\begin{tabularx}{\linewidth}{|X|X|}
\hline
\rowcolor{popblueL} \textbf{Arguments en faveur de la psychanalyse} & \textbf{Critiques adressées à la psychanalyse} \\
\hline
Elle rend compte de faits troublants (rêves, lapsus, symptômes) que la conscience seule n'explique pas. & \textbf{Sartre} : l'inconscient détruit la liberté et la responsabilité ; la mauvaise foi suffit à expliquer le mensonge à soi. \\
\hline
Elle a une efficacité thérapeutique : la parole libère certains patients de leurs symptômes. & \textbf{Popper} : théorie irréfutable, donc non scientifique au sens strict ; elle explique tout et son contraire. \\
\hline
Elle a enrichi la compréhension de l'enfance, de l'éducation, de l'art et de la culture. & \textbf{Sciences cognitives} : les mécanismes freudiens (Ça, Surmoi, Œdipe) ne sont pas observables ni mesurables expérimentalement. \\
\hline
\end{tabularx}
\end{center}

\courssub{Bilan nuancé : une herméneutique plutôt qu'une science expérimentale}

Faut-il alors rejeter la psychanalyse ? Un bilan équilibré conduit à une position nuancée. La psychanalyse n'est peut-être pas une \emph{science expérimentale} au sens de la physique ou de la biologie : ses concepts ne se mesurent pas en laboratoire. Mais elle peut être comprise comme une \cle{herméneutique}, c'est-à-dire une méthode d'\emph{interprétation du sens} : elle aide l'homme à se comprendre lui-même, à donner une signification à ses souffrances, ses rêves et ses conduites. Comme le suggère Paul Ricœur, Freud est moins un savant de laboratoire qu'un grand lecteur de l'âme humaine.

Sa postérité est considérable : Jacques \cle{Lacan} (1901-1981) a relu Freud à la lumière du langage (« l'inconscient est structuré comme un langage »), et l'ensemble de la \emph{psychologie des profondeurs} (Jung, Adler, puis les psychothérapies contemporaines) hérite de l'idée freudienne selon laquelle notre vie psychique dépasse largement notre conscience claire. Même discutée, la psychanalyse a définitivement transformé l'image que l'homme se fait de lui-même : après Freud, nul ne peut plus prétendre être totalement transparent à ses propres yeux.

\begin{savaistu}[La psychanalyse en Afrique et au Cameroun]
L'idée que nos conduites ont des causes cachées n'est pas étrangère aux cultures africaines : bien avant Freud, les traditions de guérison et d'écoute de la parole (chez le guérisseur, le sage, le griot) accordaient déjà une place centrale au récit de soi et à l'interprétation des rêves. Aujourd'hui, la psychologie et la psychiatrie se développement au Cameroun, notamment à Yaoundé et à Douala, où des praticiens formés aux thérapies issues de la psychanalyse accompagnent des patients. Cela montre que la question freudienne — « que sais-je vraiment de moi-même ? » — est une question universelle.
\end{savaistu}

\begin{attention}[Erreur fréquente à éviter]
Ne présentez jamais la psychanalyse comme une \emph{vérité scientifique établie}. Elle reste une \emph{théorie discutée} : contestée par Sartre au nom de la liberté, par Popper au nom de la scientificité, et par les sciences cognitives au nom de l'expérimentation. À l'examen, une copie qui ignore ces critiques est une copie incomplète ; une copie qui ne retient que les critiques est tout aussi partiale. Autre piège : ne confondez pas le \emph{préconscient} (accessible à volonté) avec l'\emph{inconscient} (inaccessible sans travail d'interprétation).
\end{attention}

\begin{methode}[Disserter sur Freud : thèse, antithèse, synthèse]
\begin{enumerate}
\item \textbf{Thèse} : exposer la position de Freud (l'inconscient existe, il détermine une partie de nos conduites ; preuves : rêves, actes manqués, symptômes).
\item \textbf{Antithèse} : opposer les critiques (Sartre : la liberté et la mauvaise foi ; Popper : l'irréfutabilité).
\item \textbf{Synthèse personnelle justifiée} : dépasser l'opposition en distinguant les plans — la psychanalyse comme herméneutique du sens et comme pratique thérapeutique, sans prétendre au statut de science expérimentale. Justifier son point de vue par des arguments et des exemples, non par une simple opinion.
\end{enumerate}
\end{methode}

\begin{exercice}[Pour s'entraîner]
Un candidat au Probatoire technique déclare après son échec : « Ce n'est pas de ma faute, mon inconscient ne voulait pas que je réussisse. »
\begin{enumerate}
\item Expliquez ce que Freud pourrait dire de cet échec (utilisez les notions d'acte manqué et de refoulement).
\item Montrez ensuite, avec Sartre, en quoi ce raisonnement peut relever de la mauvaise foi.
\item Concluez : l'inconscient peut-il servir d'excuse ? Justifiez votre réponse.
\end{enumerate}
\end{exercice}

\begin{corrige}[Exercice : l'inconscient comme excuse]
\begin{enumerate}
\item Pour Freud, l'échec peut s'analyser comme un \emph{acte manqué} à grande échelle : si le candidat a, par exemple, négligé ses révisions, perdu sa convocation ou paniqué le jour de l'épreuve, ces « accidents » peuvent exprimer une peur inconsciente de l'examen, \emph{refoulée} parce qu'inavouable (peur de décevoir sa famille, sentiment d'illégitimité). Le désir inconscient d'éviter l'épreuve aurait alors saboté sa conduite.
\item Pour Sartre, invoquer l'inconscient est une attitude de \emph{mauvaise foi} : le candidat se fait passer pour une chose déterminée (« mon inconscient a décidé ») afin d'échapper à sa responsabilité. Or il était libre de réviser, de se présenter, de gérer son stress : il est responsable de ce qu'il a fait de sa situation.
\item Conclusion possible : l'inconscient peut \emph{expliquer} une part de nos conduites, mais il ne peut pas \emph{excuser} le sujet, car expliquer n'est pas justifier. Reconnaître les forces qui nous habitent est précisément le premier pas pour en devenir moins esclave — ce qui rejoint, paradoxalement, l'exigence sartrienne de liberté.
\end{enumerate}
\end{corrige}

\begin{aretenir}[L'essentiel de la section]
\begin{itemize}
\item Première topique : conscient, préconscient, inconscient (dynamique).
\item Seconde topique : le Ça (pulsions, principe de plaisir), le Moi (principe de réalité, médiation), le Surmoi (interdits intériorisés, culpabilité).
\item Le refoulement maintient hors de la conscience les désirs inacceptables ; ceux-ci reviennent sous forme de rêves, d'actes manqués, de symptômes.
\item Méthode : libre association, interprétation des rêves (« voie royale »), actes manqués, transfert.
\item Critiques : Sartre (l'inconscient détruit la liberté ; la mauvaise foi suffit) ; Popper (théorie irréfutable, donc non scientifique).
\item Bilan : la psychanalyse est une herméneutique du sens et une pratique thérapeutique, dont la postérité (Lacan, psychologie des profondeurs) est immense.
\end{itemize}
\end{aretenir}

\coursec{Synthèse : conscience et inconscient, complémentaires ou adversaires ?}

Au terme de notre étude de la psychanalyse freudienne et de ses critiques, une question demeure, qui est sans doute la plus importante de toute la leçon : si l'inconscient existe et gouverne une grande partie de notre vie psychique, que devient la conscience, que la tradition philosophique présentait comme le siège de la liberté, de la raison et de la responsabilité ? Faut-il choisir entre la conscience et l'inconscient, comme on choisirait entre deux adversaires ? Ou bien ces deux dimensions de la vie psychique peuvent-elles se compléter ? C'est à ce débat central que cette synthèse est consacrée.

\courssub{Le débat central : l'homme est-il « maître dans sa propre maison » ?}

Commençons par une intuition simple. Chacun de nous se vit comme un être qui décide : j'ai choisi ma filière technique, j'ai choisi mes amis, je décide de me lever le matin pour aller au lycée. Cette évidence immédiate semble prouver que la conscience gouverne ma vie. Mais l'expérience quotidienne apporte aussi des contre-épreuves troublantes : pourquoi ai-je répondu avec une colère disproportionnée à une simple remarque ? Pourquoi ai-je oublié le nom de ce voisin précisément au moment de le présenter ? Pourquoi certains élèves, malgré une volonté sincère de réussir, « sabotent-ils » leurs examens par des oublis et des retards répétés ?

Freud a donné à ce débat sa formulation la plus célèbre. Dans son essai \emph{Une difficulté de la psychanalyse} (1917), il affirme que la psychanalyse a porté à l'orgueil humain une \cle{troisième blessure narcissique} : après Copernic, qui a montré que la Terre n'est pas le centre de l'univers, et après Darwin, qui a montré que l'homme descend de l'animal, la psychanalyse montre que \og le Moi n'est pas maître dans sa propre maison \fg{}. Autrement dit, la conscience, que l'homme croyait souveraine, ne serait qu'une petite partie éclairée d'un psychisme dont l'essentiel lui échappe.

\begin{savaistu}[Les trois blessures de l'orgueil humain selon Freud]
Freud comparait sa découverte à deux révolutions scientifiques antérieures. Première blessure, \emph{cosmologique} : Copernic (XVI\textsuperscript{e} siècle) décentre la Terre, qui tourne autour du Soleil. Deuxième blessure, \emph{biologique} : Darwin (XIX\textsuperscript{e} siècle) rattache l'homme au règne animal. Troisième blessure, \emph{psychologique} : la psychanalyse décentre le Moi, qui découvre qu'il n'est pas maître chez lui, car ses pensées et ses actes dépendent de forces inconscientes. Pour Freud, chaque progrès de la science a ainsi obligé l'humanité à plus de modestie.
\end{savaistu}

La question philosophique est donc précise : la conscience est-elle le fondement de la vie morale et de la connaissance (thèse rationaliste), ou n'est-elle que la surface visible d'un psychisme gouverné par l'inconscient (thèse freudienne) ? Et surtout : peut-on dépasser cette opposition ?

\courssub{La position rationaliste : la conscience, fondement de la connaissance et de la morale}

La tradition philosophique classique fait de la conscience le principe absolu de la vie humaine.

\textbf{Descartes} en donne le fondement avec le \emph{cogito} : \og je pense, donc je suis \fg{}. Même si je doute de tout, je ne peux pas douter que je pense, et donc que j'existe en tant que chose pensante. La conscience de soi est ainsi la première vérité, indubitable, sur laquelle toute connaissance peut être reconstruite. Pour Descartes, l'esprit est entièrement transparent à lui-même : il n'y a rien en moi que je ne puisse, en principe, connaître par la réflexion.

\textbf{Sartre}, au XX\textsuperscript{e} siècle, radicalise cette position dans \emph{L'Être et le Néant} (1943). Pour lui, la conscience est \cle{translucide} : elle est entièrement ce qu'elle se sait être. Il n'existe pas d'inconscient, car un psychisme dont une partie échapperait à la conscience serait une contradiction : qui ferait le travail de censure, sinon une conscience qui sait ce qu'elle censure ? Sartre parle alors de \cle{mauvaise foi} : l'homme se ment à lui-même en se faisant passer pour déterminé (\og je ne peux pas m'en empêcher, c'est plus fort que moi \fg{}), alors qu'il est en réalité \og condamné à être libre \fg{}. Chacun est responsable de la totalité de ses actes, car chacun se choisit à chaque instant.

\begin{aretenir}[Thèse rationaliste]
Pour Descartes et Sartre, la conscience est le fondement de la connaissance (le cogito) et de la morale (la responsabilité). Sans conscience claire de soi, pas de vérité certaine ni de liberté. L'inconscient freudien est donc rejeté : il menacerait à la fois la connaissance et la responsabilité morale.
\end{aretenir}

\courssub{La position freudienne : la conscience n'est que la surface du psychisme}

Face à cette tradition, Freud soutient une thèse inverse. La conscience n'est qu'une petite partie du psychisme, comparable à la partie émergée d'un iceberg : l'essentiel — le Ça (réservoir des pulsions) et une grande partie du Surmoi (interdits intériorisés) — est inconscient. Le Moi conscient est tiraillé entre ces forces qu'il ne contrôle pas.

Les preuves cliniques invoquées par Freud sont nombreuses : les rêves, les lapsus, les actes manqués, les symptômes névrotiques montrent que nos conduites obéissent à des motivations cachées, refoulées parce qu'inacceptables pour la conscience morale. L'homme ne sait pas toujours pourquoi il agit ; il se raconte après coup des raisons (\og rationalisations \fg{}) qui masquent les vraies causes pulsionnelles.

\begin{attention}[Piège à éviter : deux « choses » séparées]
Ne présente jamais la conscience et l'inconscient comme deux objets distincts, deux boîtes ou deux petits personnages qui se disputeraient dans la tête. Ce sont deux \emph{dimensions d'une même vie psychique} : le même désir, le même souvenir, peut être tantôt conscient, tantôt refoulé dans l'inconscient. Dire \og mon inconscient m'a fait faire cela \fg{} comme s'il s'agissait d'un étranger en moi, c'est déjà mal comprendre Freud : l'inconscient, c'est encore \emph{moi}.
\end{attention}

\courssub{Schéma-bilan : le Moi entre conscience et inconscient}

\begin{popfigure}
\begin{center}
\begin{tikzpicture}[scale=1.0]
% Socle de la balance
\fill[popdark] (0,-0.15) rectangle (0.5,0);
\draw[popdark,thick] (0.25,0) -- (0.25,2.2);
\draw[popdark,thick] (-3.4,2.2) -- (3.9,2.2);
% Fléau : point d'équilibre = le Moi
\fill[popgold] (0.25,2.2) circle (0.14);
\node[above,popdark] at (0.25,2.42) {\textbf{Le Moi}};
% Plateau gauche : conscience
\draw[popblue,thick] (-2.9,2.2) -- (-3.4,1.35);
\draw[popblue,thick] (-2.9,2.2) -- (-2.4,1.35);
\draw[popblue,thick,fill=popblueL] (-3.55,1.35) arc (180:360:0.65) -- cycle;
\node[popblue,align=center] at (-2.9,0.35) {\textbf{Conscience}\\ \footnotesize liberté\\ \footnotesize responsabilité\\ \footnotesize raison};
% Plateau droit : inconscient
\draw[poppurple,thick] (3.4,2.2) -- (2.9,1.35);
\draw[poppurple,thick] (3.4,2.2) -- (3.9,1.35);
\draw[poppurple,thick,fill=poppurpleL] (2.75,1.35) arc (180:360:0.65) -- cycle;
\node[poppurple,align=center] at (3.4,0.35) {\textbf{Inconscient}\\ \footnotesize déterminisme\\ \footnotesize pulsions\\ \footnotesize refoulement};
% Flèches de tension
\draw[-{Stealth},popmuted] (-2.9,1.15) -- (-2.9,0.85);
\draw[-{Stealth},popmuted] (3.4,1.15) -- (3.4,0.85);
\end{tikzpicture}
\end{center}
\legende{Schéma-bilan : le Moi (point d'équilibre, en doré) est suspendu entre les exigences de la conscience (liberté, responsabilité, raison) et les forces de l'inconscient (déterminisme, pulsions, refoulement). La vie psychique est cette tension permanente, non la victoire d'un plateau sur l'autre.}
\end{popfigure}

Ce schéma résume le débat : ni la toute-puissance de la conscience (le plateau gauche écraserait tout), ni l'esclavage absolu sous l'inconscient (le plateau droit l'emporterait) ne rendent compte de l'expérience humaine. Le Moi vit dans la tension entre les deux.

\courssub{Vers une synthèse : la liberté éclairée}

L'opposition entre rationalistes et freudiens semble donc totale. Pourtant, une synthèse est possible, et c'est Freud lui-même qui en donne la clé : le but de la psychanalyse n'est pas de constater notre esclavage, mais de le réduire. Sa formule célèbre — \og où était le Ça, le Moi doit advenir \fg{} (\emph{Wo Es war, soll Ich werden}) — signifie que l'analyse vise à rendre conscient ce qui était inconscient, donc à \emph{élargir} le domaine de la conscience.

\begin{definition}[La liberté éclairée]
On appelle \cle{liberté éclairée} une liberté qui s'accroît par la connaissance de ses propres déterminations. Je ne suis pas libre en ignorant les forces qui m'agissent, mais en les connaissant : comprendre mes peurs, mes colères, mes désirs refoulés me permet de ne plus les subir aveuglément. Cette idée fait écho à Spinoza, pour qui nous sommes d'autant plus libres que nous connaissons les causes de nos passions (complément utile pour le Baccalauréat technique).
\end{definition}

La synthèse se formule alors ainsi : l'inconscient \emph{limite} la conscience, mais la connaissance de l'inconscient \emph{élargit} la conscience. Le rationaliste a raison de dire que la conscience est le fondement de la morale ; le freudien a raison de dire que cette conscience est naturellement étroite et trompée. La conclusion pratique : la liberté n'est pas un donné, c'est une \emph{conquête} — elle se gagne par le travail de réflexion sur soi, et parfois par le travail analytique.

\courssub{Application à la vie courante : se connaître pour mieux se gouverner}

Cette synthèse n'est pas abstraite : elle transforme la manière de vivre des situations quotidiennes.

\begin{itemize}
\item \textbf{La gestion de la colère.} Un élève explose de colère pour une remarque anodine d'un camarade. S'il prend le temps de réfléchir, il découvrira peut-être que cette colère déborde la situation présente : elle réactive une humiliation ancienne, une tension familiale. Comprendre cette motivation cachée ne supprime pas la colère, mais permet de ne plus la déverser sur un innocent.
\item \textbf{Les peurs inexpliquées.} La peur de prendre la parole en public, de passer un examen, de conduire : ces peurs ont souvent des racines que la conscience ignore. Les identifier (souvenir d'un échec, parole blessante d'un adulte) est le premier pas pour les surmonter.
\item \textbf{Les choix d'orientation.} Choisir une filière technique pour plaire à son père, ou pour fuir une matière où l'on a été humilié, ce n'est pas choisir librement au sens fort. S'interroger honnêtement sur ses motivations profondes rend le choix plus authentique — donc plus libre.
\end{itemize}

\begin{exoresolu}[Analyse d'une situation de vie courante]
\textbf{Énoncé.} Marlyse, élève de Première technique, rate systématiquement ses évaluations de mathématiques alors qu'elle travaille beaucoup. Elle dit : \og C'est le destin, je suis nulle, je n'y peux rien. \fg{} En discutant avec un conseiller d'orientation, elle se souvient qu'en classe de Sixième, un enseignant l'avait ridiculisée devant toute la classe pour une erreur au tableau. Analyse cette situation à l'aide des notions de la leçon.
\tcblower
\textbf{Solution détaillée.}
\begin{enumerate}
\item \emph{Repérer le discours de mauvaise foi.} En invoquant le \og destin \fg{}, Marlyse se présente comme une chose déterminée, ce qui la dispense de chercher la vraie cause. Sartre appellerait cela de la mauvaise foi : elle se croit non libre pour ne pas avoir à affronter le problème.
\item \emph{Identifier le travail de l'inconscient.} L'échec répété malgré le travail suggère un blocage : la situation d'évaluation réactive inconsciemment le souvenir refoulé de l'humiliation de Sixième. L'angoisse paralyse ses capacités au moment de l'épreuve.
\item \emph{Appliquer la synthèse.} En prenant conscience, grâce au dialogue, de l'origine de son blocage, Marlyse élargit sa conscience : ce qui était subi (l'angoisse inexpliquée) devient connu et nommé. Elle peut alors travailler ce souvenir, se détacher du jugement de l'ancien enseignant, et retrouver progressivement ses moyens.
\item \emph{Conclure.} Cette situation illustre la liberté éclairée : Marlyse ne devient pas libre en niant ses déterminations, mais en les connaissant. L'inconscient limitait sa conscience ; la prise de conscience élargit sa liberté.
\end{enumerate}
\end{exoresolu}

\courssub{Enjeu moral et civique : responsabilité et tolérance}

Le débat entre conscience et inconscient a des conséquences directes en morale et en droit.

D'un côté, la \cle{responsabilité} suppose la conscience : on ne peut être tenu pour responsable que de ce qu'on a fait en connaissance de cause. C'est le fondement de toute vie morale et de toute justice : punir quelqu'un qui ne savait pas ce qu'il faisait serait injuste.

De l'autre côté, la découverte de l'inconscient invite à la \emph{tolérance} et à l'\emph{écoute} : si l'homme n'est pas totalement maître de lui-même, alors le jugement moral doit être prudent. Celui qui agit mal n'est pas toujours un être pervers qui aurait froidement choisi le mal ; il est parfois prisonnier de mécanismes qui le dépassent.

\begin{exemplebox}[L'expertise psychiatrique en droit]
Le droit reconnaît implicitement les limites de la conscience. Au Cameroun, comme dans la plupart des systèmes juridiques, le Code pénal prévoit qu'un trouble psychique peut atténuer, voire abolir, la responsabilité pénale : avant le procès d'une personne soupçonnée de troubles mentaux, le juge peut ordonner une \cle{expertise psychiatrique}. Si les experts concluent que l'accusé était atteint, au moment des faits, d'un trouble ayant aboli son discernement, il ne peut être condamné : il sera orienté vers des soins plutôt que vers la prison. Si le trouble a seulement \emph{altéré} le discernement, la peine peut être atténuée. Le droit admet ainsi que la conscience, fondement de la responsabilité, n'est pas toujours entière — ce qui rejoint, sans le dire, la leçon de Freud.
\end{exemplebox}

L'enjeu civique est donc double : ne pas utiliser l'inconscient comme excuse universelle (\og ce n'est pas moi, c'est mon inconscient \fg{} — ce serait une nouvelle mauvaise foi), mais ne pas non plus juger autrui avec la sévérité naïve de celui qui croit que tout le monde est toujours pleinement maître de ses actes. La maladie mentale, en particulier, doit être accueillie comme une souffrance à soigner, non comme une faute à punir ou une honte à cacher — attitude encore trop répandue dans certaines de nos communautés.

\courssub{Carte mentale récapitulative de la leçon}

\begin{popfigure}
\begin{center}
\begin{tikzpicture}[
  mindmap,
  grow cyclic,
  every node/.style={concept},
  root concept/.append style={concept color=popdark, font=\large\bfseries},
  level 1/.append style={level distance=3.4cm, sibling angle=90, font=\bfseries},
  level 2/.append style={level distance=2.4cm, sibling angle=45, font=\footnotesize},
  concept color=popblue
]
\node[root concept]{Conscience et inconscient}
  child[concept color=popblue]{ node {Conscience}
    child { node {Types : spontanée, réfléchie, morale} }
    child { node {Caractéristiques : intentionnalité, clarté, unité} }
    child { node {Valeur : connaissance, liberté, responsabilité} }
    child { node {Limites : erreurs, illusions, oublis} }
  }
  child[concept color=poppurple]{ node {Inconscient}
    child { node {Découverte : hypnose, rêves, lapsus} }
    child { node {Types : neurologique, freudien, collectif} }
    child { node {Psychanalyse : Ça, Moi, Surmoi ; refoulement} }
    child { node {Critiques : Sartre, scientificité discutée} }
  };
\end{tikzpicture}
\end{center}
\legende{Carte mentale récapitulative : les deux branches de la leçon et leurs quatre sous-parties respectives. La synthèse (liberté éclairée) relie les deux branches : connaître l'inconscient élargit la conscience.}
\end{popfigure}

\courssub{Méthode : traiter une dissertation sur ce thème}

Les sujets de dissertation possibles sur cette leçon sont nombreux : \og Suis-je maître de moi-même ? \fg{}, \og La conscience suffit-elle à définir l'homme ? \fg{}, \og La connaissance de l'inconscient nous rend-elle plus libres ? \fg{}

\begin{methode}[Structure type d'une dissertation dialectique sur ce thème]
\begin{enumerate}
\item \textbf{Introduction} : partir d'une situation concrète (un lapsus, une colère inexpliquée), dégager le problème : la conscience gouverne-t-elle réellement notre vie ? Annoncer le plan.
\item \textbf{Thèse — Toute-puissance de la conscience} : le cogito de Descartes, la conscience comme fondement de la connaissance et de la morale ; Sartre et la responsabilité totale. Sans conscience, pas de sujet libre.
\item \textbf{Antithèse — Mise en échec par l'inconscient} : les faits (rêves, lapsus, névroses) ; la topique freudienne ; \og le Moi n'est pas maître dans sa propre maison \fg{}. La conscience se trompe sur ses propres motifs.
\item \textbf{Synthèse — Une conscience élucidée} : l'inconscient limite la conscience, mais le connaître l'élargit ; la liberté comme conquête (liberté éclairée, écho à Spinoza) ; conséquences morales : responsabilité maintenue, mais jugement tempéré par la tolérance.
\item \textbf{Conclusion} : répondre nettement à la question posée et ouvrir (par exemple vers la mémoire et le temps).
\end{enumerate}
\end{methode}

\begin{attention}[Erreurs fréquentes en dissertation]
\begin{itemize}
\item Ne transforme pas la deuxième partie en exposé admiratif de Freud : il faut \emph{discuter} (critiques de Sartre, débat sur la scientificité de la psychanalyse).
\item N'affirme jamais que l'inconscient \og excuse tout \fg{} : la synthèse doit maintenir la responsabilité, sinon la conclusion devient immorale.
\item Évite la synthèse artificielle (\og les deux ont raison \fg{}) : il faut montrer \emph{comment} la connaissance de l'inconscient renforce la conscience au lieu de la détruire.
\end{itemize}
\end{attention}

\courssub{Conclusion et ouverture}

Conscience et inconscient ne sont donc ni de purs adversaires ni de simples associés : ils sont les deux dimensions d'une même vie psychique, unies par un rapport de tension féconde. L'inconscient détrône la conscience de sa prétention à la toute-puissance ; mais la conscience, en accueillant la vérité de l'inconscient, devient plus vaste, plus lucide et plus libre qu'elle ne l'était. L'homme n'est peut-être pas maître dans sa propre maison, mais il peut apprendre à en connaître les pièces obscures — et c'est déjà une forme supérieure de maîtrise.

Cette réflexion ouvre naturellement vers la leçon suivante, consacrée au \cle{temps}. Car comment l'inconscient conserve-t-il les souvenirs refoulés, sinon par la \emph{mémoire} ? Et comment la conscience se rapporte-t-elle à son passé, à son présent et à son avenir ? La mémoire est précisément le lieu où la conscience rencontre son passé — y compris ce passé refoulé que la psychanalyse nous a appris à écouter. Comprendre le temps, ce sera comprendre comment nous devenons ce que nous sommes.

\newpage

\coursec{Activité d'intégration}

\coursec{La conscience et l'inconscient}

\courssub{Objectifs de l'activité}
\begin{itemize}
    \item \textbf{Mobiliser} les notions de \cle{conscience} (immédiate, réflexive, morale) et d'\cle{inconscient} (topique, dynamique, freudien) pour analyser des conduites concrètes.
    \item \textbf{Construire} et \textbf{légender} les deux schémas fondamentaux de la topique freudienne : l'\cle{iceberg} (topographie) et les \cle{trois instances} (Ça, Moi, Surmoi).
    \item \textbf{Confronter} l'explication causaliste de Freud (déterminisme psychique, refoulement) à l'explication finaliste de Sartre (\cle{mauvaise foi}, liberté absolue) sur le même terrain : les \cle{actes manqués}.
    \item \textbf{Rédiger} une prise de position argumentée (15 lignes), structurée par une thèse, des arguments étayés par les documents et une ouverture critique.
\end{itemize}

\courssub{Situation}
\label{sit:tp_inconscient}
Vous êtes élèves en classe de Première Technique (spécialités tertiaires, industrielles, sciences et technologies). Dans le cadre du module « \emph{La conscience et l'inconscient} », votre professeur de Philosophie organise un TP d'intégration à partir d'un dossier documentaire. L'objectif est de dépasser la simple définition des concepts pour les mettre à l'épreuve du réel camerounais.

\textbf{Dossier documentaire}

\textbf{Document 1 : Trois récits de la vie courante à Yaoundé et Douala}

\begin{enumerate}
    \item \textbf{Récit A (Marché Mokolo, Yaoundé).} Monsieur \textsc{Kamga}, commerçant en tissus depuis 20 ans, présente son nouveau stock à un client important. Au moment de citer le nom de son principal concurrent, \textsc{M. Fouda}, installé à 50 mètres, il bloque : « \emph{Euh... le monsieur d'en face, vous savez, celui qui...} ». Le nom lui revient 10 minutes plus tard, seul dans son bureau. M. Kamga affirme n'avoir « \emph{aucune raison d'oublier ce nom} ».
    \item \textbf{Récit B (Lycée technique de Nkolbisson).} \textsc{Ngo}, élève en Terminale G2, prépare le Probatoire. Depuis trois semaines, elle fait le même cauchemar : elle arrive dans la salle d'examen, le sujet est distribué, elle ouvre la copie... et réalise qu'elle a révisé le mauvais programme (Mathématiques au lieu de Physique). Elle se réveille en sueur, le cœur battant. Le jour de l'épreuve de Physique, elle est calme et réussit.
    \item \textbf{Récit C (Quartier Bonamoussadi, Douala).} \textsc{M. Essomba}, 45 ans, cadre bancaire, surprend son fils cadet (14 ans) en train de fumer une cigarette derrière le garage. Il entre dans une colère violente, le gifle, le prive de sortie pendant un mois et lui crie : « \emph{Tu salis le nom de la famille ! Moi, à ton âge, je n'aurais jamais osé !} ». Son épouse, témoin de la scène, lui murmure plus tard : « \emph{Toi non plus tu n'as pas osé... tu le faisais en cachette chez ton oncle à Bafia.} » M. Essomba pâlit, se tait, mais ne revient pas sur la sanction.
\end{enumerate}

\textbf{Document 2 : Extrait de \emph{La Psychopathologie de la vie quotidienne} (1901), S. Freud}
\begin{quote}
« \emph{L'oubli de noms propres, même de ceux qui nous sont familiers, n'est pas un phénomène accidentel. Il repose sur un motif de refoulement. Le nom oublié est en liaison avec un complexe refoulé ; l'oubli est la conséquence d'une résistance qui s'oppose à la mise en conscience de ce complexe. L'acte manqué n'est pas une défaillance de la machine psychique, mais l'expression d'une intention cachée, d'un désir inconscient qui contredit l'intention consciente.} » 
\end{quote}
\textit{(Adaptation pédagogique : Freud y analyse l'oubli temporaire du nom « Signorelli » par un voyageur, lié à une association refoulée avec la mort et la sexualité.)}

\textbf{Document 3 : Extrait de \emph{L'Être et le Néant} (1943), J.-P. Sartre — La mauvaise foi}
\begin{quote}
« \emph{La mauvaise foi est une fuite devant la responsabilité. L'homme veut se croire déterminé (comme un objet, un \emph{en-soi}) pour ne pas assumer sa liberté (son \emph{pour-soi}). Celui qui dit "Je n'ai pas pu faire autrement, c'est plus fort que moi" ment à lui-même : il \emph{sait} qu'il choisit, mais il \emph{veut} ne pas le savoir. L'inconscient freudien est une ruse pour sauver la responsabilité : on invente un "Ça" qui décide à ma place. Mais la conscience est toujours conscience \emph{de} quelque chose ; elle est transparence. Il n'y a pas de "derrière" de la conscience.} » 
\end{quote}
\textit{(Adaptation pédagogique : Sartre récuse l'inconscient comme instance causale ; il y voit un projet de fuite de la liberté.)}

\textbf{Document 4 : Données synthétiques — Sommeil, rêve et activité psychique}
\begin{center}
\begin{tabularx}{\linewidth}{|l|X|X|}\hline
\textbf{Paramètre} & \textbf{Donnée moyenne (Adulte)} & \textbf{Interprétation psychanalytique} \\ \hline
Durée totale de sommeil & 7 h à 8 h / 24 h & Temps de « suspension » du Moi conscient. \\ \hline
Durée du sommeil paradoxal (REM) & $\approx$ 1 h 30 à 2 h (20--25\% du sommeil) & Phase principale du rêve ; « voie royale vers l'inconscient » (Freud). \\ \hline
Nombre de cycles de sommeil/nuit & 4 à 6 cycles de 90 min & Alternance travail psychique conscient / travail inconscient. \\ \hline
Pourcentage de rêves oubliés au réveil & $>$ 90\% & Efficacité du \cle{refoulement} / censure du Moi. \\ \hline
Activité cérébrale (IRMf) en REM & Hyperactivité limbique (émotions), hypoactivité cortex préfrontal (contrôle, logique) & Base neurologique de l'irruption du \cle{Ça} (pulsions) sans frein du \cle{Moi} (réalité). \\ \hline
\end{tabularx}
\end{center}
\textit{Source : Synthèse de données INSERM / Société Française de Recherche sur le Sommeil, adaptée.}

\courssub{Consignes}
Travail en groupes de 4 élèves (20 min), puis restitution et correction collective (25 min), rédaction individuelle finale (15 min).

\begin{enumerate}
    \item \textbf{Analyse des récits (Doc 1).} Pour chacun des trois récits (A, B, C), identifiez : 
    \begin{itemize}
        \item ce qui relève de la \cle{conscience} (intention explicite, perception, jugement moral affiché) ;
        \item ce qui pourrait relever de l'\cle{inconscient} (désir refoulé, angoisse symbolisée, identification projetée, acte manqué).
    \end{itemize}
    Justifiez chaque identification par un élément précis du texte.
    
    \item \textbf{Modélisation graphique (Doc 1 \& 2).} Sur papier millimétré (ou au tableau), tracez deux schémas distincts, légendés et fléchés :
    \begin{enumerate}
        \item Le schéma de l'\cle{iceberg} (Topographie : Conscient / Préconscient / Inconscient). Placez-y : le nom de M. Fouda (Récit A), le contenu du cauchemar (Récit B), le souvenir refoulé de M. Essomba (Récit C).
        \item Le schéma des \cle{trois instances} (Ça, Moi, Surmoi). Placez-y les \emph{moteurs} de l'action pour chaque récit : ex. « Pulsion agressive » (Ça), « Crainte du ridicule / adaptation » (Moi), « Idéal moral / interdiction » (Surmoi).
    \end{enumerate}
    
    \item \textbf{Débat interprétatif (Doc 2 \& 3).} En vous appuyant sur les extraits de Freud et de Sartre, discutez la question : 
    \emph{« Faut-il expliquer la conduite de M. Kamga (oubli), de Ngo (cauchemar récurrent) et de M. Essomba (sévérité disproportionnée) par l'inconscient freudien (déterminisme, refoulement) ou par la mauvaise foi sartrienne (choix de se mentir, fuite de la liberté) ? »}
    Notez les arguments clés pour les deux thèses.
    
    \item \textbf{Rédaction individuelle (Synthèse).} Rédigez une conclusion argumentée de \textbf{15 lignes maximum} (environ 180 mots) répondant à la question du débat. Votre texte doit comporter : une thèse claire, deux arguments distincts appuyés sur les documents, une objection (antithèse) et une réponse (synthèse), une ouverture.
\end{enumerate}

\courssub{Ressources à mobiliser}
\begin{definition}[Conscience — Types et caractéristiques]
\textbf{Conscience immédiate (spontanée) :} présence à soi vécue sans réflexion (ex. : « J'ai mal », « Je vois rouge »).
\textbf{Conscience réflexive :} le sujet se prend lui-même pour objet de connaissance (ex. : « Je sais que je sais », « J'analyse ma colère »).
\textbf{Conscience morale :} capacité à juger ses actes au regard de valeurs (bien/mal, juste/injuste). \textbf{Caractères :} unité (un seul « Je »), transparence (idéal cartésien : « Je ne peux pas ne pas savoir ce que je pense »), intentionalité (toujours conscience \emph{de} quelque chose).
\end{definition}

\begin{definition}[Inconscient — Types et découverte]
\textbf{Inconscient topique (Freud, 1900) :} lieu psychique (le « Ça ») où sont refoulés les désirs inacceptables. \textbf{Inconscient dynamique :} processus de refoulement, de déplacement, de condensation. \textbf{Découverte :} hypnose (Charcot, Janet), actes manqués, rêves, symptômes névrotiques. \textbf{Preuve freudienne :} l'inconscient \emph{agit} (déterminisme psychique) sans que le sujet en ait la maîtrise.
\end{definition}

\begin{propriete}[Modèles freudiens — Démonstration exigible au Probatoire : savoir tracer et légender les deux schémas]
\textbf{1. Première topique (1900) — L'Iceberg :}
\begin{itemize}
    \item \textbf{Conscient (sommet) :} perceptions, pensées actuelles.
    \item \textbf{Préconscient (ligne de flottaison) :} souvenirs accessibles, refoulables.
    \item \textbf{Inconscient (base immergée) :} refoulé, pulsions, traces mnésiques chargées d'affect. \emph{Règle :} l'inconscient ne devient conscient qu'après transformation (censure).
\end{itemize}
\textbf{2. Seconde topique (1923) — Les Trois Instances :}
\begin{itemize}
    \item \textbf{Ça (Id) :} pôle pulsionnel, inné, amoral, principe de plaisir, énergie libre.
    \item \textbf{Moi (Ego) :} instance de réalité, médiatrice, principe de réalité, défense (refoulement, sublimation, formation réactionnelle).
    \item \textbf{Surmoi (Superego) :} instance morale, intériorisation des interdits parentaux/sociaux, idéal du moi, culpabilité.
\end{itemize}
\textbf{Conflit :} Le Moi sert « trois maîtres sévères » (Monde extérieur, Ça, Surmoi) $\rightarrow$ angoisse $\rightarrow$ mécanismes de défense.
\end{propriete}

\begin{aretenir}[Sartre vs Freud : Le statut de l'inconscient]
\textbf{Freud :} L'inconscient est une \textbf{cause} (déterminisme). Le sujet est \emph{divisé} : « Le Moi n'est pas maître en sa propre maison. » Responsabilité atténuée.
\textbf{Sartre :} L'inconscient est une \textbf{illusion} (mauvaise foi). Le sujet est \emph{un} (pour-soi), radicalement libre. « L'homme est condamné à être libre. » Responsabilité totale.
\end{aretenir}

\courssub{Démarche et corrigé commenté}
\begin{exoresolu}[Corrigé guidé de l'activité d'intégration]
\textbf{Énoncé de l'activité :} Analyse des trois récits (Doc 1) à l'aide des concepts freudiens (Doc 2, 4) et sartriens (Doc 3), modélisation graphique (iceberg, trois instances), débat d'interprétation et rédaction synthétique (15 lignes).
\tcblower
\textbf{Étape 1 — Analyse des récits (Consigne 1)}

\begin{center}
\begin{tabularx}{\linewidth}{|l|X|X|}\hline
\textbf{Récit} & \textbf{Conscience (Manifeste)} & \textbf{Inconscient (Latent / Refoulé)} \\ \hline
\textbf{A. M. Kamga (Oubli nom)} & Intention : présenter son stock, relation commerciale normale. Perception : le client est là. Jugement : « Je n'ai pas de raison d'oublier. » & \textbf{Acte manqué (lapsus memoriae).} Refoulement de l'hostilité/concurrence envers M. Fouda. Le nom « Fouda » est lié à un complexe (peur de l'échec, jalousie) que le Moi censure. Le retour différé (10 min) marque la levée temporaire de la censure (Moi détendu). \\ \hline
\textbf{B. Ngo (Cauchemar)} & Inquiétude consciente : réussite au Probatoire. Volonté : réviser, se préparer. Jugement : « Je suis prête. » & \textbf{Travail du rêve (via Doc 4).} Contenu manifeste : erreur de salle/matière. Pensées latentes (inconscientes) : \emph{angoisse d'incompétence}, \emph{désir inconscient d'échapper à l'épreuve} (régression), \emph{peur du jugement parental}. Répétition = insistance du refoulé. Réveil en sueur = angoisse non symbolisée. Succès réel = Moi fort qui a géré l'angoisse par le rêve (fonction de garde du sommeil). \\ \hline
\textbf{C. M. Essomba (Sévérité)} & Colère morale affichée : « Tu salis la famille. » Jugement de valeur : « Moi, je n'aurais jamais osé. » Volonté éducative : sanctionner. & \textbf{Formation réactionnelle + Projection.} M. Essomba a \emph{réellement} fumé jeune (témoignage épouse). Son Surmoi tyrannique lui interdit d'assumer ce passé. Il \emph{projette} sa propre faute refoulée sur le fils. La sévérité excessive (gifle, 1 mois) = surcompensation du Moi pour faire taire la culpabilité inconsciente (Ça : « J'ai joui de transgresser » ; Surmoi : « Tu es impur »). Pâlement = reconnaissance inconsciente (Préconscient) de la vérité. \\ \hline
\end{tabularx}
\end{center}

\textbf{Étape 2 — Modélisation graphique (Consigne 2)}

\begin{popfigure}
\centering
\begin{tikzpicture}[scale=0.9, every node/.style={font=\small}]
% Iceberg container
\draw[popblue, very thick] (0,0) -- (0,-5) -- (5,-5) -- (5,0) -- cycle;
\fill[popblue!20] (0,0) rectangle (5,-5);
% Iceberg shape
\draw[popdark, thick, fill=popblue!10] (1.5,-0.5) .. controls (1,-1) and (1,-3) .. (1.5,-4) .. controls (2.5,-4.5) and (4,-4) .. (4,-3) .. controls (4.5,-2) and (4.5,-1) .. (4,-0.5) -- cycle;
% Waterline
\draw[popblue, dashed, thick] (0,-1) -- (5,-1) node[right, text=popblue] {Ligne de flottaison = Seuil de conscience};
% Labels zones
\node[popdark, align=center] at (2.75, -0.3) {\textbf{CONSCIENT}\\(Perceptions, pensées actuelles)};
\node[popdark, align=center] at (2.75, -1.8) {\textbf{PRÉCONSCIENT}\\(Souvenirs accessibles, refoulables)};
\node[popdark, align=center] at (2.75, -3.5) {\textbf{INCONSCIENT}\\(Refoulé, Pulsions, Traces mnésiques)};
% Arrows for items
\draw[->, poporange, thick] (5.2, -0.7) -- (4.5, -0.7) node[right, align=left, text width=3cm] {Récit A : Intention commerciale\\ (M. Kamga)};
\draw[->, poporange, thick] (5.2, -1.5) -- (4.5, -1.5) node[right, align=left, text width=3cm] {Récit A : Nom « Fouda »\\ (avant retour)};
\draw[->, poporange, thick] (5.2, -2.5) -- (4.5, -2.5) node[right, align=left, text width=3cm] {Récit B : Contenu latent\\ du cauchemar (angoisse)};
\draw[->, poporange, thick] (5.2, -3.8) -- (4.5, -3.8) node[right, align=left, text width=3cm] {Récit C : Souvenir fumé\\ refoulé (M. Essomba)};
\end{tikzpicture}
\legende{Schéma de l'iceberg (Première topique freudienne) : localisation des éléments des récits.}
\end{popfigure}

\begin{popfigure}
\centering
\begin{tikzpicture}[scale=0.85, every node/.style={font=\small}]
% Three instances - Venn diagram style
\def\r{2.2}
\def\offset{1.5}
% Ca (Id) - Bottom
\draw[poporange, thick, fill=poporange!15] (0,0) circle (\r) node[below=0.2cm] {\textbf{ÇA (Id)}};
\node[poporange, align=center, text width=2.8cm] at (0, -0.5) {Pulsions (Vie/Mort)\\\emph{Principe de plaisir}\\\textbf{Énergie libre}};
% Moi (Ego) - Middle right
\draw[popblue, thick, fill=popblue!15] (\offset, 0.8) circle (\r) node[above=0.2cm] {\textbf{MOI (Ego)}};
\node[popblue, align=center, text width=2.8cm] at (\offset, 0.3) {Réalité / Adaptation\\\emph{Principe de réalité}\\\textbf{Défenses (Refoulement...)}};
% Surmoi (Superego) - Top middle
\draw[poppurple, thick, fill=poppurple!15] (0.75, 2.2) circle (\r) node[above=0.2cm] {\textbf{SURMOI}};
\node[poppurple, align=center, text width=2.8cm] at (0.75, 1.7) {Idéal moral / Interdits\\\emph{Culpabilité / Orgueil}\\\textbf{Intériorisation parentale}};

% Intersection label
\node[popdark, font=\tiny, align=center] at (0.75, 0.2) {CONFLIT};
% Arrows for narratives
\draw[->, poporange, thick] (-3.5, -0.5) -- (-1.8, -0.2) node[left, align=right, text width=3cm] {Récit B : Pulsion\\de fuite (Ça)};
\draw[->, popblue, thick] (4.5, 1.5) -- (3.0, 1.2) node[right, align=left, text width=3cm] {Récit A : Censure\\du nom (Moi)};
\draw[->, poppurple, thick] (0.75, 4.5) -- (0.75, 3.5) node[above, align=center, text width=3cm] {Récit C : Idéal moral\\tyrannique (Surmoi)};
\draw[->, popdark, thick] (-3.5, 2.5) -- (-1.5, 1.8) node[left, align=right, text width=3cm] {Récit C : Projection\\de la faute (Moi)};
\end{tikzpicture}
\legende{Schéma des trois instances (Seconde topique) : moteurs psychiques en jeu dans les récits.}
\end{popfigure}

\textbf{Étape 3 — Débat : Freud vs Sartre (Consigne 3)}

\begin{center}
\begin{tabularx}{\linewidth}{|>{\bfseries}X|X|X|}\hline
\textbf{Thèse} & \textbf{Arguments Freud (Doc 2, 4)} & \textbf{Arguments Sartre (Doc 3)} \\ \hline
\textbf{Explication par l'Inconscient (Déterminisme)} & \begin{itemize}\item \textbf{Causalisme :} L'oubli (A), le rêve (B), la projection (C) ont une \emph{cause} psychique antérieure (refoulement, pulsion). Le sujet \emph{subit} son histoire.
\item \textbf{Preuve structurelle :} La répétition (cauchemar B), le retour du refoulé (nom A, souvenir C), l'efficacité de la censure (Doc 4 : 90\% rêves oubliés, hypoactivité préfrontale) prouvent un fonctionnement \emph{autonome} de l'inconscient.
\item \textbf{Thérapie :} Comprendre la cause (levée du refoulement) apaise (Ngo calme le jour J). Le sujet n'est pas responsable de l'émergence du symptôme.
\end{itemize} & \begin{itemize}\item \textbf{Finalisme :} Ces conduites ont un \emph{sens} choisi : M. Kamga \emph{veut} ignorer son rival (stratégie d'évitement), Ngo \emph{veut} se rassurer par l'échec imaginaire (magie), M. Essomba \emph{veut} se croire pur en punissant l'autre.
\item \textbf{Transparence :} M. Essomba \emph{sait} qu'il a fumé (pâlement). La « mauvaise foi » est un projet conscient de se mentir : « Je me fais croire que je n'ai pas fumé. »
\item \textbf{Responsabilité :} Dire « C'est mon inconscient » est une excuse pour fuir l'angoisse de la liberté (choisir d'être commerçant anxieux, élève angoissée, père hypocrite).
\end{itemize} \\ \hline
\end{tabularx}
\end{center}

\textbf{Étape 4 — Rédaction type (Consigne 4) — \emph{Exemple de copie élève (15 lignes)}}

\begin{quote}
\emph{Thèse : Si l'inconscient freudien éclaire la structure profonde des actes manqués, la mauvaise foi sartrienne reste nécessaire pour rendre compte de la responsabilité du sujet face à ses symptômes.}

\emph{Argument 1 (Freud) : L'oubli du nom chez M. Kamga (Doc 1A) et le cauchemar récurrent de Ngo (Doc 1B) ne sont pas des hasards. Comme l'explique Freud (Doc 2), ils sont l'effet d'une cause psychique refoulée (hostilité, angoisse) qui s'impose au Moi par des mécanismes autonomes (déplacement, condensation), confirmés par la neurobiologie du sommeil paradoxal (Doc 4 : hyperactivité limbique sans contrôle cortical). Le sujet ne « choisit » pas son symptôme, il le subit.}

\emph{Argument 2 (Sartre) : Pourtant, M. Essomba (Doc 1C) \emph{sait} confusément qu'il a fumé (pâlement, silence). Sa sévérité n'est pas un automatisme : c'est un \emph{projet} de se construire une innocence morale en projetant sa faute sur son fils. C'est la définition même de la mauvaise foi (Doc 3) : se mentir à soi-même pour fuir la responsabilité de son passé.}

\emph{Synthèse : L'inconscient explique la \emph{genèse} déterminée du désir (le « pourquoi » causal), mais la mauvaise foi décrit la \emph{manière} dont le sujet, ici et maintenant, \emph{assume} ou \emph{refuse} ce désir (le « comment » finaliste).}

\emph{Ouverture : La psychanalyse ne libère-t-elle pas justement en transformant l'inconscient en conscience, c'est-à-dire en rendant au sujet la propriété de ses actes ( « Là où était le Ça, le Moi doit advenir ») ?}
\end{quote}
\end{exoresolu}

\courssub{Bilan de l'activité}
\begin{aretenir}[Ce que cette activité a permis de mettre en évidence]
\begin{enumerate}
    \item \textbf{Opérativité des concepts :} La distinction \cle{Conscient / Préconscient / Inconscient} (topique) et \cle{Ça / Moi / Surmoi} (dynamique) n'est pas abstraite : elle permet de \emph{découper} une situation concrète (oubli, rêve, violence éducative) en instances psychiques repérables.
    \item \textbf{Preuve par l'acte manqué :} L'oubli de nom (Récit A) est le paradigme freudien : un acte conscient (parler) est détourné par une intention inconsciente (refouler la rivalité). La \cle{résistance} (oubli) prouve l'existence du refoulé.
    \item \textbf{Le rêve comme « voie royale » :} Le Récit B + Doc 4 montrent que le sommeil n'est pas l'arrêt de la psyché mais un \emph{travail} (travail du rêve : condensation, déplacement) où l'inconscient s'exprime sous la censure du Moi. L'angoisse réveillée signale l'échec partiel de la censure.
    \item \textbf{Projection et formation réactionnelle :} Le Récit C illustre deux mécanismes de défense majeurs : M. Essomba \emph{projette} sa propre faute refoulée sur le fils (Ça $\rightarrow$ Autre) et utilise une \emph{formation réactionnelle} (sévérité excessive = inverse de la tentation) pour calmer son Surmoi.
    \item \textbf{Enjeu philosophique majeur :} Le débat Freud/Sartre ne porte pas sur l'existence des phénomènes (oubli, rêve, hypocrisie) mais sur leur \emph{statut ontologique} : \textbf{Determinisme psychique} (le sujet est divisé, l'inconscient est une cause) vs \textbf{Liberté radicale} (le sujet est un, l'inconscient est un leurre choisi). Le Probatoire attend que l'élève sache \emph{articuler} les deux sans les confondre.
    \item \textbf{Compétence rédactionnelle :} La contrainte des « 15 lignes » force à la \emph{structuration} (Thèse / Argument 1 / Argument 2 / Objection-Réponse / Ouverture) et à la \emph{précision} (citation doc, vocabulaire technique : refoulement, projection, mauvaise foi, principe de réalité).
\end{enumerate}
\end{aretenir}

\begin{attention}[Pièges à éviter pour le Probatoire]
\begin{itemize}
    \item Confondre \emph{Inconscient topique} (lieu) et \emph{Inconscient dynamique} (processus).
    \item Dire « L'inconscient, c'est ce dont on n'est pas conscient » (trop vague : inclut le préconscient, l'ignorance banale). \textbf{Précision :} L'inconscient freudien est \emph{refoulé}, \emph{dynamique}, \emph{résistant}.
    \item Attribuer à Sartre l'idée que « l'inconscient n'existe pas ». Sartre dit : « L'inconscient \emph{en tant qu'instance causale} n'existe pas ; ce qu'on appelle ainsi est un \emph{projet de mauvaise foi}. »
    \item Oublier de placer les éléments du récit \emph{sur le schéma} (consigne 2) : le schéma vide ne rapporte pas de points.
    \item Dépasser 15 lignes : sanction de hors-sujet / non-respect de consigne.
\end{itemize}
\end{attention}

\newpage

\coursec{Méthodes et exercices résolus}

\coursec{La conscience et l’inconscient}

\courssub{La conscience : notion et types}

\begin{definition}[Conscience]
La \cle{conscience} (du latin \emph{cum-scire}, « savoir avec ») est la connaissance immédiate, intuitive et non médiate que le sujet a de lui-même (ses pensées, ses sentiments, ses actes) et du monde extérieur. Elle se caractérise par deux propriétés fondamentales : l'\cle{intentionnalité} (toute conscience est conscience \emph{de} quelque chose, Brentano/Husserl) et la \cle{réflexivité} (capacité de se prendre soi-même pour objet).
\end{definition}

\begin{aretenir}[Les trois types de conscience]
\begin{itemize}
\item \textbf{Conscience spontanée (ou immédiate) :} présence au monde et à l'action sans retour sur soi. Ex. : je marche, je conduis « machinalement ». C'est la conscience \emph{vécue} (Bergson).
\item \textbf{Conscience réfléchie (ou réflexive) :} retour de l'esprit sur ses propres opérations. Je \emph{sais} que je pense, j'analyse mes motifs. Condition de la \cle{connaissance de soi} (« Connais-toi toi-même », Socrate) et de la liberté critique.
\item \textbf{Conscience morale :} faculté de juger intérieurement la valeur morale de ses actes (bien/mal), d'éprouver le devoir, la culpabilité ou le remords. Rousseau : « conscience, instinct divin » ; Kant : voix de la \cle{raison pratique}.
\end{itemize}
\textbf{Repère} : dans un sujet de dissertation, précisez toujours \emph{quel} type de conscience vous mobilisez (spontanée, réfléchie, morale). L'ambiguïté affaiblit l'argumentation.
\end{aretenir}

\courssub{Caractéristiques de la conscience}

\begin{propriete}[Caractéristiques structurales de la conscience]
\begin{enumerate}
\item \textbf{Intentionnalité :} pas de conscience « vide » ; elle vise toujours un objet (perception, pensée, désir). « Toute conscience est conscience de quelque chose » (Brentano, repris par Husserl et Sartre).
\item \textbf{Transparence (thèse cartésienne) :} pour Descartes (\emph{Méditations}), l'esprit se connaît lui-même avec une évidence absolue : « Je pense, donc je suis ». Rien de ce qui est en moi ne peut m'échapper (\emph{cogito}).
\item \textbf{Unité / Synthèse :} la conscience unifie la multiplicité des impressions en une expérience cohérente (Kant : « Je pense » doit pouvoir accompagner toutes mes représentations).
\item \textbf{Temporalité :} conscience de la durée (Bergson), rétention-présentation-protention (Husserl), projet vers l'avenir (Heidegger/Sartre).
\end{enumerate}
\end{propriete}

\courssub{Valeur de la conscience}

\begin{aretenir}[Valeur anthropologique et morale]
\begin{itemize}
\item \textbf{Fondement de la responsabilité :} on n'est responsable que de ce qu'on fait \emph{en conscience} (liberté éclairée). Droit camerounais : irresponsabilité pénale en cas de trouble psychique abolissant le discernement (Art. 64 Code Pénal).
\item \textbf{Condition de l'autonomie morale :} Kant lie conscience morale et \cle{devoir} (impératif catégorique) ; la conscience est la loi que la raison se donne à elle-même.
\item \textbf{Dignité humaine :} la conscience réflexive distingue l'homme de l'animal (Descartes : « l'âme est sans parties », animal-machine) et fonde les droits de l'homme (sujet de droit).
\item \textbf{Outil de connaissance :} la conscience réfléchie permet la science (méthode cartésienne : doute, analyse, synthèse) et la sagesse (examen de soi, stoïcisme).
\end{itemize}
\end{aretenir}

\courssub{Limites de la conscience}

\begin{aretenir}[Les limites intrinsèques et extrinsèques]
\begin{enumerate}
\item \textbf{Illusions et erreurs (Descartes, Malebranche) :} les sens trompent ; la mémoire est fallacieuse ; le préjugé déforme le jugement. La conscience n'est pas un critère de vérité infaillible.
\item \textbf{L'inconscient dynamique (Freud) :} désirs refoulés, pulsions, traumatismes déterminent nos choix à notre insu. « Le moi n'est pas maître dans sa propre maison. »
\item \textbf{L'inconscient cognitif (Leibniz, sciences cognitives) :} « petites perceptions » (Leibniz), traitements automatiques, priming, biais cognitifs (Kahneman) : la majeure partie du traitement mental est non-consciente.
\item \textbf{Conditionnement social et historique (Marx, Durkheim, Bourdieu) :} la conscience individuelle est traversée par des déterminismes extérieurs (idéologie, habitus, structures) qu'elle ne maîtrise pas.
\item \textbf{Relativité culturelle :} ce que la conscience morale juge « bien » varie selon les sociétés (anthropologie philosophique), remettant en cause son universalité naturelle.
\end{enumerate}
\end{aretenir}

\courssub{La découverte de l'inconscient et ses types}

\begin{savaistu}[Genèse philosophique de l'inconscient]
Avant Freud, l'inconscient est une \emph{idée philosophique} :
\begin{itemize}
\item \textbf{Leibniz (\emph{Nouveaux Essais}, 1704) :} les \cle{petites perceptions} — perceptions infiniment petites, confuses, dont nous n'avons pas conscience mais qui composent notre vie mentale (ex. : bruit de la mer perçu sans être aperçu).
\item \textbf{Schopenhauer (\emph{Le Monde comme volonté et comme représentation}, 1819) :} la \cle{Volonté} aveugle, irrationnelle, fond métaphysique de l'être, dont l'intellect n'est que le serviteur.
\item \textbf{Nietzsche (\emph{Généalogie de la morale}, 1887) :} les affects, instincts, « arrière-monde » physiologique qui commandent nos valeurs ; la conscience est un « outil » tardif et superficiel.
\item \textbf{Pierre Janet (psychologie française, fin XIX\textsuperscript{e}) :} \cle{désagrégation psychique}, subconscient pathologique (hystérie), automatismes psychologiques.
\end{itemize}
\textbf{Apport de Freud :} il transforme l'intuition philosophique en \cle{théorie scientifique} (topique, dynamique, économique) et en \cle{pratique thérapeutique} (psychanalyse : cure par la parole, association libre, interprétation des rêves).
\end{savaistu}

\begin{definition}[Types d'inconscient (repérage pour le Bac)]
\begin{itemize}
\item \textbf{Inconscient dynamique / freudien :} refoulé, pulsionnel, conflictuel (Ça vs Surmoi). Il \emph{veut} rester caché (résistance).
\item \textbf{Inconscient descriptif / topique :} tout ce qui n'est pas conscient à l'instant $t$ (préconscient = accessible ; inconscient proprement dit = refoulé).
\item \textbf{Inconscient cognitif / adaptatif (sciences cognitives modernes) :} processus automatiques, efficaces, non conflictuels (marche, lecture, grammaire, biais implicites). Pas de refoulement, pas de sens caché.
\item \textbf{Inconscient existentiel (Sartre) :} \textbf{rejeté}. Pour Sartre, la conscience est transparence ; l'inconscient est une \cle{mauvaise foi} pour fuir la liberté. Il n'y a que du \emph{non-thétique} (conscience non positionnelle de soi).
\end{itemize}
\end{definition}

\courssub{La psychanalyse freudienne et sa critique}

\begin{definition}[La première topique (1900) : Niveaux de l'appareil psychique]
\begin{itemize}
\item \textbf{Conscient (Cs) :} perceptions, pensées actuelles, surface.
\item \textbf{Préconscient (Pcs) :} souvenirs, savoirs accessibles par un effort d'attention (pas de refoulement). Réservoir de la mémoire disponible.
\item \textbf{Inconscient (Ics) :} contenus \cle{refoulés} (désirs interdits, traumatismes, pulsions sexuelles/agressives). Inaccessibles directement ; s'expriment par \cle{formations de l'inconscient} : rêves (voie royale), lapsus, actes manqués, symptômes névrotiques, humour.
\end{itemize}
\textbf{Mécanismes :} \cle{Refoulement} (opération originaire, « Vertreibung »), déplacement, condensation, symbolisation (travail du rêve).
\end{definition}

\begin{definition}[La seconde topique (1923) : Instances psychiques]
\begin{itemize}
\item \textbf{Le Ça (Es) :} pôle pulsionnel, inné, irrationnel, intemporel, régi par le \cle{principe de plaisir} (recherche satisfaction immédiate). « Ça » = l'inconnu en moi.
\item \textbf{Le Moi (Ich) :} instance consciente/préconsciente, médiateur, régi par le \cle{principe de réalité}. Adaptation au monde, défense (mécanismes de défense : refoulement, projection, rationalisation, sublimation).
\item \textbf{Le Surmoi (Über-Ich) :} intériorisation des interdits parentaux/sociaux (idéal du moi, conscience morale). Siège de la culpabilité, juge sévère. Héritier du complexe d'Œdipe résolu.
\end{itemize}
\textbf{Conflit :} Le Moi est le « pauvre moi » servi par trois maîtres exigeants : le monde extérieur, le Ça, le Surmoi. L'angoisse signale le danger.
\end{definition}

\begin{aretenir}[Critiques majeures de la psychanalyse (exigibles au Bac)]
\begin{enumerate}
\item \textbf{Critique épistémologique (Karl Popper, \emph{Conjectures et Réfutations}) :} la psychanalyse n'est pas \cle{réfutable} (falsifiable). Elle explique tout \emph{a posteriori} (si le patient va bien : cure réussie ; s'il va mal : résistance). Donc : \emph{pseudo-science} (métaphysique), non science au sens poppérien.
\item \textbf{Critique philosophique (Sartre, \emph{L'Être et le Néant}) :} l'inconscient freudien dissout la \cle{liberté} et la \cle{responsabilité}. L'homme est « condamné à être libre » ; invoquer l'inconscient, c'est de la \cle{mauvaise foi} pour fuir l'angoisse du choix. La conscience est transparence (non-thétique).
\item \textbf{Critique herméneutique (Paul Ricœur, \emph{De l'interprétation}) :} la psychanalyse n'est pas une science causale mais une \cle{sémantique} : elle \emph{interprète} des symboles (comme l'archéologie ou l'exégèse). Elle vise la \emph{compréhension} (Verstehen) du sens, pas l'explication (Erklären) par des lois naturelles.
\item \textbf{Critique féministe / culturelle :} universalisation du modèle viennois bourgeois (complexe d'Œdipe phallocentrique), négligence des déterminismes sociaux (race, classe, genre).
\end{enumerate}
\textbf{Nuance} : malgré les critiques, Freud a opéré une « blessure narcissique » décisive : l'homme n'est pas transparent à lui-même. La psychanalyse reste une pratique clinique efficace pour beaucoup (thérapie par la parole, élargissement du Moi : « Wo Es war, soll Ich werden »).
\end{aretenir}

\courssub{Synthèse : conscience et inconscient, complémentaires ou adversaires ?}

\begin{exemplebox}[Synthèse dialectique pour la dissertation]
\begin{center}
\begin{tabularx}{\linewidth}{|l|X|X|}\hline
\textbf{Dimension} & \textbf{Thèse : Souveraineté de la conscience} & \textbf{Antithèse : Hégémonie de l'inconscient} \\ \hline
\textbf{Anthropologie} & Homme = sujet rationnel, maître de ses actes (Descartes, Kant). & Homme = être divisé, déterminé par des forces obscures (Freud, Nietzsche). \\ \hline
\textbf{Moralité} & Responsabilité pleine ; conscience = juge intérieur. & Responsabilité atténuée ; l'acte a des motifs inconnus du sujet. \\ \hline
\textbf{Connaissance de soi} & Possible par la réflexion (introspection, examen de conscience). & Illusoire sans tiers (psychanalyste) ; résistance au savoir. \\ \hline
\textbf{Synthèse (Ricœur, Merleau-Ponty, éthique moderne)} & \textbf{Conscience et inconscient sont co-originaires.} La conscience émerge \emph{sur} un fond inconscient (corps, histoire, langage). La maîtrise de soi n'est pas un \emph{don} (transparence cartésienne) mais une \emph{conquête} (lucidité, travail sur soi, thérapie, éducation). L'inconscient n'est pas seulement l'ennemi : il est aussi réserve de créativité (sublimation), d'intuition, de vie. \\ \hline
\end{tabularx}
\end{center}
\end{exemplebox}

% ============================================================
% PARTIE 2 : SAVOIR-FAIRE (MÉTHODES ET ENTRAÎNEMENTS BAC)
% ============================================================

\courssub{Savoir-faire 1 : Appliquer les notions de conscience et d'inconscient à une situation concrète}

\begin{methode}[Analyser une situation de la vie courante avec les notions de la leçon]
Pour appliquer les notions de \cle{conscience} et d'\cle{inconscient} à une situation concrète (un acte manqué, un mensonge, un rêve, une décision, un symptôme), suivre la démarche suivante :
\begin{enumerate}
\item \textbf{Lire attentivement la situation} et dégager les faits observables (paroles, gestes, comportements, réactions d'autrui).
\item \textbf{Identifier la notion pertinente} : s'agit-il d'un acte volontaire et réfléchi (conscience réfléchie/morale) ou d'un acte involontaire, inexpliqué par le sujet lui-même (inconscient dynamique / formation de l'inconscient) ?
\item \textbf{Mobiliser les caractéristiques de la conscience} : intentionnalité, réflexivité, unité, temporalité, transparence (ou opacité).
\item \textbf{Mobiliser la théorie freudienne} si nécessaire : distinguer \cle{conscient}, \cle{préconscient} et \cle{inconscient} (1\textsuperscript{re} topique) ; repérer les instances (Ça, Moi, Surmoi) et les formations de l'inconscient (rêve, lapsus, acte manqué, symptôme, humour).
\item \textbf{Interpréter la situation} en reliant explicitement le fait observé à la notion : « ce lapsus manifeste un désir refoulé car... » ; « cette hésitation révèle un conflit Moi/Ça... ».
\item \textbf{Conclure} en formulant un jugement nuancé : la conscience suffit-elle à expliquer le comportement, ou faut-il admettre l'inconscient comme instance déterminante ?
\end{enumerate}
\textbf{Réflexe Bac} : toujours citer un élément précis de la situation (un mot, un geste, un silence) pour justifier l'usage de la notion. Une application sans ancrage dans le cas concret est hors sujet.
\end{methode}

\begin{exoresolu}[Le lapsus du proviseur : analyse freudienne]
\textbf{Énoncé.} Pendant la cérémonie de remise des prix du lycée, le proviseur déclare : « Je félicite tous les élèves qui ont \emph{échoué}... pardon, \emph{réussi} leur examen. » Les élèves éclatent de rire. À l'aide de la théorie freudienne de l'inconscient, analysez cette situation et montrez ce qu'elle révèle du fonctionnement psychique.
\tcblower
\textbf{Solution détaillée.}
\begin{enumerate}
\item \textbf{Faits observables.} Le proviseur prononce un mot (« échoué ») contraire à son intention consciente (« réussi »). Il s'agit d'un \cle{lapsus} (erreur de langage involontaire). Le sujet se corrige aussitôt (« pardon »), preuve qu'il n'a pas voulu dire ce mot. Le rire des élèves signale la perception d'une vérité cachée.
\item \textbf{Notion pertinente.} L'acte est involontaire, non voulu par le \cle{Moi} (instance consciente, principe de réalité). La thèse cartésienne de la transparence de l'esprit à lui-même (« je pense, donc je suis » implique maîtrise) est mise en défaut. Il faut faire appel à l'\cle{inconscient} (1\textsuperscript{re} topique) et au \cle{Ça} (2\textsuperscript{e} topique).
\item \textbf{Analyse freudienne.} Selon Freud (\emph{Psychopathologie de la vie quotidienne}, 1901), le lapsus n'est pas un hasard ni une simple fatigue : c'est une \cle{formation de l'inconscient}. Un contenu refoulé (désir, pensée, hostilité, angoisse) franchit la \cle{censure} (barrière Moi/Inconscient) et se manifeste déguisé dans le langage. Le mot prononcé trahit ce que le sujet pense ou souhaite \emph{réellement} au niveau du Ça.
\item \textbf{Interprétation contextuelle (Cameroun).} Le lapsus du proviseur peut révéler : (a) une angoisse refoulée concernant le taux de réussite réel de l'établissement (pression hiérarchique) ; (b) une pensée critique inconsciente envers une partie des élèves ; (c) un désir masochiste de saboter son propre discours. Le rire des élèves agit comme une \cle{validation sociale} de la vérité de l'inconscient : le mot « échappé » dit tout haut ce que la politesse officielle (Surmoi) taisait.
\item \textbf{Conclusion.} Cette situation illustre la célèbre formule : « Le moi n'est pas maître dans sa propre maison » (Freud, \emph{Introduction à la psychanalyse}, 1917). La conscience ne suffit pas à expliquer nos conduites ; le lapsus est une voie d'accès privilégiée à l'inconscient, qui agit en nous à notre insu.
\end{enumerate}
\textbf{Phrase de conclusion type.} Le lapsus du proviseur confirme la thèse freudienne : nos erreurs involontaires de la vie quotidienne sont des messages codés de l'inconscient, et l'analyse de la vie quotidienne valide l'existence d'une vie psychique qui échappe au contrôle volontaire.
\end{exoresolu}

\courssub{Savoir-faire 2 : Rédiger une introduction de dissertation problématisée}

\begin{methode}[Construire une introduction en quatre étapes (Standard Bac)]
Une introduction de dissertation philosophique au Probatoire/Baccalauréat technique comporte quatre moments obligatoires :
\begin{enumerate}
\item \textbf{L'accroche (amorce) :} un fait, une expérience commune, une citation courte ou une réflexion générale qui met en relation avec le sujet. \textit{Exemple :} « Nous croyons tous savoir ce que nous pensons et pourquoi nous agissons. »
\item \textbf{La présentation du sujet et la définition des termes :} définir brièvement et rigoureusement les notions clés du sujet (ex. : conscience = connaissance immédiate/réflexive ; inconscient = ensemble des processus psychiques échappant à la conscience, refoulés).
\item \textbf{La problématique :} formuler le problème sous forme d'une ou deux questions interrogatives qui opposent deux thèses possibles (thèse/antithèse). \textit{Exemple :} « La conscience suffit-elle à rendre compte de notre vie psychique ? L'homme est-il maître de lui-même, ou est-il déterminé par des forces inconscientes ? »
\item \textbf{L'annonce du plan :} présenter les deux ou trois grandes parties du devoir (thèse, antithèse, synthèse) en une ou deux phrases, sans développer les arguments.
\end{enumerate}
\textbf{Réflexes Bac} : une introduction fait environ 8 à 12 lignes manuscrites ; elle ne contient \textbf{ni} exemple détaillé, \textbf{ni} citation longue, \textbf{ni} opinion personnelle tranchée (« je pense que... » est à réserver à la conclusion). Le ton est objectif et analytique.
\end{methode}

\begin{exoresolu}[Sujet : « Suis-je maître de moi-même ? »]
\textbf{Énoncé.} Rédigez l'introduction complète d'une dissertation sur le sujet : \emph{Suis-je maître de moi-même ?}
\tcblower
\textbf{Solution détaillée.}

\textbf{Étape 1 — Accroche.} Partir d'une évidence commune : chacun se sent libre de ses pensées et de ses actes. « Il semble aller de soi que je sais ce que je pense, ce que je veux et pourquoi j'agis : je me sens spontanément maître de moi-même. »

\textbf{Étape 2 — Définitions.} « La \cle{conscience} désigne la connaissance immédiate que le sujet a de ses états intérieurs et du monde ; elle se double d'une conscience réflexive, par laquelle je me sais pensant et agissant. Mais la psychologie moderne, avec Freud, a mis au jour l'\cle{inconscient} : un ensemble de désirs refoulés, de pulsions et de processus psychiques qui déterminent nos conduites à notre insu. »

\textbf{Étape 3 — Problématique.} « Si une partie de ma vie psychique m'échappe, puis-je encore prétendre me connaître et me gouverner ? La conscience est-elle la souveraine de nos conduites, ou n'est-elle que la partie émergée d'une vie psychique dominée par l'inconscient ? Autrement dit : \emph{suis-je véritablement maître de moi-même ?} »

\textbf{Étape 4 — Annonce du plan.} « Nous verrons d'abord que la conscience, par la réflexion et la volonté, fonde notre maîtrise de nous-mêmes (thèse cartésienne et kantienne) ; nous montrerons ensuite que la découverte de l'inconscient ébranle radicalement cette prétention (thèse freudienne) ; nous demanderons enfin si conscience et inconscient ne peuvent pas se concilier dans une maîtrise relative, conquise par le travail de lucidité (synthèse : Ricœur, Sartre, éthique du care). »

\textbf{Texte rédigé (modèle complet pour copie).}

Il semble aller de soi que je sais ce que je pense, ce que je veux et pourquoi j'agis : chacun se sent spontanément maître de lui-même. La \cle{conscience} désigne en effet la connaissance immédiate que le sujet a de ses états intérieurs et du monde ; elle se double d'une conscience réflexive, par laquelle je me sais pensant et agissant. Mais la psychologie moderne, avec Freud, a mis au jour l'\cle{inconscient} : un ensemble de désirs refoulés et de processus psychiques qui déterminent nos conduites à notre insu. Dès lors, un problème se pose : si une partie de ma vie psychique m'échappe, puis-je encore prétendre me connaître et me gouverner ? La conscience est-elle souveraine, ou n'est-elle que la partie émergée d'une vie psychique dominée par l'inconscient ? Suis-je véritablement maître de moi-même ? Nous verrons d'abord que la conscience fonde notre maîtrise de nous-mêmes ; nous montrerons ensuite que la découverte de l'inconscient ébranle cette prétention ; nous demanderons enfin si une maîtrise relative, mais réelle, demeure possible.

\textbf{Justification de la démarche.} Ce texte respecte les quatre étapes : l'accroche part du sens commun ; les deux notions du sujet sont définies ; la problématique oppose deux thèses (souveraineté de la conscience / détermination inconsciente) sous forme interrogative ; le plan dialectique est annoncé en trois temps.
\end{exoresolu}

\courssub{Savoir-faire 3 : Rédiger un paragraphe argumenté (Thèse, Argument, Exemple, Auteur)}

\begin{methode}[La structure d'un paragraphe de dissertation (Modèle T.D.A.A.T.)]
Chaque paragraphe du développement suit une structure rigoureuse en cinq temps :
\begin{enumerate}
\item \textbf{Thèse (Idée directrice) :} une phrase qui annonce l'argument principal du paragraphe.
\item \textbf{Définition / Explicitation :} préciser la notion mobilisée (type de conscience, concept d'inconscient).
\item \textbf{Argument (Raisonnement) :} un enchaînement logique qui justifie la thèse (connecteurs : \emph{car, en effet, dès lors, par conséquent, or, cependant}).
\item \textbf{Auteur / Référence philosophique :} citer brièvement un auteur et sa thèse précise (Descartes, Freud, Sartre, Nietzsche, Kant, Ricœur...), sans longue citation. Le nom de l'auteur vaut argument d'autorité.
\item \textbf{Exemple concret + Transition :} illustrer par un cas précis, de préférence tiré de la vie courante ou du contexte camerounais, puis conclure le paragraphe par une phrase qui ouvre sur l'argument suivant.
\end{enumerate}
\textbf{Formule à retenir} : \textbf{T}hèse + \textbf{D}éfinition + \textbf{A}rgument + \textbf{A}uteur + \textbf{E}xemple + \textbf{T}ransition. Un paragraphe sans argument est une affirmation gratuite ; un paragraphe sans exemple reste abstrait ; un paragraphe sans auteur manque de profondeur philosophique.
\end{methode}

\begin{exoresolu}[Rédiger un paragraphe sur la valeur de la conscience morale]
\textbf{Énoncé.} Rédigez un paragraphe argumenté montrant que la conscience est une valeur fondamentale pour la vie morale et la responsabilité.
\tcblower
\textbf{Solution détaillée.}

\textbf{Construction pas à pas (T.D.A.A.T.).}
\begin{enumerate}
\item \textbf{Thèse} : La conscience morale est le fondement indispensable de la responsabilité humaine.
\item \textbf{Définition} : Par \cle{conscience morale}, on entend la faculté par laquelle le sujet juge intérieurement le bien et le mal, s'approuve ou se condamne lui-même (remords, culpabilité).
\item \textbf{Argument} : La responsabilité suppose l'imputabilité : on ne peut reprocher à un homme que ce qu'il a voulu \emph{en connaissance de cause}. Un acte accompli sans conscience (somnambulisme, folie, contrainte physique) n'est pas imputable, car l'agent n'en était pas l'auteur conscient.
\item \textbf{Auteur} : Kant (\emph{Fondements de la métaphysique des mœurs}) fait de la conscience morale la voix de la \cle{raison pratique} qui ordonne le devoir (impératif catégorique) ; Rousseau (\emph{Émile}, « Profession de foi du vicaire savoyard ») la décrit comme un « \cle{instinct divin} », juge infaillible du bien et du mal, source de la dignité.
\item \textbf{Exemple + Transition} : Ainsi, dans les tribunaux camerounais, un accusé reconnu atteint de troubles psychiques abolissant le discernement au moment des faits est déclaré \cle{pénalement irresponsable} (Art. 64 Code Pénal) : le droit positif lui-même lie responsabilité et conscience. La conscience apparaît donc comme le fondement de la dignité morale de l'homme ; \textbf{mais cette souveraineté est-elle sans limites ?}
\end{enumerate}

\textbf{Paragraphe rédigé (modèle pour copie).}

La conscience morale possède une valeur fondamentale, car elle est la condition même de la vie morale et de la responsabilité. Par \cle{conscience morale}, on entend la faculté par laquelle le sujet juge intérieurement le bien et le mal, et s'approuve ou se condamne lui-même. Or la responsabilité morale et juridique suppose la conscience : on ne peut reprocher à un homme que ce qu'il a voulu et compris ; un acte commis dans l'ignorance, la folie ou le somnambulisme n'est pas imputable, car l'agent n'en était pas l'auteur conscient. C'est pourquoi Kant fait de la conscience morale la voix de la raison pratique qui ordonne le devoir, et Rousseau y voit un « instinct divin », juge infaillible du bien et du mal. Ainsi, dans les tribunaux camerounais, un accusé reconnu atteint de troubles psychiques au moment des faits est déclaré pénalement irresponsable : le droit lui-même reconnaît que sans conscience, il n'y a ni culpabilité ni mérite. La conscience apparaît donc comme le fondement de la dignité morale de l'homme ; mais cette souveraineté est-elle sans limites ?

\textbf{Justification.} Le paragraphe suit les cinq étapes (T.D.A.A.T.), mobilise deux auteurs classiques, un raisonnement déductif clair (responsabilité $\rightarrow$ conscience) et un exemple juridique concret et localisé (Cameroun), et s'achève par une transition explicite vers les limites de la conscience (antithèse).
\end{exoresolu}

\courssub{Savoir-faire 4 : Présenter et discuter la théorie freudienne (Explication de texte / Dissertation)}

\begin{methode}[Exposer puis critiquer la psychanalyse freudienne (Standard Bac)]
Pour traiter la \cle{psychanalyse freudienne} dans un devoir (explication de texte ou dissertation), respecter cet ordre :
\begin{enumerate}
\item \textbf{Situer historiquement :} Freud, médecin viennois, fonde la psychanalyse à la fin du XIX\textsuperscript{e} siècle (\emph{L'Interprétation des rêves}, 1900 ; \emph{Psychopathologie de la vie quotidienne}, 1901 ; \emph{Introduction à la psychanalyse}, 1917).
\item \textbf{Exposer la première topique (1900) :} distinguer les trois \cle{niveaux} : \cle{Conscient} (surface), \cle{Préconscient} (mémoire disponible), \cle{Inconscient} (refoulé, pulsions). Mécanisme clé : le \cle{refoulement}. Formations de l'inconscient : rêves (voie royale), lapsus, actes manqués, symptômes, humour.
\item \textbf{Exposer la seconde topique (1923) :} distinguer les trois \cle{instances} : le \cle{Ça} (pulsions, principe de plaisir, inconscient), le \cle{Moi} (médiateur, principe de réalité, conscient/préconscient), le \cle{Surmoi} (interdits moraux intériorisés, idéal du moi, conscience morale). Conflit permanent $\rightarrow$ angoisse $\rightarrow$ mécanismes de défense.
\item \textbf{Présenter les critiques majeures (exigibles) :}
      \begin{itemize}
      \item \textbf{Karl Popper (épistémologie) :} non-réfutabilité = pseudo-science. Toute objection = « résistance ».
      \item \textbf{Jean-Paul Sartre (ontologie/éthique) :} l'inconscient = \cle{mauvaise foi}. La conscience est transparence (non-thétique). L'homme est liberté absolue.
      \item \textbf{Paul Ricœur (herméneutique) :} psychanalyse = \emph{herméneutique du soupçon} (avec Marx, Nietzsche). Elle \emph{interprète} des symboles, elle ne \emph{prouve} pas des causes. « Archéologie du sujet ».
      \item \textbf{Critiques féministes/sociales :} universalisation du modèle occidental bourgeois, phallocentrisme, oubli des déterminismes sociaux (Bourdieu, Fanon).
      \end{itemize}
\item \textbf{Nuancer / Héritage :} même critiquée, la psychanalyse a montré que la conscience n'est pas transparente ; elle a transformé la psychiatrie (thérapie par la parole), la littérature, l'anthropologie. « \emph{Wo Es war, soll Ich werden} » (Là où était le Ça, le Moi doit advenir) : la cure vise à élargir la conscience.
\end{enumerate}
\textbf{Réflexe Bac} : ne jamais présenter Freud comme ayant « découvert » l'inconscient \emph{ex nihilo}. Rappeler l'intuition philosophique antérieure (Leibniz, Schopenhauer, Nietzsche, Janet). Freud lui a donné un statut \emph{clinique}, une \emph{topique} et une \emph{méthode} (cure).
\end{methode}

\begin{exoresolu}[Explication de texte courte : « Le moi n'est pas maître dans sa propre maison »]
\textbf{Énoncé.} Freud écrit : « Le moi n'est pas maître dans sa propre maison. » (\emph{Introduction à la psychanalyse}, 1917, 17\textsuperscript{e} conférence). Expliquez cette formule, dégagez sa signification philosophique et discutez-la.
\tcblower
\textbf{Solution détaillée.}

\textbf{1. Explication littérale de la formule.}
La « maison » désigne la vie psychique totale de l'individu ; le « moi » désigne ici l'instance consciente (le Moi de la 2\textsuperscript{e} topique, ou le Conscient de la 1\textsuperscript{re}). Freud affirme que la conscience ne règne pas sur l'ensemble du psychisme : la plus grande partie de notre vie mentale — désirs refoulés, souvenirs traumatiques oubliés, pulsions sexuelles et agressives (le \cle{Ça}) — se déroule dans l'\cle{inconscient}, hors du contrôle du Moi. Le sujet agit souvent sans connaître les véritables motifs de ses actes ; il se croit libre alors qu'il est déterminé.

\textbf{2. Signification philosophique (Portée anthropologique).}
Cette formule renverse la tradition rationaliste et humaniste. Pour \textbf{Descartes}, la conscience est transparence absolue : « Je pense, donc je suis » signifie que rien de ce qui est en moi ne peut m'échapper ; le moi est roi. Freud inflige à l'humanité une troisième « \cle{blessure narcissique} » (après Copernic — la Terre n'est pas le centre — et Darwin — l'homme descend de l'animal) : \textbf{l'homme n'est même pas maître chez lui}. L'inconscient se manifeste par des voies détournées (formations de l'inconscient) : rêves, lapsus, actes manqués, symptômes névrotiques, humour. La raison n'est pas la souveraine ; elle est le serviteur de pulsions qu'elle ignore.

\textbf{3. Discussion critique (Arguments pour/contre).}
\begin{itemize}
\item \textbf{Force explicative (Pour) :} Cette thèse rend compte de phénomènes que la seule conscience ne saurait expliquer : pourquoi un élève brillant oublie-t-il soudain la date de son examen ? (Acte manqué = désir refoulé d'éviter l'épreuve). Pourquoi un homme aimant blesse-t-il sa femme par un lapsus ? (Hostilité refoulée). La psychanalyse donne du sens à l'absurde apparent.
\item \textbf{Critique de la scientificité (Popper) :} La théorie n'est pas \cle{réfutable}. Si le patient va bien : confirmation. S'il va mal : « résistance » = confirmation. Elle explique tout \emph{a posteriori}, donc ne prédit rien. Statut de \emph{métaphysique} ou de \emph{mythe} heuristique, pas de science au sens poppérien.
\item \textbf{Critique de la liberté (Sartre) :} L'inconscient freudien (Ça) fonctionne comme un \cle{déterminisme} interne qui dépossède le sujet. Or pour Sartre, « l'homme est condamné à être libre » : il choisit toujours, même de fuir sa liberté (mauvaise foi). Invoquer l'inconscient, c'est se dédouaner.
\item \textbf{Critique herméneutique (Ricœur) :} Freud ne fait pas de la \emph{physique} du psychisme, mais de la \emph{sémantique}. Il déchiffre des symboles. La cure est une \emph{interprétation} qui libère le sens, pas une chirurgie qui retire une cause.
\end{itemize}

\textbf{4. Conclusion synthétique.}
La formule de Freud a le mérite historique et philosophique de briser l'illusion d'une conscience toute-puissante et transparente : nous ne sommes pas totalement transparents à nous-mêmes. Mais reconnaître l'inconscient ne signifie pas renoncer à toute maîtrise. La cure psychanalytique elle-même vise à \emph{rendre conscient} ce qui était refoulé (« Là où était le Ça, le Moi doit advenir »), donc à élargir le règne du Moi. La maîtrise de soi n'est pas un \emph{donné} ontologique (cartésien), mais une \emph{tâche} éthique et un \emph{travail} jamais achevé (Ricœur, éthique moderne).
\end{exoresolu}

\courssub{Savoir-faire 5 : Rédiger une conclusion et justifier une position}

\begin{methode}[Construire une conclusion de dissertation (Standard Bac)]
La conclusion est le moment où l'élève \textbf{justifie une réponse personnelle argumentée} à la problématique. Démarche en quatre temps :
\begin{enumerate}
\item \textbf{Bilan (Synthèse) :} rappeler en deux phrases maximum les deux thèses confrontées (ex. : valeur de la conscience / défi de l'inconscient). Ne pas répéter les arguments, juste les grandes lignes.
\item \textbf{Réponse à la problématique :} reprendre la question initiale et y répondre clairement, de façon nuancée (« La conscience ne suffit donc pas à... mais elle n'est pas non plus... » ; « La maîtrise n'est pas absolue, elle est relative... »).
\item \textbf{Justification (Le « pourquoi » de la réponse) :} donner la raison philosophique principale du choix (ex. : la synthèse : conscience et inconscient sont co-originaires ; la lucidité est un progrès, jamais un état acquis ; la responsabilité se construit).
\item \textbf{Ouverture :} élargir vers une question voisine (la liberté, la responsabilité, la connaissance de soi, l'éthique, l'intelligence artificielle) \textbf{sans} traiter ce nouveau sujet. Une phrase suffit.
\end{enumerate}
\textbf{Attention Bac} : la conclusion ne doit introduire \textbf{ni} nouvel argument, \textbf{ni} nouvel auteur, \textbf{ni} nouvel exemple. Elle mesure environ 6 à 8 lignes manuscrites. Elle doit montrer que le devoir a \emph{progressé}.
\end{methode}

\begin{exoresolu}[Conclure le sujet « Suis-je maître de moi-même ? »]
\textbf{Énoncé.} Rédigez la conclusion de la dissertation : \emph{Suis-je maître de moi-même ?} (Le développement a montré : I. La conscience fonde la maîtrise ; II. L'inconscient la limite radicalement ; III. Une maîtrise relative par le travail de lucidité).
\tcblower
\textbf{Solution détaillée.}

\textbf{Analyse de la tâche.} Le devoir a opposé deux thèses : la conscience comme fondement de la maîtrise de soi (Descartes, Kant) et l'inconscient comme remise en cause de cette maîtrise (Freud, Nietzsche). La conclusion doit dépasser cette opposition par une synthèse justifiée.

\textbf{Texte rédigé (modèle pour copie).}

Notre parcours a montré que la conscience, par la réflexion et la volonté, fonde la responsabilité morale et la maîtrise de soi ; mais il a aussi établi, avec Freud, qu'une part irréductible de notre vie psychique — désirs refoulés, pulsions, souvenirs oubliés — détermine nos conduites à notre insu. Faut-il alors renoncer à toute maîtrise ? Non : l'inconscient ne supprime pas la conscience, il en fixe les limites structurelles. Je ne suis pas maître de moi-même de façon absolue et spontanée, comme le voulait Descartes, mais je peux le devenir davantage : la connaissance de soi, l'examen de conscience, l'éducation, et même la cure psychanalytique, élargissent le domaine de la lucidité. La maîtrise de soi n'est donc pas un \emph{état} naturel, mais une \emph{conquête} éthique, un travail permanent sur soi. Reste alors une question décisive : si je ne suis pas pleinement maître de mes motifs profonds, dans quelle mesure puis-je être tenu pour pleinement responsable de mes actes ? C'est le problème de la \cle{liberté} et de son articulation avec la \cle{responsabilité}.

\textbf{Justification de la démarche.} Le texte respecte les quatre étapes : 1. Bilan des deux thèses (phrase 1) ; 2. Réponse nuancée à la problématique (« Non... mais je peux le devenir ») ; 3. Justification par l'idée de lucidité comme conquête/travail (synthèse Ricœur/Sartre/Freud) ; 4. Ouverture vers le problème de la liberté et de la responsabilité, sans le traiter.
\end{exoresolu}

\courssub{Savoir-faire 6 : Distinguer les types de conscience dans un raisonnement (Exercice d'application)}

\begin{methode}[Utiliser correctement la typologie de la conscience (Anti-confusion)]
Pour éviter les confusions qui coûtent des points dans une copie de dissertation ou d'explication :
\begin{enumerate}
\item \textbf{Conscience spontanée (ou immédiate) :} simple présence à soi et au monde, sans retour sur soi. \textit{Mots-clés :} « je vis », « j'agis », « je perçois », « sans y penser ». Exemple : je conduis ma moto en discutant.
\item \textbf{Conscience réfléchie (réflexive) :} retour de l'esprit sur lui-même ; je me sais pensant/agissant. \textit{Mots-clés :} « je me demande », « je me rends compte », « j'analyse », « pourquoi ? ». Condition de la \cle{connaissance de soi} (Socrate) et de la \cle{liberté critique}.
\item \textbf{Conscience morale :} jugement intérieur du bien et du mal, sentiment du devoir, culpabilité, remords. \textit{Mots-clés :} « j'aurais dû », « c'est mal », « je me reproche », « conscience tranquille ». Fondement de la \cle{responsabilité} (Kant, Rousseau).
\item \textbf{Procédé d'identification systématique :} face à un exemple, se demander : \textbf{1.} Le sujet s'observe-t-il penser/agir ? $\rightarrow$ Réfléchie. \textbf{2.} Le sujet juge-t-il la valeur morale de son acte ? $\rightarrow$ Morale. \textbf{3.} Le sujet agit-il simplement sans se poser la question ? $\rightarrow$ Spontanée.
\end{enumerate}
\textbf{Réflexe Bac} : dans votre copie, citez le type de conscience \emph{avec précision} (ex. : « Ici, c'est la conscience \emph{réfléchie} qui opère... »). Écrire seulement « la conscience » sans préciser affaiblit l'argumentation et montre une maîtrise approximative du vocabulaire technique.
\end{methode}

\begin{exoresolu}[Identifier les types de conscience : Le benskin à Douala]
\textbf{Énoncé.} Un chauffeur de taxi-moto (benskin) à Douala roule en discutant avec un passager, sans penser à sa conduite. Arrivé à un carrefour, il hésite : « Pourquoi ai-je accéléré tout à l'heure ? ». Plus tard, ayant frôlé un piéton, il se reproche : « J'aurais dû ralentir, j'ai mal agi. » Identifiez et justifiez les trois types de conscience présents dans ce récit.
\tcblower
\textbf{Solution détaillée.}
\begin{enumerate}
\item \textbf{Premier moment : Conscience spontanée.} Le chauffeur roule « sans penser à sa conduite » : il est présent au monde (route, circulation) et à ses actes (accélérer, freiner), mais sans retour sur lui-même. C'est la \cle{conscience spontanée} (ou immédiate), celle de l'action ordinaire, fluide, non réfléchie (Bergson : « l'action qui se suffit à elle-même »).
\item \textbf{Deuxième moment : Conscience réfléchie.} « Pourquoi ai-je accéléré ? » : le chauffeur prend ses propres actes et motivations pour objet de pensée. Il opère un \cle{retour réflexif} de l'esprit sur lui-même : c'est la \cle{conscience réfléchie}, qui rend possible la \cle{connaissance de soi} (« Connais-toi toi-même », Socrate) et l'analyse des causes de son comportement.
\item \textbf{Troisième moment : Conscience morale.} « J'aurais dû ralentir, j'ai mal agi » : le chauffeur juge son acte du point de vue de la valeur (bien/mal), éprouve un \cle{remords} et formule un \cle{devoir} (« j'aurais dû »). C'est la \cle{conscience morale}, juge intérieur de nos conduites (Rousseau : « instinct divin » ; Kant : voix de la raison pratique).
\end{enumerate}
\textbf{Conclusion.} Ce récit montre que la conscience n'est pas univoque : elle est d'abord présence immédiate au monde (spontanée), puis retour réflexif sur soi (réfléchie), enfin jugement éthique (morale). La vie humaine passe constamment de l'un à l'autre de ces registres ; la philosophie doit les distinguer pour ne pas les confondre.
\end{exoresolu}

\begin{attention}[Les 5 erreurs qui coûtent des points au Probatoire / Baccalauréat]
\begin{enumerate}
\item \textbf{Confondre inconscient et non-conscient / préconscient.} L'inconscient freudien n'est pas le simple fait de « ne pas penser à quelque chose » (cela relève du \cle{préconscient} : accessible) ni de l'inattention : c'est un ensemble de contenus \emph{refoulés} (interdits, pulsions) qui agissent \emph{activement} sur le sujet. Ne jamais écrire « inconscient » pour désigner un état de coma, de sommeil profond ou de distraction.
\item \textbf{Présenter Freud comme l'inventeur absolu de l'inconscient.} Oublier que Leibniz (petites perceptions), Schopenhauer (Volonté), Nietzsche (instincts) ou Janet (subconscient) en avaient eu l'intuition philosophique ou clinique avant lui est une imprécision historique sanctionnée. Freud a \emph{systématisé} l'inconscient, lui a donné une \emph{topique}, une \emph{dynamique} (pulsions) et une \emph{méthode thérapeutique} (psychanalyse).
\item \textbf{Affirmer sans argumenter (Dogmatisme).} Écrire « la conscience a des limites » ou « l'inconscient existe » sans expliquer \emph{pourquoi} (illusions des sens, erreurs de la mémoire, déterminismes inconscients, refoulement) ni donner d'exemple concret : une thèse non argumentée n'est pas un paragraphe philosophique, c'est une opinion.
\item \textbf{Oublier la problématique interrogative dans l'introduction.} Une introduction qui définit les termes mais ne pose \textbf{aucune question} (point d'interrogation) ne peut pas annoncer de plan dialectique : la dissertation philosophique part \emph{toujours} d'un problème formulé par une question qui oppose deux thèses.
\item \textbf{Introduire un nouvel auteur, un nouvel argument ou un nouvel exemple dans la conclusion.} La conclusion fait le \emph{bilan}, \emph{répond} à la problématique et \emph{ouvre} ; toute nouveauté argumentative y est hors de propos (hors structure) et révèle un devoir mal construit (plan non tenu).
\end{enumerate}
\end{attention}

\newpage

\coursec{Exercices}

\courssub{Je vérifie mes connaissances}

\begin{exercice}[QCM : notions fondamentales]
Pour chaque question, une seule réponse est exacte. Recopiez le numéro et la lettre choisie.
\begin{enumerate}
    \item La conscience \emph{morale} se distingue de la conscience \emph{psychologique} par le fait qu'elle :
    \begin{itemize}
        \item[A.] Permet de percevoir le monde extérieur.
        \item[B.] Juge la valeur éthique de nos actes (bien/mal).
        \item[C.] Est l'apanage des philosophes.
        \item[D.] Correspond au \emph{Ça} freudien.
    \end{itemize}
    \item Dans la \emph{seconde topique} de Freud, l'instance qui représente les interdits intériorisés, l'idéal moral et la conscience morale est :
    \begin{itemize}
        \item[A.] Le Ça (\emph{Es}).
        \item[B.] Le Moi (\emph{Ich}).
        \item[C.] Le Surmoi (\emph{Über-Ich}).
        \item[D.] Le Préconscient.
    \end{itemize}
    \item Le \emph{lapsus} (acte manqué) est interprété par la psychanalyse comme :
    \begin{itemize}
        \item[A.] Une simple erreur de distraction sans signification.
        \item[B.] La manifestation directe du Surmoi.
        \item[C.] Un compromis entre une intention consciente et un désir refoulé (inconscient).
        \item[D.] Une preuve que le Moi est maître dans sa propre maison.
    \end{itemize}
    \item Quel philosophe a formulé la critique célèbre : « Le Moi n'est pas maître dans sa propre maison » ?
    \begin{itemize}
        \item[A.] René Descartes.
        \item[B.] Jean-Paul Sartre.
        \item[C.] Sigmund Freud.
        \item[D.] Friedrich Nietzsche.
    \end{itemize}
\end{enumerate}
\end{exercice}

\begin{exercice}[Vrai ou Faux justifié]
Indiquez si l'affirmation est VRAIE ou FAUSSE, puis justifiez votre réponse en une phrase précise en mobilisant le vocabulaire du cours.
\begin{enumerate}
    \item La conscience réflexive est la capacité de se connaître soi-même comme sujet pensant.
    \item L'inconscient freudien (\emph{Ça}) est un réservoir de souvenirs accessibles volontairement par la mémoire.
    \item Pour Sartre, l'inconscient n'existe pas car la conscience est transparence à elle-même ; ce qu'on appelle « inconscient » est de la \emph{mauvaise foi}.
    \item Le refoulement est un mécanisme de défense conscient par lequel on décide d'oublier un traumatisme.
\end{enumerate}
\end{exercice}

\begin{exercice}[Définitions et distinctions]
Définissez précisément les notions suivantes en deux lignes maximum chacune, en donnant un exemple concret pour chacune :
\begin{enumerate}
    \item \cle{Conscience immédiate} (ou conscience spontanée).
    \item \cle{Conscience réflexive}.
    \item \cle{Inconscient topique} (premi\`ere topique).
    \item \cle{Transfert} (dans la cure psychanalytique).
\end{enumerate}
\end{exercice}

\begin{exercice}[Identification des instances et mécanismes]
Complétez le tableau ci-dessous en associant chaque situation vécue à l'instance freudienne (Ça, Moi, Surmoi) ou au mécanisme de défense principalement concerné.

\begin{center}
\begin{tabularx}{\linewidth}{|l|X|X|}\hline
\textbf{Situation} & \textbf{Instance / Mécanisme} & \textbf{Justification brève} \\ \hline
Un élève a très faim en cours (Ça) mais attend la récréation pour manger. &  &  \\ \hline
Une jeune fille oublie systématiquement l'anniversaire de son père autoritaire. &  &  \\ \hline
Un candidat au Probatoire ressent une culpabilité intense après avoir triché. &  &  \\ \hline
Un homme battu par son chef se met à crier sur sa femme en rentrant chez lui. &  &  \\ \hline
\end{tabularx}
\end{center}
\end{exercice}

\courssub{Je m'entraîne}

\begin{exercice}[Analyse d'un cas : le lapsus révélateur]
Lors d'une réunion parents-professeurs au Lycée Technique de Douala, le proviseur veut dire : « Je souhaite la \emph{réussite} de tous les élèves. » Il dit par erreur : « Je souhaite la \emph{retraite} de tous les élèves. »
\begin{enumerate}
    \item Expliquez cet acte manqué selon la théorie freudienne (refoulement, conflit Ça/Moi/Surmoi, satisfaction substitutive).
    \item Comment Sartre (dans \emph{L'Être et le Néant}) interpréterait-il ce même « lapsus » ? Mobilisez la notion de \cle{mauvaise foi}.
    \item En quoi ces deux interprétations révèlent-elles deux conceptions opposées de la \cle{liberté} humaine ?
\end{enumerate}
\end{exercice}

\begin{exercice}[La conscience morale face au dilemme]
Un infirmier dans un hôpital de district au Cameroun découvre que le stock de médicaments antipaludéens (donnés par une ONG) est détourné par le directeur pour être revu en pharmacie privée. L'infirmier hésite à dénoncer car il a besoin de son salaire pour soigner sa propre mère.
\begin{enumerate}
    \item Identifiez le conflit de \cle{valeurs} chez cet infirmier (devoir moral vs intérêt personnel/solidarité familiale).
    \item Analysez la position de la \cle{conscience morale} (au sens kantien : impératif catégorique) face à ce dilemme.
    \item Analysez la position de la \cle{conscience psychologique} (au sens freudien : conflit Surmoi/Ça, angoisse morale) face à ce dilemme.
    \item Proposez une issue éthique argumentée en 5 lignes.
\end{enumerate}
\end{exercice}

\begin{exercice}[Critiques de la psychanalyse : confrontation d'arguments]
Le tableau ci-dessous présente trois critiques majeures de la psychanalyse freudienne. Complétez la colonne « Argument principal » et la colonne « Auteur/École » associé.

\begin{center}
\begin{tabularx}{\linewidth}{|l|X|X|}\hline
\textbf{Type de critique} & \textbf{Argument principal (à compléter)} & \textbf{Auteur / École (à compléter)} \\ \hline
Critique \emph{épistémologique} (statut scientifique) &  & Karl Popper / \emph{Falsifiabilité} \\ \hline
Critique \emph{philosophique} (liberté, subjectivité) & L'inconscient est une fiction ; l'homme est une passion inutile, il choisit son être en toute transparence. &  \\ \hline
Critique \emph{sociologique / anthropologique} & L'inconscient n'est pas individuel/atomisé mais \emph{collectif} et culturel ; l'Œdipe n'est pas universel. &  \\ \hline
\end{tabularx}
\end{center}
\end{exercice}

\begin{exercice}[Schéma et application : la seconde topique]
\begin{enumerate}
    \item Dessinez (schéma simple à la main ou description structurée) le modèle de la \emph{seconde topique} freudienne : \cle{Ça}, \cle{Moi}, \cle{Surmoi}. Indiquez les relations de force (flèches) et la part consciente/inconsciente de chaque instance (hachurage ou mention).
    \item En vous appuyant sur ce schéma, analysez la situation suivante : \emph{« Paul, élève en 1re T, veut réviser son Bac (Moi), mais il ressent une pulsion irrésistible de scroller TikTok pendant 3h (Ça). Ensuite, il se sent coupable et "nul" (Surmoi). »}
    \item Quel est le rôle du \cle{Moi} idéalement ? Pourquoi échoue-t-il souvent ici ?
\end{enumerate}
\end{exercice}

\courssub{Je me prépare au Bac}

\begin{exercice}[Sujet de dissertation type Baccalauréat Technique]
\textbf{Sujet :} \og La conscience suffit-elle à expliquer la vie psychique humaine ? \fg
\begin{enumerate}
    \item Analysez le sujet : définissez les termes, dégagez la problématique, annoncez le plan (2 grandes parties, 2 sous-parties chacune).
    \item Rédigez l'introduction complète (accroche, définition, problématique, annonce du plan).
    \item Rédigez la \textbf{Première partie} : \og La conscience, lumière nécessaire mais insuffisante de la vie psychique \fg (deux sous-parties argumentées avec auteurs : Descartes/Leibniz/Kant puis limites : Leibniz/perceptions obscures, Freud/inconscient).
    \item Rédigez la \textbf{Deuxième partie} : \og L'inconscient, dimension structurante et irréductible de la vie psychique \fg (deux sous-parties : Freud/topiques et mécanismes ; Critiques/Alternatives : Sartre/mauvaise foi, Jung/inconscient collectif, Ricœur/herméneutique).
    \item Rédigez une conclusion synthétique et une ouverture sur la responsabilité juridique au Cameroun (ex: irresponsabilité pénale pour trouble psychique).
\end{enumerate}
\end{exercice}

\begin{exercice}[Explication de texte : Freud, \emph{Métapsychologie} (1915)]
\textbf{Texte :} \og Il faut bien se rendre compte que le Moi n'est pas maître dans sa propre maison. Il dépend de trois puissances dures et il est menacé de trois côtés : du monde extérieur, de la libido du Ça, de la sévérité du Surmoi. [...] Le Moi sert trois maîtres sévères et fait de son mieux pour concilier leurs exigences. \fg
\begin{enumerate}
    \item \textbf{Analyse} : Identifiez la thèse principale du texte. Dégagez les trois « maîtres » du Moi et expliquez la nature de leurs exigences respectives.
    \item \textbf{Expliquez} les concepts suivants \emph{dans le contexte du texte} : \cle{Ça}, \cle{Surmoi}, \cle{libido}, \cle{conflit psychique}.
    \item \textbf{Discussion} : Ce texte signifie-t-il que l'homme est \emph{totalement} déterminé et sans liberté ? Confrontation avec Sartre (liberté absolue) et Ricœur (liberté comme capacité/attestation).
    \item \textbf{Exemple concret} : Illustrez le conflit « Moi / Ça / Surmoi » par une situation vécue par un élève en classe de Première Technique (ex: tricherie, absentéisme, choix d'orientation).
\end{enumerate}
\end{exercice}

\begin{exercice}[Question de réflexion guidée : Responsabilité et Inconscient]
\textbf{Sujet :} \og Peut-on être responsable de ses actes inconscients ? \fg
\textbf{Consignes :} Rédigez une réponse argumentée d'environ 20 lignes (200-250 mots). Votre réponse doit comporter :
\begin{enumerate}
    \item Une thèse claire (Oui / Non / Nuancé) annoncée dès la première phrase.
    \item Deux arguments philosophiques distincts (ex: distinction fait/droit, notion d'imputabilité, rôle de la psychanalyse, éthique de la lucidité).
    \item \textbf{Deux exemples précis tirés de la vie courante au Cameroun} (ex: accident de la route par somnolence/révélation d'un désir de mort ; vol par kleptomanie ; réaction violente "incontrôlée" liée à un traumatisme passé).
    \item Une conclusion qui rouvre sur la question de la \cle{guérison} ou de l'\cle{assomption} de son histoire.
\end{enumerate}
\end{exercice}

\newpage

\coursec{Corrigés des exercices}

\begin{corrige}[Exercice 1 — QCM : notions fondamentales]
\begin{enumerate}
    \item \textbf{B.} La conscience morale est la faculté de juger la valeur éthique de nos actes (bien/mal) ; la conscience psychologique est la simple présence à soi et au monde. A décrit la conscience spontanée, C est faux (tout homme a une conscience morale), D confond avec l'inconscient pulsionnel.
    \item \textbf{C.} Le \cle{Surmoi} (\emph{Über-Ich}) est l'instance des interdits parentaux et sociaux intériorisés, de l'idéal du Moi et de la culpabilité. Le Ça est le pôle pulsionnel, le Moi l'instance de médiation.
    \item \textbf{C.} Pour Freud (\emph{Psychopathologie de la vie quotidienne}), le lapsus est un \emph{acte manqué} : un compromis entre l'intention consciente et un désir refoulé qui trouve un déguisement pour s'exprimer.
    \item \textbf{C.} C'est \cle{Freud} qui déclare que « le Moi n'est pas maître dans sa propre maison », formule de la troisième blessure narcissique infligée à l'homme (après Copernic et Darwin).
\end{enumerate}
\textbf{Points de vigilance :} ne pas confondre \emph{conscience morale} (jugement de valeur) et \emph{Surmoi} (instance psychique inconsciente en grande partie) ; ne pas attribuer la formule à Nietzsche, qui a préparé la critique de la conscience sans la formuler ainsi.
\end{corrige}

\begin{corrige}[Exercice 2 — Vrai ou Faux justifié]
\begin{enumerate}
    \item \textbf{VRAI.} La \cle{conscience réflexive} est le retour de l'esprit sur lui-même : le sujet se prend lui-même pour objet de pensée (\emph{je pense que je pense}), ce qui fonde le \emph{Cogito} cartésien.
    \item \textbf{FAUX.} Le Ça freudien est un réservoir de pulsions et de contenus \emph{refoulés}, inaccessibles à la mémoire volontaire ; ce qui est accessible volontairement relève du conscient, et ce qui est rappelable relève du \emph{préconscient}.
    \item \textbf{VRAI.} Dans \emph{L'Être et le Néant}, Sartre refuse l'inconscient : la conscience est translucide, totalement présente à elle-même ; prétendre « ne pas savoir » ce qu'on sait est de la \cle{mauvaise foi}, mensonge à soi-même.
    \item \textbf{FAUX.} Le \cle{refoulement} est précisément un mécanisme de défense \emph{inconscient} et involontaire : le Moi ne « décide » pas d'oublier ; le contenu pénible est chassé hors de la conscience sans que le sujet en ait conscience. L'oubli délibéré serait de la suppression consciente.
\end{enumerate}
\textbf{Points de vigilance :} la justification doit mobiliser le vocabulaire du cours (refoulement, préconscient, mauvaise foi) et non une simple reformulation de l'affirmation.
\end{corrige}

\begin{corrige}[Exercice 3 — Définitions et distinctions]
\begin{enumerate}
    \item \textbf{Conscience immédiate (spontanée) :} présence simple et directe du sujet au monde et à ses actes, sans retour sur soi. \emph{Exemple :} le chauffeur de taxi-brousse à Douala qui conduit en parlant, absorbé par la route.
    \item \textbf{Conscience réflexive :} capacité de l'esprit à se prendre lui-même pour objet, à se connaître comme sujet pensant. \emph{Exemple :} l'élève qui s'interroge : « pourquoi ai-je menti au proviseur ? »
    \item \textbf{Inconscient topique (première topique) :} dans le modèle Inconscient/Préconscient/Conscient, le système des contenus refoulés (pulsions, souvenirs pénibles) soumis à la censure et s'exprimant de façon déguisée. \emph{Exemple :} un rêve de chute qui exprime une angoisse refoulée.
    \item \textbf{Transfert :} dans la cure psychanalytique, déplacement sur la personne du psychanalyste des sentiments (amour, hostilité) éprouvés autrefois envers des figures parentales. \emph{Exemple :} le patient qui se met en colère contre son analyste comme il le faisait contre son père.
\end{enumerate}
\textbf{Points de vigilance :} chaque définition doit comporter un trait distinctif (retour sur soi, refoulement, déplacement) et un exemple \emph{concret}, pas une définition par synonyme.
\end{corrige}

\begin{corrige}[Exercice 4 — Identification des instances et mécanismes]
\begin{center}
\begin{tabularx}{\linewidth}{|l|X|X|}\hline
\rowcolor{popblueL}\textbf{Situation} & \textbf{Instance / Mécanisme} & \textbf{Justification brève} \\ \hline
Un élève a très faim en cours mais attend la récréation pour manger. & \textbf{Moi} (principe de réalité) & Le Moi diffère la satisfaction pulsionnelle du Ça pour tenir compte des contraintes du réel (règlement, horaires). \\ \hline
Une jeune fille oublie systématiquement l'anniversaire de son père autoritaire. & \textbf{Refoulement} (acte manqué) & L'oubli répété n'est pas fortuit : il exprime de façon déguisée une hostilité refoulée envers le père. \\ \hline
Un candidat au Probatoire ressent une culpabilité intense après avoir triché. & \textbf{Surmoi} & Le Surmoi, interdit moral intériorisé, juge et punit le Moi par le sentiment de culpabilité. \\ \hline
Un homme battu par son chef se met à crier sur sa femme en rentrant. & \textbf{Déplacement} & L'agressivité ne pouvant s'exercer contre le chef (objet dangereux), elle est reportée sur un objet substitutif plus sûr. \\ \hline
\end{tabularx}
\end{center}
\textbf{Points de vigilance :} distinguer \emph{instance} (Ça, Moi, Surmoi — seconde topique) et \emph{mécanisme de défense} (refoulement, déplacement, sublimation, projection) ; la justification doit nommer le processus, pas seulement décrire la scène.
\end{corrige}

\begin{corrige}[Exercice 5 — Analyse d'un cas : le lapsus révélateur]
\begin{enumerate}
    \item \textbf{Interprétation freudienne.} Pour Freud, le lapsus n'est jamais un hasard : c'est un \emph{acte manqué}, c'est-à-dire un compromis entre une intention consciente (souhaiter la réussite) et un désir refoulé (le proviseur, fatigué, désire secrètement sa propre retraite, ou souhaite inconsciemment se débarrasser des élèves). Le désir refoulé, maintenu hors de la conscience par la censure, profite d'une baisse d'attention pour s'exprimer de façon déguisée : le mot substitué procure une \emph{satisfaction substitutive}. Le Moi, croyant parler librement, est trahi par le Ça ; le Surmoi peut ensuite produire gêne et honte.
    \item \textbf{Interprétation sartrienne.} Sartre refuse l'hypothèse d'un désir caché dans un « inconscient » : pour lui, la conscience est entièrement transparente à elle-même. Le proviseur \emph{sait}, au fond, qu'il aspire à la retraite ; feindre de l'ignorer est de la \cle{mauvaise foi} : il se ment à lui-même en se donnant l'excuse de l'« erreur » pour ne pas assumer son désir. Le lapsus est donc un choix déguisé, non une détermination mécanique.
    \item \textbf{Deux conceptions de la liberté.} Chez Freud, la liberté du sujet est \emph{limitée} : le Moi n'est pas maître dans sa propre maison, il est déterminé par des forces pulsionnelles qui le dépassent ; la liberté ne s'acquiert que par la prise de conscience (la cure). Chez Sartre, la liberté est \emph{totale et inconditionnée} : « l'homme est condamné à être libre », aucune instance inconsciente ne peut l'excuser ; invoquer l'inconscient est une fuite de la responsabilité. Le même fait (le lapsus) est donc lu comme \emph{déterminisme psychique} chez l'un et comme \emph{choix de mauvaise foi} chez l'autre.
\end{enumerate}
\textbf{Points de vigilance :} bien articuler les trois termes freudiens demandés (refoulement, conflit des instances, satisfaction substitutive) ; pour Sartre, la mauvaise foi doit être définie (mensonge à soi-même) et non seulement nommée.
\end{corrige}

\begin{corrige}[Exercice 6 — La conscience morale face au dilemme]
\begin{enumerate}
    \item \textbf{Conflit de valeurs.} L'infirmier est déchiré entre le \emph{devoir moral} (dénoncer le détournement, protéger les malades, respecter la justice et le serment professionnel) et l'\emph{intérêt personnel et la solidarité familiale} (conserver son salaire pour soigner sa mère). C'est un conflit entre l'universel (le droit des patients) et le particulier (l'attachement filial).
    \item \textbf{Conscience morale kantienne.} Pour Kant, la conscience morale obéit à l'\emph{impératif catégorique} : « Agis de telle sorte que la maxime de ton action puisse être érigée en loi universelle. » Or on ne peut vouloir universellement que les médicaments des malades soient détournés, ni que chacun se taise par intérêt. Le devoir commande donc de dénoncer, quelles que soient les conséquences : la moralité d'un acte réside dans l'intention conforme au devoir, non dans le calcul des avantages.
    \item \textbf{Conscience psychologique freudienne.} Freud décrirait un conflit d'instances : le \emph{Surmoi} (interdits et idéaux intériorisés : honnêteté, serment) exige la dénonciation ; le \emph{Ça} et les intérêts du \emph{Moi} (peur de perdre le salaire, amour filial) poussent au silence. De ce conflit naît l'\emph{angoisse morale} et la culpabilité : quelle que soit la décision, une exigence sera frustrée, d'où le malaise psychique de l'infirmier.
    \item \textbf{Issue éthique proposée.} L'infirmier peut dénoncer le détournement par la voie hiérarchique ou auprès de l'ONG, éventuellement de manière anonyme ou collective (avec d'autres infirmiers) pour limiter le risque de licenciement. Cette solution concilie le devoir de justice envers les malades — dont la vie est en jeu, valeur supérieure — et la prudence légitime à l'égard de sa famille. Taire un détournement mortel serait une complicité ; la conscience morale commande d'agir, mais la sagesse commande de choisir les moyens les moins destructeurs.
\end{enumerate}
\textbf{Points de vigilance :} ne pas confondre conscience morale (norme universelle) et conscience psychologique (vécu du conflit) ; l'issue proposée doit être \emph{argumentée} (hiérarchie des valeurs) et non seulement pratique.
\end{corrige}

\begin{corrige}[Exercice 7 — Critiques de la psychanalyse : confrontation d'arguments]
\begin{center}
\begin{tabularx}{\linewidth}{|l|X|X|}\hline
\rowcolor{popblueL}\textbf{Type de critique} & \textbf{Argument principal} & \textbf{Auteur / École} \\ \hline
Critique \emph{épistémologique} (statut scientifique) & Une théorie scientifique doit être \emph{falsifiable} (réfutable par l'expérience). Or la psychanalyse explique tout et son contraire : toute objection est interprétée comme « résistance », donc aucun fait ne peut l'invalider ; elle n'est pas une science. & Karl Popper / \emph{Falsifiabilité} \\ \hline
Critique \emph{philosophique} (liberté, subjectivité) & L'inconscient est une fiction ; l'homme est une passion inutile, il choisit son être en toute transparence. & \textbf{Jean-Paul Sartre} / existentialisme (\emph{L'Être et le Néant}) \\ \hline
Critique \emph{sociologique / anthropologique} & L'inconscient n'est pas individuel/atomisé mais \emph{collectif} et culturel ; l'Œdipe n'est pas universel. & \textbf{Carl Gustav Jung} (inconscient collectif, archétypes) ; appui ethnologique de \textbf{Malinowski} (matriarcat des Trobriand) \\ \hline
\end{tabularx}
\end{center}
\textbf{Points de vigilance :} pour Popper, le mot-clé exigible est \emph{falsifiabilité/réfutabilité} ; pour la critique anthropologique, Jung (inconscient collectif) est l'auteur du cours, Malinowski fournit l'argument ethnographique contre l'universalité du complexe d'Œdipe.
\end{corrige}

\begin{corrige}[Exercice 8 — Schéma et application : la seconde topique]
\begin{enumerate}
    \item \textbf{Schéma attendu.} Trois zones : le \textbf{Ça} (entièrement \emph{inconscient}, réservoir des pulsions, régi par le principe de plaisir) ; le \textbf{Surmoi} (en grande partie \emph{inconscient}, interdits et idéal, source de culpabilité) ; le \textbf{Moi} (en partie \emph{conscient}, en partie inconscient, médiateur régi par le principe de réalité). Flèches : le Ça pousse ses pulsions vers le Moi ; le Surmoi presse le Moi d'interdits ; le monde extérieur contraint le Moi. Le Moi est au centre, « serviteur de trois maîtres ».
\begin{popfigure}
\begin{center}
\begin{tikzpicture}[scale=1]
    \draw[fill=poppurpleL,draw=poppurple] (0,0) circle (1.6cm);
    \node[poppurple] at (0,0.9) {\textbf{Ça} (inconscient)};
    \node[poppurple,font=\small] at (0,0.4) {pulsions};
    \draw[fill=popblueL,draw=popblue] (2.6,0.6) circle (1.2cm);
    \node[popblue] at (2.6,1.1) {\textbf{Moi}};
    \node[popblue,font=\small] at (2.6,0.6) {partie consciente};
    \draw[fill=popgoldL,draw=popgold] (2.6,2.6) circle (1.1cm);
    \node[popgold] at (2.6,3.0) {\textbf{Surmoi}};
    \node[popgold,font=\small] at (2.6,2.5) {interdits};
    \draw[-{Stealth},thick,poppurple] (1.5,0.3) -- (1.6,0.5);
    \draw[-{Stealth},thick,poppurple] (1.6,0.2) -- (1.5,0.55);
    \draw[-{Stealth},thick,popgold] (2.6,1.5) -- (2.6,1.75);
    \draw[-{Stealth},thick,popdark] (4.6,0.6) -- (3.85,0.6);
    \node[popdark,font=\small] at (5.4,0.6) {Réalité};
\end{tikzpicture}
\end{center}
\legende{Seconde topique freudienne : le Moi, médiateur entre les pulsions du Ça, les interdits du Surmoi et les exigences de la réalité extérieure.}
\end{popfigure}
    \item \textbf{Analyse de la situation de Paul.} Le \emph{Ça} exige le plaisir immédiat (scroller TikTok, principe de plaisir) ; le \emph{Moi} conscient pose le projet rationnel de réviser le Bac (principe de réalité : différer la satisfaction pour un gain futur) ; le \emph{Surmoi}, après la faute, condamne le Moi par la culpabilité et le sentiment de dévalorisation (« je suis nul »). Paul illustre le Moi « malmené » : il a cédé à la pulsion, puis subi le verdict du Surmoi.
    \item \textbf{Rôle idéal du Moi.} Le Moi doit \emph{concilier} les trois exigences : accorder au Ça des satisfactions réalistes et socialement admises (pause détente limitée), sans encourir les foudres du Surmoi. Il échoue souvent ici parce que la pulsion numérique est immédiate et puissante (renforcement constant), tandis que la récompense de la révision est lointaine ; le Moi, faible et dépendant, n'a pas la force de différer. La \emph{sublimation} (transformer l'énergie pulsionnelle en travail) serait l'issue freudienne idéale.
\end{enumerate}
\textbf{Points de vigilance :} le schéma doit indiquer la part consciente/inconsciente de chaque instance ; l'analyse doit nommer les principes (plaisir/réalité) et non seulement raconter la scène.
\end{corrige}

\begin{corrige}[Exercice 9 — Sujet de dissertation type Baccalauréat Technique]
\textbf{1. Analyse du sujet.} \emph{Conscience} : présence du sujet à lui-même et au monde (psychologique) et faculté de juger ses actes (morale). \emph{Suffire} : être la condition unique et totale. \emph{Vie psychique} : ensemble des phénomènes de la vie mentale (pensées, désirs, rêves, émotions). \textbf{Problématique :} la conscience, longtemps tenue pour l'essence même de l'esprit, épuise-t-elle la réalité psychique, ou celle-ci comporte-t-elle une dimension inconsciente irréductible ? \textbf{Plan :} I. La conscience, lumière nécessaire mais insuffisante ; II. L'inconscient, dimension structurante mais discutée.

\textbf{2. Introduction (rédigée).} Depuis toujours, l'homme se définit comme un être qui pense et qui se sait pensant : « Je pense, donc je suis », affirme Descartes. La conscience semble être le propre de l'homme, la lumière qui éclaire tous ses actes et fonde sa responsabilité. Pourtant, qui n'a jamais été surpris par un rêve étrange, un oubli inexplicable, un lapsus révélateur ? Ces phénomènes suggèrent qu'une part de notre vie mentale nous échappe. Dès lors, la conscience suffit-elle à expliquer la vie psychique humaine ? Autrement dit, l'esprit se réduit-il à ce qu'il sait de lui-même, ou comporte-t-il des zones obscures qui nous déterminent à notre insu ? Nous verrons d'abord que la conscience est une lumière nécessaire mais insuffisante, puis que l'inconscient constitue une dimension structurante de la vie psychique, encore que discutée.

\textbf{3. Première partie (rédigée).} \emph{A. La conscience, fondement de la subjectivité.} Pour Descartes, le \emph{Cogito} est la première vérité indubitable : je ne puis douter que je pense. La conscience est connaissance immédiate de soi, siège de la pensée claire et distincte, et fondement de la liberté et de la responsabilité morale. Kant en fait la condition de la moralité : la conscience morale, voix du devoir, juge nos actes à la lumière de l'impératif catégorique. Sans conscience, point de sujet, point de droit, point de morale. \emph{B. Mais ses limites.} Déjà Leibniz objecte à Descartes l'existence de \emph{petites perceptions} obscures et insensiblement actives : le bruit de la mer est fait de mille bruits de gouttes que nous n'apercevons pas distinctement. La conscience n'éclaire donc qu'une partie du psychisme. Freud radicalise cette limite : rêves, lapsus, actes manqués, symptômes névrotiques témoignent d'un \emph{inconscient} qui pense et désire à notre insu. La conscience serait la pointe émergée de l'iceberg psychique.

\textbf{4. Deuxième partie (rédigée).} \emph{A. L'inconscient freudien, instance structurante.} Freud propose deux topiques : la première (Inconscient/Préconscient/Conscient) décrit un système de contenus refoulés s'exprimant de façon déguisée (rêves comme « accomplissements de désirs ») ; la seconde (Ça/Moi/Surmoi) montre le Moi « non maître dans sa propre maison », tiraillé entre pulsions, interdits et réalité. Le refoulement, le déplacement, la sublimation organisent notre vie affective : l'inconscient n'est pas un résidu mais le moteur caché du psychisme, et la cure analytique prouve son efficacité en faisant parler les symptômes. \emph{B. Critiques et alternatives.} Sartre dénie à l'inconscient toute réalité : la conscience est transparence, et l'« inconscient » n'est que mauvaise foi, refus d'assumer sa liberté. Jung, sans nier l'inconscient, en conteste le monopole individuel : il existe un \emph{inconscient collectif} fait d'archétypes culturels. Ricœur propose une voie médiane : l'inconscient est un \emph{langage} à interpréter (herméneutique), et la liberté n'est pas pure transparence mais \emph{attestation} d'une capacité à travers ses propres opacités.

\textbf{5. Conclusion.} La conscience ne suffit donc pas à expliquer la vie psychique : elle en est la lumière nécessaire, mais non la totalité. L'inconscient, découvert par Freud, en révèle la profondeur déterminante, sans pour autant abolir la responsabilité : connaître ses déterminations, c'est déjà commencer à s'en libérer. Cette question déborde la psychologie : le droit camerounais reconnaît l'\emph{irresponsabilité pénale} en cas de trouble psychique ayant aboli le discernement — preuve que la question de la conscience demeure au cœur de la justice et de la vie sociale.

\textbf{Barème indicatif (sur 20) :} analyse du sujet et problématique : 4 pts ; introduction : 3 pts ; première partie (2 sous-parties argumentées, auteurs cités) : 5 pts ; deuxième partie (idem) : 5 pts ; conclusion et ouverture : 2 pts ; langue et cohérence : 1 pt.
\textbf{Points de vigilance :} chaque sous-partie doit articuler \emph{thèse + auteur + argument + exemple} ; ne pas opposer conscience et inconscient comme deux « choses » mais comme deux modèles du psychisme ; l'ouverture doit prolonger, non répéter.
\end{corrige}

\begin{corrige}[Exercice 10 — Explication de texte : Freud, \emph{Le Moi et le Ça} (1923) / \emph{Nouvelles Conférences d'introduction à la psychanalyse} (1933)]
\begin{enumerate}
    \item \textbf{Analyse.} Thèse principale : le Moi n'est pas le maître autonome que la tradition (Descartes) croyait ; il est un médiateur faible, dépendant de trois puissances. Les trois « maîtres » : (a) le \emph{monde extérieur}, qui impose ses contraintes matérielles et sociales (principe de réalité) ; (b) la \emph{libido du Ça}, c'est-à-dire l'énergie des pulsions qui exigent satisfaction immédiate (principe de plaisir) ; (c) la \emph{sévérité du Surmoi}, qui impose interdits et idéaux et punit par la culpabilité. Le Moi doit « concilier » ces exigences contradictoires : sa fonction est de compromis, non de commandement.
    \item \textbf{Concepts.} \emph{Ça} : instance pulsionnelle entièrement inconsciente, réservoir des désirs. \emph{Surmoi} : instance issue de l'intériorisation des interdits parentaux et sociaux, juge du Moi. \emph{Libido} : énergie psychique des pulsions (d'abord sexuelles au sens large). \emph{Conflit psychique} : état de tension permanent entre ces exigences incompatibles, source de l'angoisse et des symptômes.
    \item \textbf{Discussion.} Le texte instaure un \emph{déterminisme psychique} relatif : le Moi est contraint, mais non aboli — il « fait de son mieux pour concilier », ce qui laisse une marge de négociation, accrue par la cure (rendre conscient l'inconscient). Sartre objecterait que ce déterminisme est une excuse de mauvaise foi : l'homme est liberté absolue, « condamné à être libre ». Ricœur propose une position intermédiaire : la liberté n'est ni maîtrise transparente ni esclavage, mais \emph{capacité attestée} de se reprendre à travers ses déterminations. Freud ne supprime donc pas la liberté, il la déplace : d'un pouvoir donné à une conquête.
    \item \textbf{Exemple concret.} Un élève de Première Technique, la veille d'une évaluation, est tiraillé : le \emph{Ça} l'incite à tricher (plaisir immédiat de la réussite sans effort), le \emph{Surmoi} lui rappelle l'interdit de la fraude et l'honneur familial, la \emph{réalité} (surveillant, sanctions) limite les possibilités. Le Moi négocie : réviser sérieusement, ou céder et subir ensuite la culpabilité. La scène illustre exactement les « trois maîtres » du texte.
\end{enumerate}
\textbf{Barème indicatif :} analyse (thèse + trois maîtres) : 5 pts ; concepts : 4 pts ; discussion argumentée avec confrontation : 7 pts ; exemple : 3 pts ; expression : 1 pt.
\textbf{Points de vigilance :} définir chaque concept \emph{dans le contexte du texte} ; la discussion doit confronter trois positions nommées et conclure par une thèse personnelle justifiée.
\end{corrige}

\begin{corrige}[Exercice 11 — Question de réflexion guidée : Responsabilité et Inconscient]
\textbf{Corrigé type (environ 230 mots).}

\emph{Thèse nuancée :} on ne peut être tenu pleinement responsable d'actes dont le mobile nous échappe, mais on demeure responsable de la manière dont on répond à ce que l'inconscient révèle de nous.

Premier argument : il faut distinguer le \emph{fait} et le \emph{droit}. En fait, l'inconscient nous détermine : Freud montre que lapsus, rêves et symptômes expriment des désirs refoulés qui agissent à notre insu. Mais en droit, la responsabilité suppose l'\emph{imputabilité} : n'est punissable que celui qui avait conscience et discernement de son acte. C'est pourquoi le droit pénal camerounais reconnaît l'irresponsabilité en cas de trouble psychique ayant aboli le discernement.

Second argument : la psychanalyse elle-même fonde une responsabilité d'un nouveau type. Si le sujet n'est pas coupable de ses pulsions, il est \emph{responsable de les connaître} : l'éthique de la lucidité invite à assumer son histoire au lieu de la subir. Comme le suggère Ricœur, la liberté n'est pas la transparence sartrienne mais la capacité de se reprendre.

Exemples : un chauffeur de transport en commun qui, épuisé, s'endort au volant sur l'axe Yaoundé--Douala et provoque un accident n'a pas voulu le drame, mais sa responsabilité est engagée s'il a négligé le repos ; une kleptomane qui vole compulsivement au marché relève du soin psychiatrique plus que de la prison, car son geste échappe à sa volonté délibérée.

En conclusion, la responsabilité ne porte pas sur l'acte inconscient lui-même, mais sur l'\emph{assomption} de son histoire : guérir, c'est devenir capable de répondre de ce qu'on n'a pas choisi.

\textbf{Barème indicatif :} thèse claire dès la première phrase : 3 pts ; deux arguments distincts et justes : 6 pts ; deux exemples camerounais précis et pertinents : 6 pts ; conclusion ouvrant sur guérison/assomption : 3 pts ; expression et respect du volume : 2 pts.
\textbf{Points de vigilance :} éviter la réponse manichéenne sans nuance ; les exemples doivent être \emph{précis} (situation, lieu, enjeu) et non génériques ; la distinction fait/droit est attendue.
\end{corrige}

\newpage

\coursec{Fiche bilan}

\begin{popfigure}
\centering
\begin{tikzpicture}[
  mindmap,
  grow cyclic,
  every node/.style={concept, execute at begin node=\hskip0pt},
  concept color=popblue,
  level 1/.append style={level distance=4.5cm, sibling angle=60, font=\small\bfseries},
  level 2/.append style={level distance=3cm, sibling angle=45, font=\footnotesize},
  branch/.style={concept color=#1, edge from parent/.style={draw=#1, thick}}
]
\node[concept, text width=3.2cm, align=center, font=\normalsize\bfseries] {Conscience \\ et Inconscient}
  child[branch=poporange] { node[concept, text width=2.5cm, align=center] {1. La Conscience}
    child { node[concept, text width=2cm, align=center] {Définition : \textit{cogito}, présence à soi} }
    child { node[concept, text width=2cm, align=center] {Types : morale, psychologique, immédiate/réfléchie} }
    child { node[concept, text width=2cm, align=center] {Caractères : unité, transparence, intentionalité} }
    child { node[concept, text width=2cm, align=center] {Valeur : liberté, responsabilité, dignité} }
    child { node[concept, text width=2cm, align=center] {Limites : erreur, illusion, conditionnement} }
  }
  child[branch=popgreen] { node[concept, text width=2.5cm, align=center] {2. L'Inconscient}
    child { node[concept, text width=2cm, align=center] {Découverte : Leibniz, Janet, Freud} }
    child { node[concept, text width=2cm, align=center] {Types : freudien (refoulé), cognitif, corporel} }
    child { node[concept, text width=2cm, align=center] {Manifestations : lapsus, rêves, symptômes, actes manqués} }
  }
  child[branch=poppurple] { node[concept, text width=2.5cm, align=center] {3. Psychanalyse freudienne}
    child { node[concept, text width=2cm, align=center] {Topique : Ça, Moi, Surmoi} }
    child { node[concept, text width=2cm, align=center] {Mécanismes : refoulement, sublimation, déplacement} }
    child { node[concept, text width=2cm, align=center] {Critique : Popper (non-falsifiable), Sartre (mauvaise foi), Ricoeur (herméneutique)} }
  }
  child[branch=poppink] { node[concept, text width=2.5cm, align=center] {4. Articulation C/InC}
    child { node[concept, text width=2cm, align=center] {Conflit : le Moi n'est pas maître en sa maison} }
    child { node[concept, text width=2cm, align=center] {Complémentarité : adaptation, créativité, santé} }
    child { node[concept, text width=2cm, align=center] {Idéal : lucidité (Sartre) ou intégration (Jung/Ricoeur)} }
  }
  child[branch=popgold] { node[concept, text width=2.5cm, align=center] {5. Auteurs clés}
    child { node[concept, text width=2cm, align=center] {Descartes : certitude, transparence} }
    child { node[concept, text width=2cm, align=center] {Leibniz : petits perceptions, continuité} }
    child { node[concept, text width=2cm, align=center] {Freud : découverte scientifique, topique} }
    child { node[concept, text width=2cm, align=center] {Sartre : conscience intentionnelle, refus InC} }
    child { node[concept, text width=2cm, align=center] {Ricoeur : archéologie vs téléologie} }
  }
  child[branch=popdark] { node[concept, text width=2.5cm, align=center] {6. Méthodes Bac}
    child { node[concept, text width=2cm, align=center] {Dissertation : analyser sujet, thèse/antithèse, synthèse} }
    child { node[concept, text width=2cm, align=center] {Commentaire : expliquer, analyser, discuter} }
    child { node[concept, text width=2cm, align=center] {Exemple camerounais : rites, rêves, éducation} }
  };
\end{tikzpicture}
\legende{Carte mentale récapitulative : Conscience et Inconscient (Première Technique)}
\end{popfigure}

\begin{bilan}
\courssub{Définitions clés}
\begin{itemize}
  \item \cle{Conscience} : Connaissance immédiate et intuitive que le sujet a de ses propres états mentaux (pensées, sentiments, volontés). \textbf{Étymologie} : \textit{cum-scire} (savoir avec).
  \item \cle{Conscience morale} : Jugement de valeur sur ses actes (bien/mal), voix de la raison pratique (Kant : impératif catégorique).
  \item \cle{Conscience psychologique} : Fonction de veille, d'attention et de représentation du monde (Bergson : choix, action).
  \item \cle{Inconscient (Freud)} : Ensemble des processus psychiques refoulés, inaccessibles à la conscience mais agissants (pulsions, désirs).
  \item \cle{Inconscient cognitif} : Traitements automatiques, non conscients, de l'information (mémoire procédurale, perception).
  \item \cle{Refoulement} : Mécanisme de défense expulsant de la conscience des représentations incompatibles avec le Moi/Idéal du Moi.
\end{itemize}

\courssub{Repères théoriques (à connaître par cœur)}
\begin{center}
\begin{tabularx}{\linewidth}{|l|X|X|}\hline
\rowcolor{popblueL}
\textbf{Auteur / École} & \textbf{Thèse centrale sur Conscience / Inconscient} & \textbf{Concept clé} \\ \hline
\textbf{Descartes} & Conscience = pensée (\textit{cogito}), transparence totale, certitude absolue. Inconscient = contradiction. & \textit{Cogito, ergo sum} \\ \hline
\textbf{Leibniz} & Continuité : \cle{petites perceptions} (inconscientes) $\to$ perception consciente (seuil). Monadologie. & Seuil de conscience \\ \hline
\textbf{Freud} & Inconscient dynamique (refoulé) structuré : \textbf{Ça} (pulsions), \textbf{Moi} (réalité), \textbf{Surmoi} (idéal/moral). & Topique / Refoulement \\ \hline
\textbf{Sartre} & Refus de l'inconscient freudien. Conscience = \cle{intentionnalité} (toujours conscience \textit{de} quelque chose). \cle{Mauvaise foi} = mensonge à soi. & Pour-soi / En-soi \\ \hline
\textbf{Ricoeur} & \cle{École du soupçon} (Marx, Nietzsche, Freud). Herméneutique : déchiffrer les symboles. Archaeologie vs Téléologie. & Interprétation \\ \hline
\end{tabularx}
\end{center}

\courssub{Méthodes express}
\begin{methode}[Dissertation : « La conscience est-elle transparente à elle-même ? »]
\begin{enumerate}
  \item Analyser : \textit{transparente} = lucidité totale, absence d'ombre/erreur. \textit{Elle-même} = reflexivité.
  \item Thèse (Oui) : Descartes (cogito), conscience morale (Kant) – je sais ce que je veux/fais.
  \item Antithèse (Non) : Leibniz (petites perceptions), Freud (inconscient refoulé, lapsus), Marx (idéologie), conditionnements sociaux (Bourdieu).
  \item Synthèse : Conscience comme \textbf{conquête} (Ricoeur, Merleau-Ponty) : lucidité partielle, travail d'élucidation (psychanalyse, réflexion critique). Exemple : un chef d'entreprise camerounais prenant une décision « rationnelle » mais motivée inconsciemment par la peur de l'échec ou le désir de reconnaissance clanique.
\end{enumerate}
\end{methode}

\begin{methode}[Commentaire de texte (extrait Freud / Sartre / Descartes)]
\begin{enumerate}
  \item \textbf{Contextualiser} : Auteur, ouvrage, époque, débat visé.
  \item \textbf{Structurer} : Identifier la thèse, les arguments, les exemples, les concepts techniques (Ça, Moi, Surmoi ; Mauvaise foi ; Cogito).
  \item \textbf{Expliquer} : Définir les concepts \textit{in situ}. Montrer l'enchaînement logique.
  \item \textbf{Discuter} : Valider (cohérence, puissance explicative) / Critiquer (réductionnisme, impasses éthiques, contre-exemples concrets). Conclure sur la portée philosophique.
\end{enumerate}
\end{methode}
\end{bilan}

\begin{aretenir}[Checklist avant le Bac Probatoire]
  $\square$ Savoir définir \textbf{conscience} (psychologique vs morale) et \textbf{inconscient} (freudien vs cognitif).\\
  $\square$ Citer les \textbf{3 types de conscience} : immédiate, réfléchie, morale.\\
  $\square$ Expliquer les \textbf{4 caractéristiques} : unité, transparence (Discours de la méthode), intentionalité (Brentano/Sartre), temporalité.\\
  $\square$ Donner \textbf{2 valeurs} (liberté, responsabilité, dignité) et \textbf{3 limites} (erreur, illusion, déterminismes/inconscient).\\
  $\square$ Raconter la \textbf{découverte de l'inconscient} : Leibniz (seuil) $\to$ Janet (dissociation) $\to$ Freud (scientifique/topique).\\
  $\square$ Dessiner et nommer la \textbf{Topique freudienne} : Ça / Moi / Surmoi + principes (plaisir / réalité / moralité).\\
  $\square$ Citer \textbf{3 mécanismes de défense} : refoulement, sublimation, déplacement, projection, rationalisation.\\
  $\square$ Contraster \textbf{Freud vs Sartre} : Inconscient structurel vs Mauvaise foi (liberté totale).\\
  $\square$ Résumer la \textbf{critique de Popper} (non-falsifiabilité) et l'\textbf{herméneutique de Ricoeur} (interprétation vs causalité).\\
  $\square$ Avoir \textbf{2 exemples concrets camerounais} : rêve initiatique / interprétation traditionnelle vs psychanalyse ; lapsus en langue locale / française ; éducation familiale (Surmoi clanique).\\
  $\square$ Maîtriser le \textbf{plan type dissertation} : Intro (accroche, définitions, problème, plan) $\to$ 2 ou 3 parties argumentées (auteurs + exemples) $\to$ Conclusion (réponse + ouverture).
\end{aretenir}

\findecours

\end{document}
